\documentclass[11pt,oneside]{book}

\usepackage[T1]{fontenc}
\usepackage[utf8]{inputenc}
\usepackage{lmodern}
\usepackage{microtype}
\usepackage[letterpaper,margin=1.02in,headheight=15pt]{geometry}
\usepackage{amsmath,amssymb,amsthm,mathtools}
\usepackage{bm}
\usepackage{booktabs,longtable,array,tabularx}
\usepackage{enumitem}
\usepackage{xcolor}
\usepackage{graphicx}
\usepackage{tikz}
\usetikzlibrary{arrows.meta,calc,decorations.pathmorphing,fit,matrix,positioning,shapes.geometric}
\usepackage{pgfplots}
\pgfplotsset{compat=1.18}
\usepackage[most]{tcolorbox}
\usepackage{fancyhdr}
\usepackage{hyperref}
\usepackage[nameinlink,capitalise,noabbrev]{cleveref}

\definecolor{deepblue}{RGB}{32,63,96}
\definecolor{softblue}{RGB}{231,240,249}
\definecolor{deepgreen}{RGB}{42,92,76}
\definecolor{softgreen}{RGB}{234,245,239}
\definecolor{deeporange}{RGB}{150,78,20}
\definecolor{softorange}{RGB}{252,242,229}
\definecolor{deepred}{RGB}{135,45,52}
\definecolor{softred}{RGB}{250,234,235}
\definecolor{softgray}{RGB}{245,245,245}

\hypersetup{
  colorlinks=true,
  linkcolor=deepblue,
  citecolor=deepgreen,
  urlcolor=deeporange,
  pdftitle={Transseries for Mere Mortals},
  pdfauthor={OpenAI},
  pdfsubject={A self-contained tutorial on formal and resurgent transseries},
  pdfkeywords={transseries, asymptotics, Hahn series, resurgence, Borel summation, Stokes phenomenon}
}

\pagestyle{fancy}
\fancyhf{}
\fancyhead[L]{\nouppercase{\leftmark}}
\fancyhead[R]{\thepage}
\renewcommand{\headrulewidth}{0.4pt}

\setcounter{secnumdepth}{3}
\setcounter{tocdepth}{2}
\setlist[itemize]{topsep=3pt,itemsep=2pt,parsep=1pt}
\setlist[enumerate]{topsep=3pt,itemsep=2pt,parsep=1pt}
\allowdisplaybreaks

\newtheorem{theorem}{Theorem}[chapter]
\newtheorem{proposition}[theorem]{Proposition}
\newtheorem{lemma}[theorem]{Lemma}
\newtheorem{corollary}[theorem]{Corollary}
\theoremstyle{definition}
\newtheorem{definition}[theorem]{Definition}
\newtheorem{example}[theorem]{Example}
\newtheorem{exercise}{Exercise}[chapter]
\theoremstyle{remark}
\newtheorem{remark}[theorem]{Remark}

\newtcolorbox{bigidea}[1][]{
  enhanced,breakable,colback=softblue,colframe=deepblue,
  fonttitle=\bfseries,title=Big idea,#1}
\newtcolorbox{warningbox}[1][]{
  enhanced,breakable,colback=softred,colframe=deepred,
  fonttitle=\bfseries,title=Common trap,#1}
\newtcolorbox{recipebox}[1][]{
  enhanced,breakable,colback=softgreen,colframe=deepgreen,
  fonttitle=\bfseries,title=Working recipe,#1}
\newtcolorbox{intuitionbox}[1][]{
  enhanced,breakable,colback=softorange,colframe=deeporange,
  fonttitle=\bfseries,title=Intuition,#1}
\newtcolorbox{summarybox}[1][]{
  enhanced,breakable,colback=softgray,colframe=black!55,
  fonttitle=\bfseries,title=Chapter summary,#1}

\newcommand{\R}{\mathbb{R}}
\newcommand{\C}{\mathbb{C}}
\newcommand{\N}{\mathbb{N}}
\newcommand{\Z}{\mathbb{Z}}
\newcommand{\T}{\mathbb{T}}
\newcommand{\M}{\mathfrak{M}}
\newcommand{\m}{\mathfrak{m}}
\newcommand{\n}{\mathfrak{n}}
\newcommand{\g}{\mathfrak{g}}
\newcommand{\supp}{\operatorname{supp}}
\newcommand{\magop}{\operatorname{mag}}
\newcommand{\lc}{\operatorname{lc}}
\newcommand{\domop}{\operatorname{dom}}
\newcommand{\val}{\operatorname{v}}
\newcommand{\ord}{\operatorname{ord}}
\newcommand{\id}{\operatorname{id}}
\newcommand{\Ei}{\operatorname{Ei}}
\newcommand{\Ai}{\operatorname{Ai}}
\newcommand{\Bi}{\operatorname{Bi}}
\newcommand{\Borel}{\mathcal{B}}
\newcommand{\Laplace}{\mathcal{L}}
\newcommand{\asympseries}{\sim_{\!\mathrm{as}}}
\newcommand{\transmonomial}{\textit{transmonomial}}
\newcommand{\transseries}{\textit{transseries}}
\newcommand{\e}{\mathrm{e}}
\newcommand{\dd}{\,\mathrm{d}}
\newcommand{\eps}{\varepsilon}
\newcommand{\smallo}{o}
\newcommand{\bigO}{O}
\newcommand{\defeq}{\mathrel{:=}}
\newcommand{\eqdef}{\mathrel{=:}}
\newcommand{\precdom}{\prec}
\newcommand{\succdom}{\succ}
\newcommand{\precsimdom}{\preccurlyeq}
\newcommand{\succsimdom}{\succcurlyeq}
\newcommand{\compose}{\mathbin{\circ}}
\newcommand{\logit}[1]{\log^{\compose #1}}
\newcommand{\expit}[1]{\exp^{\compose #1}}
\newcommand{\sector}[1]{\Phi_{#1}}
\newcommand{\phihat}{\widehat{\Phi}}

\DeclarePairedDelimiter{\abs}{\lvert}{\rvert}
\DeclarePairedDelimiter{\norm}{\lVert}{\rVert}
\DeclarePairedDelimiter{\paren}{(}{)}
\DeclarePairedDelimiter{\braces}{\lbrace}{\rbrace}
\DeclarePairedDelimiter{\bracks}{[}{]}

\title{\Huge\bfseries Transseries for Mere Mortals\\[0.7em]
\Large From Asymptotic Scales to Resurgence}
\author{A step-by-step tutorial for a mathematically mature undergraduate}
\date{August 2026}

\begin{document}
\frontmatter
\hypersetup{pageanchor=false}
\maketitle
\hypersetup{pageanchor=true}

\thispagestyle{empty}
\begin{center}
\begin{minipage}{0.88\textwidth}
\small
\textbf{Scope.} This book develops two closely related meanings of the word
``transseries.'' First, a transseries is a formal generalized series built from
powers, exponentials, logarithms, and well-ordered infinite sums. Second, in
asymptotic analysis and mathematical physics, a resurgent transseries is an
organized family of usually divergent power-series sectors multiplied by
exponentially small scales. The first viewpoint explains the algebra; the second
explains how exponentially small information repairs divergent asymptotics and
Stokes jumps.

\medskip
\textbf{Prerequisites.} Single-variable calculus, elementary real analysis,
linear algebra, and comfort with formal power series. A first course in complex
analysis is useful only in the chapters on Borel summation and resurgence. Every
specialized construction is introduced from scratch.

\medskip
\textbf{Sources.} The exposition is original, but it was informed by the
materials supplied with the request, especially Edgar's survey
\cite{EdgarBeginners}, Costin's construction and summability survey
\cite{CostinTopological}, the primer of Aniceto--Ba\c{s}ar--Schiappa
\cite{AnicetoBasarSchiappa}, and papers on composition, surreal numbers, and
$H$-fields. Additional primary references are listed in the bibliography.
\end{minipage}
\end{center}
\clearpage

\chapter*{How to use this book}
\addcontentsline{toc}{chapter}{How to use this book}

There are three reasonable routes through the material.

\begin{description}[style=nextline,leftmargin=2em]
\item[The practical route]
Read Chapters~\ref{ch:asymptotic-language}--\ref{ch:first-transseries}, then
Chapters~\ref{ch:field-ops}--\ref{ch:fixed-points}, and finish with the worked
examples in Part~V. This is enough to begin calculating.

\item[The formal-algebra route]
Read Parts~I--IV in order. This gives the construction of logarithmic-exponential
transseries as a Hahn field, including the support conditions that make infinite
algebra legal.

\item[The analytic/resurgent route]
Read Part~I, the linear differential-equation example in
Chapter~\ref{ch:euler-ode}, and then Part~VI. This route explains factorial
divergence, Borel summation, Stokes phenomena, and the role of extra exponential
sectors.
\end{description}

The symbol $x$ always tends to $+\infty$ unless another limit is stated. We use
$f\precdom g$ for ``$f$ is asymptotically smaller than $g$,'' meaning
$f/g\to0$ in an analytic context and the corresponding monomial comparison in a
formal context. The two meanings agree whenever a formal transseries is realized
by an actual function.

\begin{bigidea}
A transseries is not merely a very long power series. It is a coordinate system
for asymptotic size. It can record, in one object, terms on scales such as
\[
  x^3,\quad x,\quad \log x,\quad 1,\quad x^{-1},\quad
  \e^{-x},\quad \e^{-x^2},\quad \e^{-\e^x},\ldots
\]
and it keeps them in a strict dominance order.
\end{bigidea}

\tableofcontents
\listoffigures

\chapter*{Notation at a glance}
\addcontentsline{toc}{chapter}{Notation at a glance}

\begin{longtable}{@{}>{$}l<{$} p{0.72\textwidth}@{}}
\toprule
\text{Symbol} & \text{Meaning}\\
\midrule
\endhead
f\sim g & $f/g\to1$.\\
f\precdom g & $f/g\to0$, or formally: every leading monomial of $f$ is smaller than the leading monomial of $g$.\\
f\succdom g & $g\precdom f$.\\
f=\bigO(g) & $\abs{f/g}$ is eventually bounded.\\
f=\smallo(g) & $f/g\to0$.\\
\M & An ordered multiplicative group of monomials.\\
\supp T & The set of monomials having nonzero coefficient in the formal series $T$.\\
\magop T & The largest monomial in $\supp T$.\\
\lc(T) & The coefficient of $\magop T$.\\
T^\dagger & The logarithmic derivative $T'/T$, when $T\ne0$.\\
\T & A standard field of logarithmic-exponential transseries.\\
\phihat & A formal, usually divergent, power series.\\
\Borel,\Laplace & Borel and Laplace transforms.\\
\bottomrule
\end{longtable}

\mainmatter

\part{Why asymptotics needs more than power series}

\chapter{The language of asymptotic approximation}
\label{ch:asymptotic-language}

\section{Exact formulas and asymptotic formulas answer different questions}

An exact formula says what a quantity is. An asymptotic formula says which parts
of it matter on a chosen limiting scale. These are different jobs.

For example,
\[
  \frac{x^2+3x+1}{x^2-x+2}
  = 1+\frac{4}{x}+\frac{3}{x^2}-\frac{5}{x^3}+\bigO(x^{-4})
  \qquad (x\to+\infty).
\]
The rational function on the left is exact. The expression on the right is more
useful when $x$ is large because it separates successive sizes. The first term is
of order $1$, the next of order $x^{-1}$, and so on.

\begin{definition}[Asymptotic equivalence]
For eventually nonzero functions $f$ and $g$, we write
\[
  f(x)\sim g(x) \qquad (x\to+\infty)
\]
when $f(x)/g(x)\to1$.
\end{definition}

\begin{definition}[Dominance]
We write $f\precdom g$ when $f/g\to0$, and $f\succdom g$ when $g\precdom f$.
Thus $x^{-2}\precdom x^{-1}\precdom1\precdom\log x\precdom x\precdom\e^x$.
\end{definition}

Dominance is more informative than ordinary inequalities. Although
$x^{-1}<2x^{-1}$, the two functions have the same asymptotic size. In contrast,
$x^{-2}\precdom x^{-1}$ says that $x^{-2}$ becomes negligible relative to
$x^{-1}$.

\section{Asymptotic scales}

\begin{definition}[Asymptotic scale]
A sequence $(\phi_n)_{n\ge0}$ is an asymptotic scale if
\[
  \phi_{n+1}\precdom\phi_n
  \qquad\text{for every }n.
\]
\end{definition}

The standard scale at infinity is
\[
  1,\ x^{-1},\ x^{-2},\ldots.
\]
But many problems naturally use other scales:
\[
  x^\alpha,\ x^{\alpha-1},\ldots;
  \qquad
  1,\frac1{\log x},\frac1{(\log x)^2},\ldots;
  \qquad
  1,\e^{-x},\e^{-2x},\ldots.
\]
A transseries allows several such scales to be interleaved and nested.

\begin{figure}[ht]
\centering
\begin{tikzpicture}[x=1.08cm,y=1cm,>=Latex]
  \draw[->,thick] (0,0)--(11.2,0) node[above left]{larger asymptotic size};
  \foreach \x/\lab in {0.4/{\e^{-\e^x}},2/{\e^{-x^2}},3.5/{\e^{-x}},5.1/{x^{-2}},6.2/{x^{-1}},7.2/{1},8.2/{\log x},9.2/{x},10.4/{\e^x}}{
    \draw (\x,0.12)--(\x,-0.12) node[below=5pt] {$\lab$};
  }
  \draw[decorate,decoration={brace,amplitude=5pt},deepblue]
    (0.35,0.55)--(3.55,0.55) node[midway,above=7pt]{beyond all algebraic orders};
  \draw[decorate,decoration={brace,amplitude=5pt},deepgreen]
    (5.05,0.55)--(9.25,0.55) node[midway,above=7pt]{power-logarithmic scales};
\end{tikzpicture}
\caption{A tiny part of the dominance hierarchy at $x\to+\infty$.}
\label{fig:growth-ladder}
\end{figure}

\section{Poincar\'e asymptotic expansions}

\begin{definition}[Asymptotic expansion]
Let $(\phi_n)$ be an asymptotic scale. We write
\[
  f(x)\asympseries \sum_{n=0}^{\infty} a_n\phi_n(x)
\]
if, for every $N\ge1$,
\[
  f(x)-\sum_{n=0}^{N-1}a_n\phi_n(x)=\smallo(\phi_{N-1}(x)).
\]
Equivalently, after subtracting the first $N$ terms, the remainder is smaller
than the last retained scale.
\end{definition}

The infinite sum in this definition need not converge. It is shorthand for an
infinite list of finite-accuracy statements.

\begin{example}[A geometric expansion]
For $x>1$,
\[
  \frac1{x+1}=\frac1x\frac1{1+x^{-1}}
  =\frac1x-\frac1{x^2}+\frac1{x^3}-\cdots.
\]
Here the series actually converges, but convergence is stronger than required.
As an asymptotic statement,
\[
  \frac1{x+1}-\sum_{n=1}^{N}(-1)^{n-1}x^{-n}
  =\frac{(-1)^N}{x^N(x+1)}=\bigO(x^{-N-1}).
\]
\end{example}

\begin{example}[A divergent asymptotic expansion]
The solution of
\[
  y'+y=\frac1x
\]
that behaves like $1/x$ has the formal expansion
\[
  y(x)\asympseries \frac1x+\frac1{x^2}+\frac{2!}{x^3}
       +\frac{3!}{x^4}+\cdots.
\]
The coefficients grow factorially, so the series diverges for every finite $x$.
It is nevertheless an accurate asymptotic expansion when truncated near its
smallest term. We will analyze it completely in
Chapter~\ref{ch:euler-ode}.
\end{example}

\section{Why the coefficients are unique}

\begin{proposition}[Uniqueness in a fixed scale]
If
\[
 f\asympseries\sum_{n\ge0}a_n\phi_n
 \qquad\text{and}\qquad
 f\asympseries\sum_{n\ge0}b_n\phi_n
\]
with respect to the same asymptotic scale, then $a_n=b_n$ for all $n$.
\end{proposition}

\begin{proof}
Subtract the two expansions. At order $N=1$,
\[
  (a_0-b_0)\phi_0=\smallo(\phi_0),
\]
so $a_0-b_0=0$. Suppose $a_0=b_0,\ldots,a_{N-1}=b_{N-1}$. At the next
order,
\[
  (a_N-b_N)\phi_N=\smallo(\phi_N),
\]
which again forces $a_N=b_N$. Induction completes the proof.
\end{proof}

The phrase ``in a fixed scale'' matters. A power series in $x^{-1}$ cannot see an
$\e^{-x}$ term at all. Enlarging the scale enlarges the information that can be
recorded.

\section{Truncation, optimal truncation, and error}

For a convergent series, retaining more terms eventually improves the answer.
For a divergent asymptotic series this is usually false. If terms first decrease
and then increase, the best elementary approximation is often obtained by
stopping near the smallest term.

Consider a typical Gevrey-1 term
\[
  \frac{n!}{x^{n+1}}.
\]
The ratio of consecutive terms is $(n+1)/x$. Terms decrease while $n+1<x$ and
increase after $n+1>x$. Thus the smallest term occurs near $n\approx x$, and its
size is roughly $\e^{-x}$ by Stirling's formula. The error left after optimal
truncation is therefore exponentially small. This is the first hint that an
$\e^{-x}$ sector is not decorative: it lives exactly at the scale of the best
possible remainder of the power-series sector.

\begin{summarybox}
An asymptotic expansion is an ordered sequence of approximations, not
necessarily a convergent sum. Its coefficients are unique only relative to a
chosen scale. Divergent series can be highly useful, and their optimal remainder
often lies on a new exponentially small scale.
\end{summarybox}

\chapter{Beyond all orders: the information power series cannot see}
\label{ch:beyond-all-orders}

\section{Flat functions}

\begin{definition}[Flat relative to powers]
A function $f$ is flat at $+\infty$ relative to the power scale if
\[
  f(x)=\smallo(x^{-N})
  \qquad\text{for every }N\ge0.
\]
\end{definition}

\begin{proposition}
For every $a>0$ and every real $N$, one has
\[
  \e^{-ax}\precdom x^{-N}.
\]
\end{proposition}

\begin{proof}
The claim is equivalent to $x^N\e^{-ax}\to0$. Repeated use of l'H\^opital's
rule proves this for nonnegative integers $N$; a larger integer bounds any real
$N$. Alternatively, write $x^N\e^{-ax}=\exp(N\log x-ax)$ and observe that
$N\log x-ax\to-\infty$.
\end{proof}

Therefore the zero power series is an asymptotic expansion of $\e^{-x}$:
\[
  \e^{-x}\asympseries 0+0x^{-1}+0x^{-2}+\cdots.
\]
The same is true of $17\e^{-x}$, $\e^{-2x}$, and
$\e^{-x}(1+x^{-1}+x^{-2}+\cdots)$ whenever the latter expression is interpreted
as an asymptotic expansion.

\begin{warningbox}
The statement ``two functions have the same asymptotic power series'' does not
mean that their difference is zero. It means only that their difference is
smaller than every power retained by that scale.
\end{warningbox}

\section{An entire family hidden behind one power series}

Suppose $f$ has the power-series expansion
\[
  f(x)\asympseries\sum_{n\ge0}a_nx^{-n}.
\]
Then for every constant $C$,
\[
  f(x)+C\e^{-x}\asympseries\sum_{n\ge0}a_nx^{-n}.
\]
Thus a power series can represent an entire one-parameter family. In a
differential equation that parameter is often an integration constant.

This phenomenon is not rare. A differential equation of order $r$ normally has
$r$ free constants, but a purely algebraic asymptotic expansion may appear
unique. The ``missing constants'' often multiply exponentially small solutions
of the linearized equation. Transseries restore those sectors.

\section{A toy implicit equation that creates exponentials}

Consider
\begin{equation}
  y=x+\e^{-y}.
  \label{eq:toy-implicit}
\end{equation}
If we seek only a power series in $x^{-1}$, the correction $y-x$ is invisible:
it is smaller than every power. Set $\delta=y-x$. Then
\[
  \delta=\e^{-x}\e^{-\delta}.
\]
This already suggests a series in powers of the small transmonomial $\e^{-x}$:
\[
  \delta=c_1\e^{-x}+c_2\e^{-2x}+c_3\e^{-3x}+\cdots.
\]
In Chapter~\ref{ch:implicit-exp} we will calculate
\[
  y=x+\e^{-x}-\e^{-2x}+\frac32\e^{-3x}
     -\frac83\e^{-4x}+\frac{125}{24}\e^{-5x}+\cdots.
\]
A power series says only $y\asympseries x$. The transseries resolves the whole
exponential correction.

\section{Nested smallness}

Exponential scales themselves have scales beneath them:
\[
  \e^{-x^2}\precdom\e^{-x}\precdom x^{-N}
  \qquad\text{for every fixed }N,
\]
and
\[
  \e^{-\e^x}\precdom\e^{-x^2}.
\]
Inside an $\e^{-x}$ sector we may also have an ordinary inverse-power series:
\[
  \e^{-x}\left(1+\frac{2}{x}+\frac{7}{x^2}+\cdots\right).
\]
This is the characteristic nested structure of a transseries: a series whose
``coefficients'' may themselves be series on a smaller scale.

\section{Why ``nonperturbative'' does not mean ``mysterious''}

In physics, the inverse-power sector is often called perturbative and terms such
as $\e^{-Ax}$ are called nonperturbative. Mathematically, both are formal pieces
of one transseries. The terminology describes how the terms arise in a given
approximation method, not a difference in logical status.

\begin{intuitionbox}
A microscope with magnification $x^N$ cannot distinguish two functions whose
difference is $\e^{-x}$. A transseries changes microscopes. After factoring out
$\e^{-x}$, it asks for the next layer of detail, perhaps a power series in
$x^{-1}$. Then it can zoom again to $\e^{-2x}$, and so on.
\end{intuitionbox}

\begin{summarybox}
Exponentially small terms are invisible to every finite algebraic order, yet
they often contain integration constants and the optimal remainder of divergent
series. Transseries enlarge the asymptotic scale so that this information becomes
visible and computable.
\end{summarybox}

\chapter{Your first transseries}
\label{ch:first-transseries}

\section{The standard sector form}

The most common transseries in applications has the form
\begin{equation}
  \Phi(x;\sigma)
  =\sum_{k=0}^{\infty}
     \sigma^k\e^{-kAx}x^{k\beta}\sector{k}(x),
  \qquad
  \sector{k}(x)=\sum_{n=0}^{\infty}a_{k,n}x^{-n},
  \label{eq:one-parameter-transseries}
\end{equation}
where $A>0$, $\beta\in\R$, and $\sigma$ is a formal parameter.

Read this expression from outside inward.

\begin{enumerate}
\item The index $k$ labels exponential sectors.
\item Neighboring sectors differ by the small ratio
      $\e^{-Ax}x^\beta$.
\item Inside each sector is an inverse-power series.
\item The parameter $\sigma$ usually becomes an integration constant or boundary
      datum after an analytic summation prescription is chosen.
\end{enumerate}

The $k=0$ sector is
\[
  \sector{0}(x)=a_{0,0}+a_{0,1}x^{-1}+\cdots.
\]
The $k=1$ sector is exponentially smaller:
\[
  \sigma\e^{-Ax}x^\beta
  \left(a_{1,0}+a_{1,1}x^{-1}+\cdots\right),
\]
and the $k=2$ sector is smaller again by another factor
$\sigma\e^{-Ax}x^\beta$.

\begin{figure}[ht]
\centering
\begin{tikzpicture}[>=Latex,node distance=7mm]
  \foreach \k in {0,...,4}{
    \node[draw,rounded corners,minimum width=2.2cm,minimum height=1.0cm,
      fill=softblue] (s\k) at (2.55*\k,0)
      {$\sigma^{\k}\e^{-\k Ax}x^{\k\beta}\Phi_{\k}(x)$};
    \ifnum\k>0
      \pgfmathtruncatemacro{\prev}{\k-1}
      \draw[->,thick] (s\prev) -- node[above]
      {$\times\,\sigma\e^{-Ax}x^\beta$} (s\k);
    \fi
  }
  \node[below=8mm of s0,align=center] {power-series\\sector};
  \node[below=8mm of s1,align=center] {one-exponential\\sector};
  \node[below=8mm of s2,align=center] {two-exponential\\sector};
\end{tikzpicture}
\caption{A one-parameter transseries arranged by exponentially decreasing sectors.}
\label{fig:sector-chain}
\end{figure}

\section{A two-dimensional coefficient array}

It is useful to picture the coefficients $a_{k,n}$ on a lattice. Moving right
increases $n$ and multiplies the scale by $x^{-1}$. Moving upward increases $k$
and multiplies it by $\e^{-Ax}x^\beta$.

\begin{center}
\begin{tikzpicture}[x=1.35cm,y=0.9cm,>=Latex]
  \draw[->] (-0.4,0)--(6.2,0) node[right]{$n$ (inverse-power order)};
  \draw[->] (0,-0.4)--(0,4.8) node[above]{$k$ (exponential sector)};
  \foreach \k in {0,...,4}{
    \foreach \nn in {0,...,5}{
      \fill (\nn,\k) circle (1.8pt);
    }
  }
  \draw[->,deepblue,thick] (1,1.18)--(2,1.18)
    node[midway,above]{$\times x^{-1}$};
  \draw[->,deepgreen,thick] (4.2,1)--(4.2,2)
    node[midway,right]{$\times\e^{-Ax}x^\beta$};
  \node[below left] at (0,0) {$a_{0,0}$};
\end{tikzpicture}
\end{center}

A finite approximation amounts to selecting a finite region of this lattice.
But ``rectangular truncation'' is not always optimal: because the exponential
step is much smaller than any fixed power step, one often computes many power
terms in the first sector before including even the leading term of the next.

\section{Three meanings of the formula}

Expression~\eqref{eq:one-parameter-transseries} can be used in three distinct
ways.

\subsection{As a purely formal object}

The symbols $x^{-1}$ and $\e^{-Ax}$ are ordered monomials. Infinite sums are
legal because their supports satisfy a well-ordering condition. No numerical
value of $x$ is substituted. This is the Hahn-field viewpoint developed in
Parts~II--IV.

\subsection{As an asymptotic specification}

Each sector may specify the asymptotic behavior of an actual function, with a
rule for how remainders compare to omitted terms. This is the viewpoint of
classical asymptotic analysis.

\subsection{As resurgent data}

The inner series $\sector{k}$ may diverge factorially. Borel--Laplace summation
then turns each formal sector into an analytic function in a direction of the
complex plane. Singularities of the Borel transform connect different values of
$k$, and crossing a Stokes direction changes the effective value of $\sigma$.
This is the resurgent viewpoint of Part~VI.

\begin{warningbox}[title=Do not silently identify the viewpoints]
A formal logarithmic-exponential transseries need not be a convergent series of
real functions. Conversely, a resurgent transseries may use complex directions,
branch cuts, and factorially divergent coefficient series not captured by a
naive notion of convergence. The viewpoints overlap, but analytic realization
requires extra hypotheses.
\end{warningbox}

\section{Multi-parameter transseries}

An $r$-parameter version is
\[
 \Phi(x;\bm\sigma)
 =\sum_{\bm k\in\N^r}
   \bm\sigma^{\bm k}
   \e^{-(\bm k\cdot\bm A)x}
   x^{\bm k\cdot\bm\beta}
   \Phi_{\bm k}(x),
\]
where
\[
  \bm\sigma^{\bm k}=\sigma_1^{k_1}\cdots\sigma_r^{k_r}.
\]
Different exponential actions $A_j$ usually correspond to different
independent solutions of a linearized equation, different saddle points of an
integral, or different instanton types. When integer relations among the $A_j$
make sectors collide, logarithms may appear; this is called resonance.

\section{A first substitution calculation}

Let
\[
  T=x+\e^{-x}+2x^{-1}\e^{-x}+3\e^{-2x}.
\]
Then, to order $\e^{-2x}$,
\begin{align*}
  \e^{-T}
  &=\e^{-x}\exp\left(-\e^{-x}-2x^{-1}\e^{-x}-3\e^{-2x}\right)\\
  &=\e^{-x}\left[
      1-(1+2x^{-1})\e^{-x}+\bigO(\e^{-2x})
    \right]\\
  &=\e^{-x}-(1+2x^{-1})\e^{-2x}+\bigO(\e^{-3x}).
\end{align*}
The computation looks exactly like ordinary formal power-series algebra, except
that the small expansion variable is itself a monomial, $\e^{-x}$, and its
coefficients may involve powers of $x^{-1}$.

\section{What the rest of the book must justify}

We have manipulated transseries informally. A rigorous theory must answer:

\begin{enumerate}
\item Which monomials are allowed?
\item Which infinite supports are allowed?
\item Why does the coefficient of a product involve only finitely many terms?
\item How are order, inverse, exponential, logarithm, derivative, integral, and
      composition defined?
\item When does a formal transseries correspond to an actual function?
\end{enumerate}

Parts~II--IV answer the first four questions. Part~VI addresses the fifth for a
large and important analytic class.

\begin{summarybox}
The most familiar transseries is a sum of exponentially graded sectors, each
carrying an inverse-power series. It may be read formally, asymptotically, or
resurgently. The formal theory must control supports; the analytic theory must
control divergence and Stokes jumps.
\end{summarybox}


\part{Generalized series and ordered supports}

\chapter{Growth rates as monomials}
\label{ch:monomials}

\section{From numerical size to asymptotic size}

Ordinary power series are built from the monomials
$1,z,z^2,\ldots$. At $x\to+\infty$ we want a much larger alphabet:
\[
  x^a,\qquad \e^{x},\qquad \e^{-x^2},\qquad
  x^{\sqrt2}\e^{-3x},\qquad
  \exp\!\left(\frac{x}{\log x}\right),\qquad\ldots
\]
The basic objects are not numerical values but formal growth rates.

\begin{definition}[Ordered monomial group]
An ordered monomial group is an abelian group $\M$ written multiplicatively,
with identity $1$, together with a total order $\succdom$ compatible with
multiplication:
\[
  \m\succdom\n \quad\Longrightarrow\quad
  \m\g\succdom\n\g
  \qquad(\g\in\M).
\]
\end{definition}

The intended reading is ``$\m$ is asymptotically larger than $\n$.'' The inverse
of a large monomial is small:
\[
  \m\succdom1 \quad\Longleftrightarrow\quad \m^{-1}\precdom1.
\]

\begin{example}[Power-exponential monomials]
For
\[
  \m=x^a\e^{bx},\qquad \n=x^c\e^{dx},
\]
we declare $\m\succdom\n$ if either $b>d$, or $b=d$ and $a>c$. This is just the
asymptotic fact
\[
  \frac{x^a\e^{bx}}{x^c\e^{dx}}
  =x^{a-c}\e^{(b-d)x}.
\]
The exponential comparison is decided first; powers break ties.
\end{example}

\section{Monomials are coordinates, not functions}

A formal monomial may be inspired by a function, but the theory initially uses
only its algebra and order. This separation is valuable. It lets us manipulate
objects before proving convergence or analytic realization.

For example, we may declare
\[
  \e^{-\e^x}\precdom \e^{-x^2}\precdom\e^{-x}\precdom x^{-N}
\]
for every fixed $N$, and use these symbols as ordered coordinates. Later, when
$x$ is a large real number, the same order is analytically correct.

\begin{intuitionbox}
Think of a monomial as a labeled ruler for one asymptotic scale. A transseries is
a coordinate vector with respect to an enormous, ordered collection of such
rulers.
\end{intuitionbox}

\section{Leading monomials determine order}

Suppose we have a formal sum
\[
  T=3\e^x-7x^4+2\log x+x^{-1}+\e^{-x}.
\]
The largest monomial is $\e^x$, so the leading term is $3\e^x$ and $T$ is
positive. The smaller terms cannot overturn the sign, just as the first nonzero
digit determines the sign of a decimal written in positional notation.

\begin{definition}[Magnitude, leading coefficient, leading term]
For a nonzero formal series $T$, if its support has a largest monomial $\m$, set
\[
  \magop(T)=\m,\qquad
  \lc(T)=[\m]T,\qquad
  \domop(T)=\lc(T)\m.
\]
Here $[\m]T$ denotes the coefficient of $\m$.
\end{definition}

We then define
\[
  T>0 \quad\Longleftrightarrow\quad \lc(T)>0.
\]
This order is lexicographic by asymptotic size.

\begin{definition}[Asymptotic comparison of series]
For nonzero series $S,T$:
\[
  S\precdom T \quad\Longleftrightarrow\quad
  \magop(S)\precdom\magop(T),
\]
\[
  S\asymp T \quad\Longleftrightarrow\quad
  \magop(S)=\magop(T),
\]
and
\[
  S\sim T \quad\Longleftrightarrow\quad
  S-T\precdom T.
\]
The last relation means that $S$ and $T$ have the same leading term.
\end{definition}

\section{Large, small, and purely large parts}

A series $T$ is called
\begin{itemize}
\item \emph{large} if $\magop(T)\succdom1$;
\item \emph{small} if every monomial in its support is $\precdom1$;
\item \emph{purely large} if every monomial in its support is $\succdom1$.
\end{itemize}
The zero series is conventionally both small and purely large, because its
support is empty.

Every transseries will have a canonical decomposition
\begin{equation}
  T=T_{\mathrm{pl}}+T_1+T_{\mathrm{sm}},
  \label{eq:canonical-decomposition-preview}
\end{equation}
where $T_{\mathrm{pl}}$ is purely large, $T_1\in\R$ is constant, and
$T_{\mathrm{sm}}$ is small. This decomposition is the key to defining
$\exp T$.

\section{Valuation notation}

Some authors reverse the order and use an additive valuation. If monomials are
written as $t^\gamma$, then smaller $\gamma$ may correspond to larger
asymptotic size. A valuation $\val(T)$ records the exponent of the leading
monomial and satisfies
\[
  \val(ST)=\val(S)+\val(T),\qquad
  \val(S+T)\ge\min\{\val(S),\val(T)\}.
\]
Cancellation can make the second inequality strict.

In this tutorial we mostly use the more visual multiplicative notation
$\magop(T)$. The two languages contain the same information:
\[
  S\precdom T
  \quad\Longleftrightarrow\quad
  \val(S)>\val(T)
\]
under the usual valuation convention.

\section{The field will be non-Archimedean}

In the ordered field of transseries, $x$ is larger than every real constant:
\[
  x>n\qquad(n\in\N),
\]
and $x^{-1}$ is a positive infinitesimal:
\[
  0<x^{-1}<\frac1n\qquad(n\ge1).
\]
Thus the Archimedean axiom of the real numbers fails. This is not a defect; it is
exactly what lets one manipulate infinity and infinitesimals algebraically.

\begin{summarybox}
The monomial group packages asymptotic growth rates into an ordered algebra.
Once a series has a largest monomial, that monomial and its coefficient determine
its leading size and sign. This turns asymptotic comparison into ordinary field
order.
\end{summarybox}

\chapter{Hahn series: how infinite algebra is made legal}
\label{ch:hahn}

\section{Why arbitrary supports do not work}

Let $\M$ be an ordered monomial group. A completely arbitrary formal expression
\[
  \sum_{\m\in\M}a_{\m}\m
\]
is too large a class. It may have no leading term, and multiplication may ask us
to add infinitely many unrelated coefficients.

For expansions at infinity, the series
\[
  1+x^{-1}+x^{-2}+\cdots
\]
behaves well: its monomials march downward from a largest one. In contrast,
\[
  1+x+x^2+x^3+\cdots
\]
has no largest monomial. It is a perfectly good power series near $x=0$, but it
is not a well-based asymptotic series at $x=+\infty$.

\begin{definition}[Well-based subset]
A subset $A\subseteq\M$ is well-based if every nonempty subset of $A$ has a
largest element with respect to $\succdom$.
\end{definition}

Equivalently, $A$ contains no infinite chain
\[
  \m_0\precdom\m_1\precdom\m_2\precdom\cdots
\]
of monomials growing without a largest member. It may contain infinite
descending chains
\[
  \m_0\succdom\m_1\succdom\m_2\succdom\cdots,
\]
which is exactly how an asymptotic expansion is written.

\begin{definition}[Hahn series]
A Hahn series over a coefficient field $K$ and monomial group $\M$ is a formal
sum
\[
  T=\sum_{\m\in\M}a_{\m}\m,
  \qquad a_{\m}\in K,
\]
whose support
\[
  \supp T=\{\m\in\M:a_{\m}\ne0\}
\]
is well-based. The collection is denoted $K((\M))$ or, in another common
notation, $K[[\M]]$.
\end{definition}

\section{Examples and nonexamples}

Assume
\[
  \e^{-x}\precdom x^{-N}\precdom1\precdom x^N\precdom\e^x
\]
for every $N>0$.

\begin{center}
\begin{tabularx}{0.94\textwidth}{@{}X X@{}}
\toprule
\textbf{Well-based} & \textbf{Not well-based}\\
\midrule
$\displaystyle\sum_{n\ge0}n!x^{-n}$ &
$\displaystyle\sum_{n\ge0}x^n$\\[0.7em]
$\displaystyle\sum_{n\ge0}\e^{-nx}$ &
$\displaystyle\sum_{n\ge0}\e^{nx}$\\[0.7em]
$\displaystyle\sum_{n\ge0}x^{-n}\e^{-n^2x}$ &
$\displaystyle\sum_{n\ge1}x^{-1/n}$\\[0.7em]
\bottomrule
\end{tabularx}
\end{center}

The last nonexample deserves attention. The monomials
$x^{-1},x^{-1/2},x^{-1/3},\ldots$ become progressively larger and approach $1$
without reaching it. The set has no largest element.

\section{Addition}

Addition is coefficientwise:
\[
  [\m](S+T)=[\m]S+[\m]T.
\]
The union of two well-based sets is well-based, so the support of $S+T$ is
legal. Cancellation may remove the old leading term and reveal a smaller one.

\begin{example}
If
\[
 S=x+3+x^{-1},\qquad T=-x+2-x^{-2},
\]
then
\[
 S+T=5+x^{-1}-x^{-2}.
\]
The leading $x$ terms cancel, and the magnitude drops from $x$ to $1$.
\end{example}

\section{Multiplication and Neumann's lemma}

Formally, multiplication should be convolution:
\begin{equation}
  [\g](ST)=\sum_{\m\n=\g}[\m]S\,[\n]T.
  \label{eq:hahn-convolution}
\end{equation}
Two facts are required:
\begin{enumerate}
\item only finitely many pairs $(\m,\n)$ contribute to a fixed $\g$;
\item the set of products of support monomials is again well-based.
\end{enumerate}
Both follow from a standard ordered-group result usually called Neumann's
lemma \cite{Hahn1907,Neumann1949,vanDerHoevenBook}.

\begin{theorem}[Neumann's lemma, two-set form]
Let $A,B\subseteq\M$ be well-based. Then
\[
  AB=\{ab:a\in A,\ b\in B\}
\]
is well-based, and every $\g\in\M$ has only finitely many representations
$\g=ab$ with $a\in A$ and $b\in B$.
\end{theorem}

\begin{proof}[Proof idea]
If $AB$ had an infinite sequence moving upward with no largest element, choose
corresponding products $a_nb_n$. By repeatedly passing to subsequences and using
the well-basedness of $A$ and $B$, one would force one factor sequence to move
upward indefinitely, contradicting the hypothesis. A similar diagonal argument
shows that infinitely many decompositions of one fixed $\g$ would force an
upward sequence in one factor and the inverse upward sequence in the other. The
full proof is an ordered-group version of the familiar fact that a coefficient
of a product of ordinary power series receives only finitely many contributions.
\end{proof}

The theorem makes \eqref{eq:hahn-convolution} an ordinary finite sum. Hence Hahn
series form a ring.

\section{Summable families}

We often need to sum a family $(T_i)_{i\in I}$ rather than a single sequence.

\begin{definition}[Summable family]
A family $(T_i)_{i\in I}$ of Hahn series is summable if
\begin{enumerate}
\item $\bigcup_i\supp T_i$ is well-based;
\item for each monomial $\m$, only finitely many $i$ have
      $[\m]T_i\ne0$.
\end{enumerate}
Then
\[
  \sum_{i\in I}T_i
\]
is defined coefficientwise.
\end{definition}

This is formal local finiteness. It does not ask for absolute convergence of
numbers. It asks that every coordinate receive only finitely many contributions.

\begin{proposition}
If $U\precdom1$, then the family $(U^n)_{n\ge0}$ is summable.
\end{proposition}

\begin{proof}[Idea]
All monomials in $U$ are small. Products of $n$ such monomials move farther down
the dominance order as $n$ grows. Neumann's lemma implies that the union of all
finite products is well-based and that a fixed monomial can occur at only
finitely many total degrees.
\end{proof}

This proposition is the formal engine behind geometric, exponential, and
logarithmic power series.

\section{Inverses}

Let $T\ne0$. Factor out its leading term:
\[
  T=a\m(1+U),
  \qquad a\in K^\times,\quad U\precdom1.
\]
Then define
\begin{equation}
  T^{-1}=a^{-1}\m^{-1}\sum_{n=0}^{\infty}(-U)^n.
  \label{eq:hahn-inverse}
\end{equation}
The family is summable because $U$ is small, and formal multiplication gives
$(1+U)\sum_{n\ge0}(-U)^n=1$.

\begin{example}
At infinity,
\begin{align*}
  (x+\log x+1)^{-1}
  &=x^{-1}\left(1+\frac{\log x+1}{x}\right)^{-1}\\
  &=x^{-1}-\frac{\log x+1}{x^2}
    +\frac{(\log x+1)^2}{x^3}-\cdots.
\end{align*}
\end{example}

Thus, once the monomial group and support condition are fixed, Hahn series form
a field.

\section{Order and the non-Archimedean triangle law}

The leading coefficient defines an order compatible with addition and
multiplication. The magnitude obeys
\[
  \magop(ST)=\magop(S)\magop(T),
\]
and
\[
  \magop(S+T)\precsimdom\max\{\magop(S),\magop(T)\},
\]
with strict inequality only when leading terms cancel. This is the
non-Archimedean analogue of a triangle inequality.

\begin{bigidea}
Well-basedness replaces numerical convergence. It guarantees a leading term,
legal multiplication, and legal substitution into ordinary formal power series
whenever the substituted object is small.
\end{bigidea}

\begin{summarybox}
A Hahn series is a generalized power series with well-based support. Neumann's
lemma ensures that products and infinite power-series substitutions are locally
finite. Factoring out the leading term reduces inversion to a geometric series.
\end{summarybox}

\chapter{Grids and ratio sets}
\label{ch:grids}

\section{Why finite generators are useful}

The full Hahn field allows arbitrary well-based supports. For explicit
calculation it is often convenient to work in a smaller class generated by a
finite list of small monomials.

\begin{definition}[Ratio set]
A ratio set is a finite tuple
\[
  \boldsymbol\mu=(\mu_1,\ldots,\mu_r)
\]
of monomials with $\mu_j\precdom1$.
For $\bm k=(k_1,\ldots,k_r)\in\Z^r$, write
\[
  \boldsymbol\mu^{\bm k}
  =\mu_1^{k_1}\cdots\mu_r^{k_r}.
\]
\end{definition}

\begin{definition}[Grid]
Given a monomial $\m$ and a lower bound $\bm k_0\in\Z^r$, the grid
\[
  \mathfrak{J}(\m,\boldsymbol\mu,\bm k_0)
  =\left\{\m\boldsymbol\mu^{\bm k}:\bm k\ge\bm k_0\right\}
\]
uses componentwise inequality. A grid-based series has support contained in
some such grid.
\end{definition}

Usually we absorb $\bm k_0$ into $\m$ and write
\[
  \mathfrak{J}=\{\m\boldsymbol\mu^{\bm n}:\bm n\in\N^r\}.
\]

\begin{example}
Take
\[
  \mu_1=x^{-1},\qquad \mu_2=\e^{-x}.
\]
Then a grid contains monomials
\[
  \m x^{-p}\e^{-qx},\qquad p,q\in\N.
\]
This is exactly the two-dimensional sector lattice from
Chapter~\ref{ch:first-transseries}.
\end{example}

\section{The exponent lattice}

A grid converts monomial bookkeeping into integer-vector bookkeeping.
Multiplication becomes addition of exponent vectors:
\[
  \boldsymbol\mu^{\bm j}\boldsymbol\mu^{\bm k}
  =\boldsymbol\mu^{\bm j+\bm k}.
\]

\begin{figure}[ht]
\centering
\begin{tikzpicture}[x=1.05cm,y=0.85cm,>=Latex]
  \draw[->] (-0.4,0)--(6.2,0) node[right]{$p$ in $x^{-p}$};
  \draw[->] (0,-0.4)--(0,5.2) node[above]{$q$ in $\e^{-qx}$};
  \foreach \p in {0,...,5}{
    \foreach \q in {0,...,4}{
      \fill[deepblue] (\p,\q) circle (1.7pt);
    }
  }
  \draw[very thick,rounded corners,deepgreen]
    (0.75,0.75) rectangle (5.25,4.25);
  \node[deepgreen,above right] at (0.75,4.25)
    {a translated quadrant};
  \draw[->,thick] (1,1.3)--(2,1.3) node[midway,above]{$\mu_1$};
  \draw[->,thick] (4.3,1)--(4.3,2) node[midway,right]{$\mu_2$};
\end{tikzpicture}
\caption{A grid generated by $x^{-1}$ and $\e^{-x}$. Moving right or up makes the monomial smaller.}
\label{fig:grid-lattice}
\end{figure}

\section{Dickson's lemma}

The key combinatorial fact is that $\N^r$ is well-partially-ordered under
componentwise comparison.

\begin{theorem}[Dickson's lemma]
Every subset $E\subseteq\N^r$ has finitely many minimal elements. Equivalently,
every infinite sequence in $\N^r$ contains indices $i<j$ with
$\bm k_i\le\bm k_j$ componentwise.
\end{theorem}

\begin{proof}
We prove the sequence form by induction on $r$. It is obvious for $r=1$. For
$r>1$, consider the first coordinates. If some value occurs infinitely often,
restrict to that subsequence and apply induction to the remaining $r-1$
coordinates. Otherwise, pass to a subsequence with strictly increasing first
coordinate, and again use induction on the remaining coordinates. The resulting
pair is nondecreasing in every coordinate.

If a set had infinitely many minimal elements, listing them would produce a
comparable pair, contradicting minimality. Conversely, every element lies above
a minimal element; otherwise one could descend forever in the sum of the
coordinates.
\end{proof}

Dickson's lemma explains why finite ratio sets are algorithmically pleasant. It
is also a basic reason for termination in monomial-ideal and Gr\"obner-basis
arguments.

\section{Grid-based supports are well-based}

Because every generator $\mu_j$ is small, increasing an exponent vector makes
the corresponding monomial smaller. Dickson's lemma prevents a subset of a grid
from having infinitely many unrelated candidates for a leading monomial.

\begin{proposition}
Every subset of a grid generated by finitely many small monomials is well-based.
\end{proposition}

\begin{proof}[Idea]
Translate the exponent set into a subset of $\N^r$. Its finitely many minimal
vectors correspond to finitely many largest candidate monomials. Since the
monomial order is total, one of these candidates is largest. The same argument
applies to every nonempty subset.
\end{proof}

\section{Relations among generators}

The map
\[
  \bm k\longmapsto\boldsymbol\mu^{\bm k}
\]
need not be injective. For example, if $\mu_2=\mu_1^2$, then
$\mu_2=\mu_1^2$ gives two exponent descriptions of the same monomial. Formal
algorithms may either
\begin{itemize}
\item choose multiplicatively independent generators when possible;
\item retain relations and combine coefficients during normalization;
\item work with ``witnesses'' that certify the required dominance statements.
\end{itemize}
The mathematics does not require unique exponent vectors, but implementations
must account for collisions.

\section{Grid-based versus well-based transseries}

Grid-based series form a proper subfield of the corresponding well-based Hahn
field. The distinction is methodological:

\begin{center}
\begin{tabularx}{0.94\textwidth}{@{}p{0.22\textwidth}XX@{}}
\toprule
& \textbf{Grid-based} & \textbf{Well-based}\\
\midrule
Support & Contained in a finitely generated translated quadrant & Any well-based subset\\
Algorithms & Finite exponent-vector representation & More flexible but harder to encode\\
Typical applications & Explicit ODE and asymptotic calculations & Structural theorems and maximal constructions\\
Main caution & May need to enlarge the ratio set & Some operations require subtler support proofs\\
\bottomrule
\end{tabularx}
\end{center}

Edgar's introductory treatment emphasizes grid-based transseries
\cite{EdgarBeginners}; van der Hoeven and later structural treatments often use
well-based variants \cite{vanDerHoevenBook,ADHBook}. We will state most formulas
in a way that applies to both, while using grids for concrete computations.

\section{A practical truncation order}

Suppose a calculation uses $\mu_1=x^{-1}$ and $\mu_2=\e^{-x}$. A request such as
``keep through $x^{-4}$ in the perturbative sector and through $x^{-2}\e^{-x}$ in
the first exponential sector'' describes a finite staircase in $\N^2$.

A good symbolic engine should not assign a single total degree $p+q$ blindly,
because one factor of $\e^{-x}$ is smaller than every fixed power of $x^{-1}$.
Instead, it should use the actual monomial order and stop at a target monomial.

\begin{recipebox}
To organize a hand calculation:
\begin{enumerate}
\item Choose a finite ratio set of genuinely small monomials.
\item Write every retained term as a leading monomial times a nonnegative power
      of the ratios.
\item Decide a target monomial below which terms will be discarded.
\item After each operation, normalize collisions and discard terms smaller than
      the target.
\end{enumerate}
\end{recipebox}

\begin{summarybox}
A grid is a translated copy of $\N^r$ in the monomial group. Dickson's lemma
ensures that grid supports are well-based. Grids turn transseries algebra into
finite-dimensional exponent bookkeeping and are especially useful for symbolic
computation.
\end{summarybox}


\part{Constructing logarithmic-exponential transseries}

\chapter{The log-free hierarchy}
\label{ch:logfree-hierarchy}

\section{Stage zero: real powers of \texorpdfstring{$x$}{x}}

We now choose a particular monomial group. At height zero, take
\[
  \M_0=\{x^a:a\in\R\},
\]
with multiplication $x^ax^b=x^{a+b}$ and order
\[
  x^a\succdom x^b \quad\Longleftrightarrow\quad a>b.
\]
Let
\[
  \T_0=\R((\M_0))
\]
be the corresponding Hahn field, or its grid-based subfield if one wants finite
ratio sets.

Elements of $\T_0$ are generalized power series such as
\[
  7x^{\sqrt2}-3x^{1/3}+5+\sum_{n\ge1}a_nx^{-n/2}.
\]
The exponents need not be integers, but the support must be well-based.

\begin{example}
The series
\[
  \sum_{m,n\ge1}x^{-m-n\sqrt2}
\]
is legal and grid-based, generated by $x^{-1}$ and $x^{-\sqrt2}$. Its
coefficients need not have a simple closed form, but its support is contained in
a finitely generated grid.
\end{example}

\section{Stage one: exponentials of height-zero giants}

A purely large $L\in\T_0$ may serve as an exponential exponent. We introduce
new monomials
\[
  x^a\e^L,
  \qquad a\in\R,
  \quad L\in\T_0\text{ purely large}.
\]
They multiply by adding exponents:
\[
  (x^a\e^L)(x^b\e^K)=x^{a+b}\e^{L+K}.
\]
Their order is lexicographic:
\[
  x^a\e^L\succdom x^b\e^K
\]
if $L>K$, or if $L=K$ and $a>b$.

Why compare $L$ and $K$ first? Because
\[
  \frac{x^a\e^L}{x^b\e^K}=x^{a-b}\e^{L-K},
\]
and a nonzero purely large leading part of $L-K$ beats every power of $x$.

Let $\M_1$ be this enlarged monomial group and
\[
  \T_1=\R((\M_1)).
\]
Examples include
\[
  \e^{-x^2+3x},\qquad x^{7/3}\e^{2x},\qquad
  x^4+\e^{-x^{3/2}},\qquad
  \sum_{n\ge0}x^{-n}\e^{-x}.
\]

\section{Higher stages}

Repeat the construction. Assuming $\T_n$ has been defined, set
\begin{equation}
  \M_{n+1}
  =\left\{x^a\e^L:
      a\in\R,\ L\in\T_n\text{ purely large}
    \right\},
  \label{eq:monomial-height-recursion}
\end{equation}
with the same multiplication and lexicographic order, and let
\[
  \T_{n+1}=\R((\M_{n+1})).
\]
The inclusions
\[
  \M_0\subseteq\M_1\subseteq\M_2\subseteq\cdots,
  \qquad
  \T_0\subseteq\T_1\subseteq\T_2\subseteq\cdots
\]
are natural because $x^a=x^a\e^0$.

\begin{definition}[Exponential height]
A log-free transseries has height at most $n$ if it lies in $\T_n$. Its exact
height is the smallest such $n$.
\end{definition}

Height counts nested exponential complexity, not the number of exponential
symbols after simplification.

\begin{center}
\begin{tabular}{@{}ll@{}}
\toprule
Expression & Typical height\\
\midrule
$x^{\sqrt2}+x^{-1}$ & $0$\\
$\e^{-x^2}+x^3\e^x$ & $1$\\
$\e^{-\e^x}$ & $2$\\
$\exp\!\bigl(-\exp(\e^x+x)\bigr)$ & $3$\\
\bottomrule
\end{tabular}
\end{center}

The log-free field is the union
\[
  \T_{\mathrm{lf}}=\bigcup_{n\ge0}\T_n.
\]
Every individual element has finite height, although the field as a whole has no
maximum height.

\section{Why the exponent is required to be purely large}

Suppose an exponent $A$ has the canonical decomposition
\[
  A=L+c+U,
\]
with $L$ purely large, $c\in\R$, and $U$ small. Then
\[
  \e^A=\e^L\e^c\e^U
  =\e^L\e^c\sum_{n\ge0}\frac{U^n}{n!}.
\]
Only the purely large part creates a genuinely new monomial. The constant is a
coefficient, and the small part expands as a Hahn series. Requiring $L$ to be
purely large therefore gives a normal form and prevents redundant monomials.

\section{Height wins}

A genuinely new positive exponential level outruns everything below it.

\begin{proposition}[Height wins]
Let $L>0$ be purely large of exact height $n$, and let $T\ne0$ have height at
most $n$. Then, for every real $a$,
\[
  x^a\e^L\succdom T,
  \qquad
  x^a\e^{-L}\precdom T.
\]
\end{proposition}

\begin{proof}[Proof sketch]
Write the leading monomial of $T$ as $x^b\e^K$, where $K$ has height at most
$n-1$ if the monomial has exact height $n$. The exact-height hypothesis implies
that the leading term of $L-K$ is the leading term of $L$, hence positive. Thus
$L>K$, and lexicographic ordering gives $x^a\e^L\succdom x^b\e^K$. Replacing
$L$ by $-L$ gives the second claim.
\end{proof}

This explains comparisons such as
\[
  \e^{-\e^x}\precdom \e^{-Nx}\precdom x^{-M}
\]
for every fixed $M,N>0$.

\section{Hereditary structure}

A monomial $\m=x^a\e^L$ contains the transseries $L$ inside its exponent, whose
monomials may themselves contain exponents, and so on. A ratio set used for
computation should therefore be \emph{hereditary}: it should include enough
ratios to represent all exponents occurring inside its generators.

This recursive tree is finite in depth because every standard log-free
transseries has finite height. It is one reason the hierarchy is suitable for
algorithms.

\begin{figure}[ht]
\centering
\begin{tikzpicture}[
  level distance=12mm,
  level 1/.style={sibling distance=45mm},
  level 2/.style={sibling distance=25mm},
  every node/.style={draw,rounded corners,fill=softblue,inner sep=4pt},
  edge from parent/.style={draw,-Latex}]
\node {$\exp(L)$}
 child {node {$L=\e^x+x^2$}
   child {node {$\e^x$}
     child {node {$x$}}}
   child {node {$x^2$}}};
\end{tikzpicture}
\caption{A tiny hereditary exponent tree. The construction terminates because height is finite.}
\label{fig:height-tree}
\end{figure}

\section{What is still missing}

The hierarchy already contains towers of exponentials, but it does not contain
$\log x$. We could try to introduce logarithms as new primitive symbols at every
stage. A cleaner route is to construct all log-free objects first and then
compose them with iterated logarithms.

\begin{summarybox}
Starting from real powers of $x$, each new height adjoins exponentials of purely
large transseries from the previous stage. Lexicographic ordering makes higher
exponential levels dominate lower ones. Every standard transseries has finite
height.
\end{summarybox}

\chapter{Adding logarithms and defining the full transline}
\label{ch:logs}

\section{Iterated logarithms as changes of scale}

Write
\[
  \log_0 x=x,\qquad
  \log_1x=\log x,\qquad
  \log_2x=\log\log x,
\]
and in general
\[
  \log_{m+1}x=\log(\log_mx).
\]
For sufficiently large real $x$, every fixed iterate is defined and tends to
$+\infty$.

If $T$ is log-free, substitute $\log_mx$ for $x$ and write
\[
  T\compose\log_m.
\]
For instance,
\[
  (x^a\e^{-x})\compose\log
  =(\log x)^a x^{-1},
\]
and
\[
  \e^{-\sqrt{x}}\compose\log_2
  =\exp\!\left(-\sqrt{\log\log x}\right).
\]

\begin{definition}[Logarithmic depth]
A transseries has depth at most $m$ if it can be written as
$T\compose\log_m$ with $T$ log-free. Its exact depth is the least such $m$ after
identifying expressions that simplify by $\exp(\log x)=x$.
\end{definition}

The standard logarithmic-exponential transseries field is
\begin{equation}
  \T=\bigcup_{m\ge0}
       \{T\compose\log_m:T\in\T_{\mathrm{lf}}\}.
  \label{eq:full-transseries-field}
\end{equation}
It is often called the \emph{transline} \cite{vanDerHoevenBook,EdgarBeginners}.

\section{Examples of finite depth}

All of the following are standard transseries:
\[
  \log x+\frac1{\log x}+\frac1{(\log x)^2}+\cdots,
\]
\[
  \exp\!\left(\frac{x}{\log x}\right),
  \qquad
  x^x=\exp(x\log x),
\]
\[
  \exp\!\left(-\exp\sqrt{\log x}\right),
  \qquad
  \frac{\log_3x}{x(\log x)^2}.
\]

The expression $x^x$ is not introduced as a primitive monomial. It is normalized
as $\exp(x\log x)$ and lies at finite height and depth.

\section{Finite depth is an elementwise condition}

The field contains $\log_mx$ for every fixed $m$, but the formal sum
\[
  x+\log x+\log_2x+\log_3x+\cdots
\]
does not belong to the standard finite-depth field because one element would
require unbounded logarithmic depth.

Similarly, a series whose support contains exponential towers of unbounded
height is not a standard LE transseries. Larger structures---nested transseries,
logarithmic hyperseries, and fields with transexponential functions---have been
constructed to accommodate such behavior; see
\cite{SchmelingThesis,vanDenDriesHoevenKaplan,PadgettTransexponential}.

\begin{warningbox}
``The field contains arbitrarily deep iterated logarithms'' does not imply ``a
single series may use unbounded depth.'' The first is a statement about a union
of stages; the second would require a larger completion.
\end{warningbox}

\section{LE and EL conventions}

The literature contains several closely related constructions.

\begin{description}
\item[LE-series]
One first builds logarithmic structure and then closes under exponentiation, or
uses an equivalent staged construction.

\item[EL-series]
One begins with exponential structure and then adds logarithmic changes of
scale.

\item[Grid-based versus well-based]
One either restricts supports to finite grids or allows arbitrary well-based
supports.
\end{description}

Under suitable identifications, the classical constructions have large common
subfields and often yield isomorphic versions of the same practical transseries
field; the comparison is subtle at the level of ambient Hahn fields and chosen
closures \cite{vanDenDriesMacintyreMarker1997,KuhlmannTressl}.

In this book, ``$\T$'' means a standard real LE transseries field with the usual
operations. Whenever a statement depends on a particular variant, we say so.

\section{What is not in the basic real transline}

\subsection{Persistent oscillation}

Functions such as $\sin x$ and $\cos x$ do not have a total eventual dominance
order: neither is eventually larger in magnitude. A real ordered field designed
for nonoscillatory growth therefore does not contain them in the naive way.
Complex or oscillatory transseries enlarge the monomial language, for example by
using $\e^{\mathrm{i}x}$, but then order must be treated differently.

\subsection{Arbitrary infinite nesting}

An expression such as
\[
  \exp(-x+\exp(-x+\exp(-x+\cdots)))
\]
requires a separate fixed-point interpretation. It is not admitted merely by
writing infinitely many exponentials.

\subsection{Uncontrolled supports}

Even if every individual monomial is meaningful, their set may fail to be
well-based. The sum
\[
  \sum_{n\ge1}x^{-1/n}
\]
remains illegal in the basic Hahn construction.

\section{The algebraic closure properties to expect}

The standard field $\T$ is much more than an ordered field. It supports
\begin{itemize}
\item real powers of positive elements;
\item total exponential and logarithm on their natural ordered domains;
\item differentiation and antiderivatives;
\item composition $T\compose S$ when $S$ is positive infinite;
\item compositional inverses of positive infinite elements;
\item solutions to broad classes of algebraic and differential equations.
\end{itemize}
Precise constructions begin in the next chapter and Part~IV. Deep structural
forms of these statements are developed in
\cite{vanDerHoevenBook,ADHBook}.

\begin{summarybox}
The full transline is obtained by composing finite-height log-free transseries
with finitely iterated logarithms. Every element has finite height and depth.
The basic real field is deliberately nonoscillatory; larger transseries and
hyperseries fields relax these restrictions.
\end{summarybox}

\chapter{The anatomy and normal forms of a transseries}
\label{ch:anatomy}

\section{Additive decomposition}

For $T\in\T$, split its support according to comparison with $1$:
\begin{equation}
  T=T_{\mathrm{pl}}+c+T_{\mathrm{sm}}.
  \label{eq:canonical-additive-decomposition}
\end{equation}
Here
\begin{itemize}
\item $T_{\mathrm{pl}}$ contains all monomials $\succdom1$;
\item $c=[1]T\in\R$;
\item $T_{\mathrm{sm}}$ contains all monomials $\precdom1$.
\end{itemize}
The decomposition is unique because the three supports are disjoint.

\begin{example}
For
\[
  T=\e^x-3x^2+\log x+7+x^{-1}+2\e^{-x},
\]
we have
\[
  T_{\mathrm{pl}}=\e^x-3x^2+\log x,
  \qquad c=7,
  \qquad T_{\mathrm{sm}}=x^{-1}+2\e^{-x}.
\]
\end{example}

\section{Multiplicative decomposition}

Every nonzero $T$ can be factored as
\begin{equation}
  T=a\m(1+U),
  \qquad
  a=\lc(T)\in\R^\times,
  \quad \m=\magop(T),
  \quad U\precdom1.
  \label{eq:canonical-multiplicative-decomposition}
\end{equation}
This is the transseries analogue of scientific notation: coefficient, scale,
and a relative correction.

\begin{example}
\[
  T=3x^2\e^x-5x\e^x+7
  =3x^2\e^x
   \left(1-\frac{5}{3x}+\frac{7}{3x^2\e^x}\right).
\]
The expression in parentheses is $1+$small.
\end{example}

The additive decomposition defines exponentiation; the multiplicative
decomposition defines inverse, logarithm, and real powers.

\section{Truncation}

If $\n\in\M$ is a cutoff monomial, define the truncation
\[
  T\big|_{\succsimdom\n}
  =\sum_{\m\succsimdom\n}[\m]T\,\m.
\]
Because the support is well-based, this retains an initial segment in dominance
order. It may be finite or infinite, depending on the cutoff and support.

A transseries field is \emph{truncation closed} if every such initial segment of
an element is again in the field. Standard transseries fields are built to have
this property. It is essential in recursive proofs and algorithms.

\section{Asymptotic neighborhoods and formal topology}

A statement
\[
  S-T\precdom\n
\]
means that $S$ and $T$ agree through every monomial at least as large as $\n$.
This gives a non-Archimedean notion of closeness: two series are close when their
first disagreement occurs very far down the support.

A sequence $T_j$ converges coefficientwise in a fixed grid if, for every
monomial, its coefficient eventually stabilizes and all supports remain in a
common lower-bounded grid. For example,
\[
  T_N=\sum_{n=0}^{N}x^{-n}
  \longrightarrow
  \sum_{n=0}^{\infty}x^{-n}.
\]
Numerical values need not converge for a fixed $x$; this is formal convergence.

\section{The logarithmic derivative}

For nonzero $T$, define
\[
  T^\dagger=\frac{T'}{T}.
\]
This operation measures relative variation. For a monomial
\[
  \m=x^a\e^L,
\]
we have
\begin{equation}
  \m^\dagger=\frac{a}{x}+L'.
  \label{eq:monomial-log-derivative}
\end{equation}
The logarithmic derivative converts multiplicative structure into additive
structure:
\[
  (ST)^\dagger=S^\dagger+T^\dagger,
  \qquad
  (T^b)^\dagger=bT^\dagger.
\]
It will organize differentiation, integration, and comparison of rapidly
varying monomials.

\section{Dominant balance}

Many equations can be attacked before any coefficient calculation by comparing
leading monomials. Suppose
\[
  F(x,T,T',\ldots)=0.
\]
One lists possible dominant monomials for each term and asks which can cancel.
A term strictly larger than every other term cannot disappear, so at least two
dominant contributions must have the same magnitude.

\begin{example}
For
\[
  y'+y=\frac1x,
\]
if $y\asymp a x^p$, then $y'\asymp apx^{p-1}$ is usually smaller than $y$.
Balancing $y$ with $x^{-1}$ gives $p=-1$ and $a=1$. The derivative then
generates the next correction. This quickly predicts the factorial recurrence
studied later.
\end{example}

\begin{recipebox}[title=Normal-form workflow]
For any nonzero transseries $T$:
\begin{enumerate}
\item Find the leading monomial $\m$ and coefficient $a$.
\item Factor $T=a\m(1+U)$ with $U\precdom1$.
\item Use geometric, binomial, exponential, or logarithmic power series on $U$.
\item Re-expand and sort monomials by dominance.
\item Truncate only after all terms capable of reaching the target order have
      been included.
\end{enumerate}
\end{recipebox}

\begin{summarybox}
Every transseries has unique additive and multiplicative normal forms. They
separate large, constant, and small parts, and isolate a leading monomial times a
relative small correction. Almost every operation reduces to applying an
ordinary formal power series to that correction.
\end{summarybox}


\part{Algebra and calculus on transseries}

\chapter{Field operations and real powers}
\label{ch:field-ops}

\section{Addition and multiplication revisited}

Once all terms are normalized into monomials, addition and multiplication are
exactly the Hahn-series operations from Chapter~\ref{ch:hahn}. The only new work
is recursive normalization of expressions inside exponential and logarithmic
monomials.

\begin{example}[Multiplication by sector order]
Let
\[
  A=1+a_1x^{-1}+\alpha\e^{-x}+\bigO(x^{-2},x^{-1}\e^{-x},\e^{-2x})
\]
and
\[
  B=1+b_1x^{-1}+\beta\e^{-x}+\bigO(x^{-2},x^{-1}\e^{-x},\e^{-2x}).
\]
Then
\[
  AB=1+(a_1+b_1)x^{-1}+(\alpha+\beta)\e^{-x}
     +\bigO(x^{-2},x^{-1}\e^{-x},\e^{-2x}).
\]
The omitted product $\alpha\beta\e^{-2x}$ lies two exponential sectors down.
\end{example}

Here the notation $\bigO(\cdots)$ is only mnemonic for a set of discarded
monomials. A rigorous formal calculation instead specifies a cutoff support.

\section{Inverse by leading-term factorization}

Given
\[
  T=a\m(1+U),\qquad U\precdom1,
\]
we have
\[
  T^{-1}=a^{-1}\m^{-1}(1-U+U^2-U^3+\cdots).
\]

\begin{example}
Let
\[
  T=x+\log x+\e^{-x}.
\]
Then
\begin{align*}
 T^{-1}
 &=x^{-1}
   \left(1+x^{-1}\log x+x^{-1}\e^{-x}\right)^{-1}\\
 &=x^{-1}-x^{-2}\log x+x^{-3}(\log x)^2
   -x^{-2}\e^{-x}
   +2x^{-3}(\log x)\e^{-x}+\cdots.
\end{align*}
The ordering tells us that all pure inverse-power-log terms dominate the first
$\e^{-x}$ term, even though we may display both for structural reasons.
\end{example}

\section{Integer and rational powers}

Integer powers use multiplication and inverse. For a positive transseries $T$
and arbitrary real $b$, use the multiplicative normal form
\[
  T=a\m(1+U),\qquad a>0,\\ U\precdom1.
\]
Define
\begin{equation}
  T^b=a^b\m^b(1+U)^b,
  \qquad
  (1+U)^b=\sum_{n=0}^{\infty}\binom{b}{n}U^n.
  \label{eq:real-power-definition}
\end{equation}
The monomial power is $\m^b=\exp(b\log\m)$.

\begin{example}[Square root]
For
\[
  T=x^2+x+\e^{-x},
\]
we find
\begin{align*}
  \sqrt{T}
  &=x\left(1+x^{-1}+x^{-2}\e^{-x}\right)^{1/2}\\
  &=x+\frac12-\frac{1}{8x}+\frac{1}{16x^2}
    +\frac12x^{-1}\e^{-x}+\cdots.
\end{align*}
\end{example}

\section{Formal identities are genuine identities}

Because the binomial, geometric, exponential, and logarithmic families are
summable on small arguments, the familiar identities hold:
\[
  T^aT^b=T^{a+b},\qquad (T^a)^b=T^{ab},
\]
for positive $T$, and
\[
  (ST)^a=S^aT^a
\]
for positive $S,T$. These are not numerical convergence claims; they are
identities in the ordered field.

\section{Real closedness}

A field is real closed if every positive element has a square root and every
odd-degree polynomial has a root, equivalently if adjoining $\mathrm{i}$ makes
it algebraically closed. Standard fields of logarithmic-exponential transseries
are real closed \cite{vanDerHoevenBook,ADHBook}.

For calculations, the practical consequence is that algebraic equations can be
solved recursively by dominant balance and formal Newton expansion. The theorem
is much deeper than the elementary square-root calculation above, but the local
mechanism is similar: choose a leading monomial, factor it out, and solve for a
small correction.

\begin{summarybox}
Field operations reduce to Hahn-series convolution and geometric inversion.
Positive transseries admit arbitrary real powers through a binomial expansion of
their relative small part. Standard transseries fields are real closed.
\end{summarybox}

\chapter{Exponential and logarithm}
\label{ch:exp-log}

\section{Exponentiating a general transseries}

Take the additive decomposition
\[
  T=L+c+U,
\]
where $L$ is purely large, $c\in\R$, and $U$ is small. Define
\begin{equation}
  \exp T=\e^L\e^c
  \sum_{n=0}^{\infty}\frac{U^n}{n!}.
  \label{eq:exp-definition}
\end{equation}
Each factor has a different role:
\begin{itemize}
\item $\e^L$ is a new transmonomial;
\item $\e^c$ is an ordinary positive real coefficient;
\item $\e^U$ is a legal Hahn series because $U\precdom1$.
\end{itemize}

\begin{example}
For
\[
 T=x+3+x^{-1}+2\e^{-x},
\]
\begin{align*}
 \e^T
 &=\e^3\e^x
   \exp(x^{-1}+2\e^{-x})\\
 &=\e^3\e^x\left(
   1+x^{-1}+\frac{1}{2x^2}+\cdots
   +2\e^{-x}+2x^{-1}\e^{-x}+\cdots
   \right).
\end{align*}
After multiplying by $\e^x$, the term $2\e^{-x}$ inside the parenthesis becomes
the constant contribution $2\e^3$.
\end{example}

\section{Logarithm of a positive transseries}

Let $T>0$ and use the multiplicative decomposition
\[
  T=a\m(1+U),\qquad a>0,\quad U\precdom1.
\]
Define
\begin{equation}
  \log T=\log a+\log\m+
  \sum_{n=1}^{\infty}\frac{(-1)^{n+1}}{n}U^n.
  \label{eq:log-definition}
\end{equation}
If $\m=x^b\e^L$, then
\[
  \log\m=b\log x+L.
\]
In the full transseries field this again lies in $\T$.

\begin{example}
\begin{align*}
 \log\left(x^3\e^{2x}
   \left(1+\frac{\log x}{x}+\e^{-x}\right)\right)
 &=2x+3\log x\\
 &\quad+\frac{\log x}{x}+\e^{-x}
 -\frac12\left(\frac{\log x}{x}+\e^{-x}\right)^2+\cdots.
\end{align*}
\end{example}

\section{Exponential and logarithm are inverses}

The formal power-series identities
\[
  \log(\e^U)=U,
  \qquad
  \e^{\log(1+U)}=1+U
\]
hold for small $U$. The large and constant factors in the normal forms are
handled exactly, so
\[
  \log(\exp T)=T
\]
for every $T\in\T$, and
\[
  \exp(\log S)=S
\]
for every positive $S\in\T$.

Thus $\exp:\T\to\T_{>0}$ is an order-preserving group isomorphism from
addition to multiplication.

\section{Some instructive expansions}

\subsection{A small exponential}
\[
  \exp(x^{-1})
  =1+x^{-1}+\frac{1}{2}x^{-2}+\frac{1}{6}x^{-3}+\cdots.
\]

\subsection{A logarithm with a large leading term}
\begin{align*}
  \log(x+\log x)
  &=\log x+\log\left(1+\frac{\log x}{x}\right)\\
  &=\log x+\frac{\log x}{x}
     -\frac{(\log x)^2}{2x^2}
     +\frac{(\log x)^3}{3x^3}-\cdots.
\end{align*}

\subsection{An exponential of mixed scales}
\begin{align*}
 \exp\left(\e^{-x}+x^{-1}\e^{-x}\right)
 &=1+(1+x^{-1})\e^{-x}
   +\frac12(1+2x^{-1}+x^{-2})\e^{-2x}+\cdots.
\end{align*}

\section{Why positivity is necessary for the real logarithm}

The ordered real field has no real logarithm of a negative element. A complex
transseries field can define branches of $\log T$ after choosing arguments, but
branch information then becomes part of the structure. In this tutorial,
$\log T$ always assumes $T>0$ unless explicitly stated otherwise.

\section{Exponentiation changes height by at most one}

If $T$ has height $n$, then $\e^T$ has height at most $n+1$. It has no greater
height when the purely large part of $T$ is absent, because then $\e^T$ is only
a power series in a small transseries times a real constant.

This illustrates a recurring principle: the large part controls genuinely new
asymptotic scales; the small part controls formal corrections on existing scales.

\begin{summarybox}
To exponentiate, split a transseries additively into purely large, constant, and
small parts. To take a logarithm, factor a positive transseries into coefficient,
leading monomial, and $1+$small. The ordinary power-series formulas then become
exact Hahn-field identities.
\end{summarybox}

\chapter{Differentiation}
\label{ch:differentiation}

\section{Recursive definition}

Start with
\[
  c'=0,\qquad x'=1,\qquad (x^a)'=ax^{a-1}.
\]
For a transmonomial $\m=x^a\e^L$, define
\begin{equation}
  \m'=\left(\frac{a}{x}+L'\right)\m.
  \label{eq:monomial-derivative}
\end{equation}
Then extend termwise:
\begin{equation}
  \left(\sum_{\m}a_{\m}\m\right)'
  =\sum_{\m}a_{\m}\m'.
  \label{eq:termwise-derivative}
\end{equation}
A support theorem, proved recursively on height, guarantees that the family on
the right is summable \cite{EdgarBeginners,vanDerHoevenBook}.

\section{Examples}

\begin{align*}
  (x^{\sqrt2}\e^{-x^2})'
  &=\left(\frac{\sqrt2}{x}-2x\right)x^{\sqrt2}\e^{-x^2},\\
  (\e^{-\e^x})'
  &=-\e^x\e^{-\e^x},\\
  (\log_3x)'
  &=\frac1{x\log x\log_2x},\\
  (x^x)'
  &=x^x(\log x+1).
\end{align*}

The last formula follows from $x^x=\exp(x\log x)$, not from treating $x$ as a
constant exponent.

\section{The usual differential rules}

The recursive definition satisfies
\[
  (S+T)'=S'+T',
\]
\[
  (ST)'=S'T+ST',
\]
\[
  (T^a)'=aT^{a-1}T'\quad(T>0),
\]
\[
  (\e^T)'=T'\e^T,
  \qquad
  (\log T)'=\frac{T'}{T}\quad(T>0).
\]
These identities can be checked first on monomials and small power-series
substitutions, then extended by summability.

\section{Logarithmic derivatives}

The formula
\[
  T^\dagger=\frac{T'}{T}
\]
is especially efficient for monomials and products. For example,
\[
  \m=x^a(\log x)^b
      \exp\!\left(cx+\frac{x}{\log x}\right)
\]
has
\begin{align*}
  \m^\dagger
  &=\frac{a}{x}+\frac{b}{x\log x}+c
    +\left(\frac{x}{\log x}\right)'\\
  &=c+\frac1{\log x}-\frac1{(\log x)^2}
    +\frac{a}{x}+\frac{b}{x\log x}.
\end{align*}
Therefore $\m'=\m\m^\dagger$ without expanding the large product.

\section{Differentiating an asymptotic series}

Formal differentiation is always defined in $\T$, but analytic termwise
differentiation of an asymptotic expansion needs hypotheses. A function can be
small while its derivative oscillates wildly. The nonoscillatory Hardy-field
setting avoids many such pathologies, and transseries differentiation is
designed to model it.

\begin{warningbox}
Do not infer an analytic derivative expansion merely from a pointwise
asymptotic expansion unless the function class controls derivatives. In the
formal field, differentiation is algebraically valid; analytic realization is a
separate theorem.
\end{warningbox}

\section{Dominance and l'H\^opital behavior}

In a Hardy-type transseries field, if $S,T$ are nonconstant and sufficiently
well behaved, comparison of $S$ and $T$ is reflected by comparison of their
derivatives. This is the algebraic shadow of l'H\^opital's rule and is encoded in
the theory of $H$-fields.

For hand calculations, a safer local rule is:
\begin{enumerate}
\item differentiate each relevant monomial using its logarithmic derivative;
\item compare the resulting products;
\item only then select the leading derivative term.
\end{enumerate}
The derivative of the leading term usually dominates, but constants disappear
and special cancellations can occur.

\section{Example: the factorial recurrence appears}

Assume
\[
  y=\sum_{n=0}^{\infty}a_nx^{-n-1}.
\]
Then
\[
  y'=\sum_{n=0}^{\infty}-(n+1)a_nx^{-n-2}.
\]
Substitution into $y'+y=x^{-1}$ gives
\[
  a_0=1,
  \qquad
  a_n=na_{n-1}\quad(n\ge1),
\]
so $a_n=n!$. A perfectly routine formal derivative thus exposes the divergence
mechanism.

\begin{summarybox}
Differentiation is defined recursively on monomials and extended termwise.
Logarithmic derivatives turn nested products and exponentials into manageable
sums. Formal differentiation is unconditional inside $\T$; termwise
differentiation of actual asymptotic functions requires analytic control.
\end{summarybox}

\chapter{Integration and asymptotic antidifferentiation}
\label{ch:integration}

\section{Elementary antiderivatives}

The familiar rules remain valid:
\[
  \int x^a\dd x=\frac{x^{a+1}}{a+1}\quad(a\ne-1),
  \qquad
  \int\frac{\dd x}{x}=\log x,
\]
\[
  \int\frac{\dd x}{x\log x}=\log_2x,
  \qquad
  \int\frac{\dd x}{x\log x\cdots\log_kx}=\log_{k+1}x.
\]
The presence of all finite iterated logarithms removes the elementary
``exceptional exponents'' that would obstruct integration in a smaller power
field.

\section{Why term-by-term monomial integration is not enough}

The antiderivative of $\e^{x^2}$ is not an elementary monomial. Nevertheless it
has a transseries expansion. Let
\[
  \m=\e^L,
  \qquad q=\m^\dagger=L'.
\]
Seek an antiderivative in the form $F=\m h$. Then
\[
  F'=\m(h'+qh).
\]
Thus $F'=\m$ is equivalent to
\begin{equation}
  h'+qh=1.
  \label{eq:integration-h-equation}
\end{equation}
If differentiation of $h$ is smaller than multiplication by $q$, solve
recursively:
\begin{equation}
  h_0=q^{-1},
  \qquad
  h_{n+1}=-q^{-1}h_n'.
  \label{eq:integration-recursion}
\end{equation}
Then
\[
  h\asympseries h_0+h_1+h_2+\cdots.
\]
Indeed,
\[
  qh_0=1,
  \qquad
  h_n'+qh_{n+1}=0,
\]
so successive residuals cancel.

\section{Example: \texorpdfstring{$\int\e^{x^2}\dd x$}{the integral of exp(x squared)}}

Here $q=2x$. The recursion gives
\[
  h_0=\frac1{2x},
  \qquad
  h_1=\frac1{4x^3},
  \qquad
  h_2=\frac3{8x^5},
  \qquad
  h_3=\frac{15}{16x^7},\ldots
\]
Therefore
\begin{equation}
  \int\e^{x^2}\dd x
  \asympseries
  \e^{x^2}\left(
   \frac1{2x}+\frac1{4x^3}+\frac3{8x^5}
   +\frac{15}{16x^7}+\cdots
  \right)+C.
  \label{eq:exp-x2-antiderivative}
\end{equation}
The general coefficient is
\[
  \frac{(2n-1)!!}{2^{n+1}x^{2n+1}}.
\]

\begin{proof}
Let
\[
  F_N=\e^{x^2}\sum_{n=0}^{N-1}
  \frac{(2n-1)!!}{2^{n+1}x^{2n+1}},
  \qquad (-1)!!=1.
\]
Differentiating shows telescopic cancellation of all retained inverse powers,
leaving
\[
  F_N'=\e^{x^2}\left(1-
  \frac{(2N-1)!!}{2^N x^{2N}}
  \right).
\]
Hence each added term improves the relative error by $x^{-2}$.
\end{proof}

\section{Example: a decaying Gaussian tail}

For $\m=\e^{-x^2}$, $q=-2x$. A convenient antiderivative tending to zero is
\[
  \int_{\infty}^{x}\e^{-t^2}\dd t
  \asympseries
  -\frac{\e^{-x^2}}{2x}
  \left(1-\frac1{2x^2}+\frac3{4x^4}
  -\frac{15}{8x^6}+\cdots\right).
\]
The alternating signs arise from the negative logarithmic derivative.

\section{When the first recursion is not the right one}

For $\m=x^{-1}$, the logarithmic derivative is $q=-x^{-1}$, and the naive
choice $q^{-1}=-x$ does not yield a decreasing correction scheme. This signals a
change of asymptotic scale rather than failure of integration: the exact
antiderivative is $\log x$.

General integration algorithms classify the leading monomial and choose a
suitable asymptotic integrator. The full theorem is that the standard field
$\T$ is closed under antiderivatives. More strongly, it is Liouville closed:
every element has an antiderivative, and every nonzero element has a logarithmic
antiderivative in the appropriate sense \cite{vanDerHoevenBook,ADHBook}.

\section{Uniqueness after fixing the constant term}

If $F'=T$ and $G'=T$, then $(F-G)'=0$. The constant field of $\T$ is $\R$, so
$F-G\in\R$. Requiring $[1]F=0$ selects a unique formal antiderivative.

\begin{recipebox}[title=Asymptotic integration of a rapidly varying monomial]
For $\m$ with logarithmic derivative $q=\m^\dagger$:
\begin{enumerate}
\item Try $F=\m h$.
\item Solve $h'+qh=1$ by $h_0=q^{-1}$ and
      $h_{n+1}=-q^{-1}h_n'$.
\item Check that $h_{n+1}\precdom h_n$.
\item If the sequence does not decrease, change scale; logarithms are the most
      common missing scale.
\end{enumerate}
\end{recipebox}

\begin{summarybox}
Integration is not merely termwise division of monomials. For rapidly varying
terms, factor out the monomial and solve a first-order correction equation. For
critical power-log terms, new logarithmic scales appear. The full transseries
field is closed under antiderivatives.
\end{summarybox}


\chapter{Composition, Taylor expansion, and compositional inverse}
\label{ch:composition}

\section{Why the inner argument must be large}

For ordinary functions, $T\compose S$ means ``replace $x$ in $T$ by $S$.'' The
same idea works formally when $S$ is positive and infinitely large:
\[
  S>0,\qquad S\succdom1.
\]
This condition preserves the asymptotic direction $x\to+\infty$.

If
\[
  T=\sum_{n\ge0}x^{-n}
\]
and we substitute the small series $S=x^{-1}$, we obtain
\[
  \sum_{n\ge0}x^n,
\]
which is not well-based at infinity. Thus arbitrary right composition cannot be
allowed.

\section{Recursive definition}

For a monomial $\m=x^a\e^L$, define
\begin{equation}
  \m\compose S=S^a\exp(L\compose S).
  \label{eq:monomial-composition}
\end{equation}
Then for
\[
  T=\sum_{\m}c_{\m}\m,
\]
define
\begin{equation}
  T\compose S=\sum_{\m}c_{\m}(\m\compose S).
  \label{eq:series-composition}
\end{equation}
The difficult part is proving that the family on the right remains summable. The
proof proceeds by induction on the height and depth of $T$; it is one of the
technically subtler foundations of the subject
\cite{EdgarComposition,vanDerHoevenBook,BerarducciMantovaGerms}.

\section{Elementary examples}

\begin{align*}
  (x^3+\log x)\compose\e^x
  &=\e^{3x}+x,\\
  \e^{-x}\compose(x+\log x)
  &=\e^{-x}x^{-1},\\
  \log\compose(x+\e^{-x})
  &=\log x+\log(1+x^{-1}\e^{-x})\\
  &=\log x+x^{-1}\e^{-x}
    -\frac12x^{-2}\e^{-2x}+\cdots.
\end{align*}

Composition respects the field operations:
\[
  (A+B)\compose S=A\compose S+B\compose S,
\]
\[
  (AB)\compose S=(A\compose S)(B\compose S),
\]
\[
  \exp(A)\compose S=\exp(A\compose S),
  \qquad
  \log(A)\compose S=\log(A\compose S)
\]
when the logarithms are defined.

\section{The chain rule}

The formal derivative satisfies
\begin{equation}
  (T\compose S)'=(T'\compose S)S'.
  \label{eq:chain-rule}
\end{equation}
It is proved first for monomials using
\eqref{eq:monomial-composition}, then extended to summable families.

\section{Taylor expansion}

For a familiar analytic function,
\[
  T(S+U)=\sum_{n\ge0}\frac{T^{(n)}(S)}{n!}U^n.
\]
In transseries the same formula holds when $U$ is sufficiently small relative to
$S$ and to the internal scales of $T$:
\begin{equation}
  T\compose(S+U)
  =\sum_{n=0}^{\infty}
    \frac{T^{(n)}\compose S}{n!}U^n.
  \label{eq:transseries-taylor}
\end{equation}
The phrase ``sufficiently small'' is important. For a fixed grid-based $T$ and
$S$, one can formulate it using a finite witness set of monomial inequalities
\cite{EdgarComposition}.

For elementary outer functions, the condition is transparent.

\begin{example}[Logarithm]
If $U\precdom S$ and $S>0$, then
\begin{align*}
  \log(S+U)
  &=\log S+\log(1+U/S)\\
  &=\log S+\sum_{n\ge1}\frac{(-1)^{n+1}}{n}
      \left(\frac{U}{S}\right)^n.
\end{align*}
This is exactly the Taylor series because
$\log^{(n)}S=(-1)^{n-1}(n-1)!S^{-n}$.
\end{example}

\begin{example}[Exponential]
For any $U$ small enough that $(U^n)$ is summable,
\[
  \exp(S+U)=\exp S\sum_{n\ge0}\frac{U^n}{n!}.
\]
\end{example}

\section{Positive infinite transseries form a composition group}

Let
\[
  \mathcal P=\{T\in\T:T>0,\ T\succdom1\}.
\]
Then $x$ is the identity element under composition, and every $T\in\mathcal P$
has a unique compositional inverse $T^{[-1]}\in\mathcal P$ satisfying
\[
  T\compose T^{[-1]}=x,
  \qquad
  T^{[-1]}\compose T=x.
\]
We use square brackets to distinguish compositional powers from ordinary powers:
\[
  T^{[2]}=T\compose T,
  \qquad T^{[-1]}=\text{compositional inverse}.
\]

\section{How to calculate an inverse}

Suppose
\[
  T=x+A,
  \qquad A\precdom x.
\]
Seek $S=x+B$ with $B\precdom x$. The equation $T\compose S=x$ becomes
\begin{equation}
  B+A\compose(x+B)=0.
  \label{eq:inverse-fixed-point}
\end{equation}
This is a fixed-point equation
\[
  B=-A\compose(x+B).
\]
Repeated substitution usually determines successively smaller terms.

Alternatively, use Newton iteration for composition:
\begin{equation}
  S_{n+1}=S_n-\frac{T\compose S_n-x}{T'\compose S_n}.
  \label{eq:composition-newton}
\end{equation}
When the initial leading approximation is correct, the error improves rapidly in
the transseries valuation.

\begin{example}[Inverse of $x+\log x$]
Let $T=x+\log x$. Write
\[
  S=x-\log x+R,
  \qquad R\precdom\log x.
\]
Substitution gives
\begin{align*}
 T\compose S
 &=x-\log x+R+\log(x-\log x+R)\\
 &=x+R-\frac{\log x-R}{x}
   -\frac{(\log x-R)^2}{2x^2}-\cdots.
\end{align*}
Solving recursively yields
\[
  T^{[-1]}
  =x-\log x+\frac{\log x}{x}
   +\frac{(\log x)^2-2\log x}{2x^2}+\cdots.
\]
This is the same large-$x$ pattern that appears in the Lambert $W$ function.
\end{example}

\section{Fractional iteration}

Composition suggests the Abel equation
\[
  V\compose T=V+1.
\]
If a suitable $V$ is found, define a fractional iterate by
\[
  T^{[s]}=V^{[-1]}\compose(V+s),
  \qquad s\in\R.
\]
Then
\[
  T^{[s+t]}=T^{[s]}\compose T^{[t]}.
\]
The existence and uniqueness of such constructions depend on the leading form
of $T$ and on normalization choices. Transseries provide a natural language for
many cases, including shallow, moderate, and deep iteration regimes studied by
Edgar \cite{EdgarFractional}.

\begin{summarybox}
Right composition is defined for positive infinite inner arguments and requires
a nontrivial support theorem. It obeys the chain rule and suitable Taylor
expansions. Positive infinite transseries form a group under composition, so
formal functional inversion and fractional iteration become algebraic problems.
\end{summarybox}

\chapter{Contractive maps and recursive equations}
\label{ch:fixed-points}

\section{The non-Archimedean contraction principle}

In ordinary analysis, a contraction makes numerical distances smaller. In a
transseries field, a map is contractive when it pushes the leading error to a
strictly smaller monomial.

\begin{definition}[Formal contraction]
On a set of transseries, a map $\mathcal F$ is contractive if
\[
  Y\ne Z
  \quad\Longrightarrow\quad
  \mathcal F(Y)-\mathcal F(Z)\precdom Y-Z.
\]
\end{definition}

Under an appropriate completeness condition on a fixed support class, a
contractive map has a unique fixed point.

\begin{theorem}[Grid-based fixed-point principle]
Let $\mathcal A$ be a closed affine set of grid-based transseries supported in a
common grid, and suppose $\mathcal F:\mathcal A\to\mathcal A$ is uniformly
contractive in the sense that each application lowers errors by monomials from a
fixed finite ratio set. Then the iteration
\[
  Y_{n+1}=\mathcal F(Y_n)
\]
converges coefficientwise to the unique fixed point $Y=\mathcal F(Y)$.
\end{theorem}

\begin{proof}[Idea]
Let $E_n=Y_{n+1}-Y_n$. Contractivity gives
\[
  E_{n+1}\precdom E_n.
\]
The uniform grid condition implies that, for each fixed monomial, only finitely
many $E_n$ can contribute to its coefficient. Hence
\[
  Y_0+\sum_{n\ge0}E_n
\]
is a summable family and defines a limit $Y$. Continuity of $\mathcal F$ in this
formal topology gives $Y=\mathcal F(Y)$. If $Y$ and $Z$ were two fixed points,
then
\[
  Y-Z=\mathcal F(Y)-\mathcal F(Z)\precdom Y-Z,
\]
a contradiction.
\end{proof}

Costin gives a broad topological version of this principle for abstract
multiseries \cite{CostinTopological}; Edgar develops a grid-based formulation
for explicit computations \cite{EdgarBeginners}.

\section{Example: \texorpdfstring{$Y=x+\e^{-Y}$}{Y = x + exp(-Y)}}

Let
\[
  \mathcal F(Y)=x+\e^{-Y}.
\]
Start with $Y_0=x$. Then
\[
  Y_1=x+\e^{-x},
\]
\[
  Y_2=x+\e^{-x-\e^{-x}}
     =x+\e^{-x}-\e^{-2x}+\frac12\e^{-3x}+\cdots.
\]
Each iteration fixes more exponential coefficients. For two candidates
$Y,Z=x+\bigO(\e^{-x})$,
\[
  \e^{-Y}-\e^{-Z}
  =-\e^{-Z}(Y-Z)+\text{smaller terms},
\]
so the error is multiplied by $\e^{-x}$ and becomes strictly smaller. The map is
contractive on the appropriate affine grid.

\section{Example: recursive logarithmic correction}

Consider
\[
  Y=x-\log Y.
\]
The first approximation is $Y_0=x$, then
\[
  Y_1=x-\log x,
\]
\[
  Y_2=x-\log(x-\log x)
  =x-\log x+\frac{\log x}{x}
   +\frac{(\log x)^2}{2x^2}+\cdots.
\]
The derivative of the right-hand side with respect to $Y$ is $-1/Y\precdom1$,
so corrections contract by roughly $x^{-1}$.

\section{Formal implicit-function method}

Suppose
\[
  F(x,Y)=0
\]
and a leading approximation $Y_0$ has been found. Write $Y=Y_0+U$ and expand
\[
  0=F(x,Y_0)+F_Y(x,Y_0)U+\frac12F_{YY}(x,Y_0)U^2+\cdots.
\]
If $F_Y(x,Y_0)$ is invertible and the nonlinear remainder is smaller, solve
\[
  U=-F_Y(x,Y_0)^{-1}F(x,Y_0)+\text{smaller corrections}.
\]
This is Newton's method in the transseries valuation.

\begin{recipebox}[title=Solving an implicit equation]
\begin{enumerate}
\item Use dominant balance to guess the leading monomial of $Y$.
\item Substitute $Y=Y_0(1+U)$ or $Y=Y_0+U$ with $U$ relatively small.
\item Expand all analytic operations in $U$.
\item Solve the largest uncancelled coefficient equation.
\item Repeat; well-basedness ensures that ``next largest'' is meaningful.
\end{enumerate}
\end{recipebox}

\section{Differential equations and free exponential parameters}

For an ODE
\[
  P(x,Y,Y',\ldots,Y^{(r)})=0,
\]
one first computes a base formal sector $\Phi_0$. Linearizing around it gives a
linear differential operator $\mathcal L$. Exponential monomials in the kernel
of $\mathcal L$ indicate missing free parameters.

For a first-order equation, if the linearized perturbation $\eta$ satisfies
\[
  \eta'=Q(x)\eta,
\]
then
\[
  \eta=C\exp\left(\int Q(x)\dd x\right).
\]
This monomial becomes the ratio between transseries sectors. Nonlinear terms
then force powers of it and couple the sector coefficients recursively.

\section{Dominant balance and Newton polygons}

When several monomials involving $Y$ and its derivatives compete, one can plot
orders and find edges whose terms can balance. This is the differential analogue
of a Newton polygon for algebraic or Puiseux series. In more advanced
transseries algorithms, dominant balances determine candidate leading
monomials, while contractive maps construct the corrections.

\begin{summarybox}
A formal contraction lowers the leading error rather than a numerical norm.
Uniform contraction in a common grid gives existence and uniqueness by
coefficientwise stabilization. This is the basic engine behind implicit
solutions, compositional inverses, and many differential-equation transseries.
\end{summarybox}

\chapter{How to compute without drowning in terms}
\label{ch:computation}

\section{Always choose the scale first}

Before expanding, decide which monomials matter. A statement such as
\[
  \text{``through }x^{-4}\text{ in sectors }k=0,1,
  \text{ and through the leading term in }k=2\text{''}
\]
is much safer than ``expand to order 4.'' There is no universal total degree
that respects exponential dominance.

\section{A finite ratio-set representation}

A practical grid-based object may be stored as
\[
  T=\m_0\sum_{\bm k\in E}c_{\bm k}\boldsymbol\mu^{\bm k},
\]
where
\begin{itemize}
\item $\m_0$ is a base monomial;
\item $\boldsymbol\mu$ is a finite tuple of small generators;
\item $E\subseteq\N^r$ is represented lazily or finitely up to a cutoff;
\item coefficients may themselves be exact rational, algebraic, or symbolic
      expressions.
\end{itemize}

Recursive exponent trees encode monomials such as
$\exp(-\e^x+x^2)$, and a comparison routine recursively compares the purely
large exponents before power factors.

\section{Truncated arithmetic}

Let $\tau$ be the smallest retained monomial. Define a projection
\[
  \pi_{\tau}(T)=T\big|_{\succsimdom\tau}.
\]
For addition,
\[
  \pi_\tau(S+T)=\pi_\tau(\pi_\tau S+\pi_\tau T).
\]
For multiplication, one must retain every pair of input terms whose product can
reach $\tau$. Since supports are well-based, only finitely many such pairs occur
in a fixed grid.

For inverse, exponential, or logarithm, determine how many powers of the small
remainder $U$ can contribute above $\tau$. If $\magop(U)=\mu\precdom1$, stop once
\[
  \mu^n\precdom \frac{\tau}{\text{leading factor}}.
\]

\section{A generic algorithm}

\begin{tcolorbox}[enhanced,breakable,colback=softgray,colframe=black!65,
 title=Truncated unary operation $F(T)$]
\begin{enumerate}
\item Normalize $T$ into the additive or multiplicative form required by $F$.
\item Separate the exact large/monomial factor from a small remainder $U$.
\item Determine a finite $N$ such that powers $U^n$ for $n>N$ lie below the
      target monomial.
\item Evaluate the ordinary formal series of $F$ through degree $N$.
\item Multiply by the exact factor, combine equal monomials, sort by dominance,
      and discard terms below the target.
\item Verify the result by substitution or differentiation to one order beyond
      the requested cutoff whenever possible.
\end{enumerate}
\end{tcolorbox}

\section{Verification by residual}

The most reliable check is to substitute the truncated answer back into the
original equation. If $Y_N$ is claimed correct through $\tau$, compute the
residual
\[
  R_N=F(x,Y_N,Y_N',\ldots).
\]
The leading monomial of $R_N$ should be strictly below the accuracy implied by
$\tau$. This catches sign errors, missing cross-sector products, and incorrect
logarithmic powers.

\section{Common bookkeeping errors}

\begin{description}
\item[Discarding too early]
A term small inside an exponent may become important after multiplication by a
large outside monomial.

\item[Using ordinary degree]
$x^{-100}$ is still much larger than $\e^{-x}$. A degree that counts both as one
step destroys the dominance order.

\item[Forgetting generator closure]
Differentiation introduces logarithmic derivatives; composition introduces
substituted generators; inversion may require a new small ratio. The ratio set
must be enlarged.

\item[Ignoring resonance]
Two distinct multi-indices may yield the same monomial. Their coefficients must
be combined, and logarithmic sectors may be forced in differential equations.

\item[Confusing formal and numerical convergence]
A formally summable family need not converge when a finite real $x$ is
substituted. Conversely, an analytic function may have a divergent but useful
formal transseries.
\end{description}

\section{Computer algebra}

Van der Hoeven developed algorithms for generalized series and transseries that
perform comparison, arithmetic, composition, differential operations, and
asymptotic solving \cite{vanDerHoevenBook,vanDerHoevenOperators}. A persistent
challenge is zero testing: after recursive operations, deciding whether a
symbolic transseries expression is identically zero may be nontrivial. Recent
work gives a zero test for an important class of differentially algebraic
transseries \cite{ChenFangHoeven2026}.

\begin{summarybox}
Effective transseries computation is target-driven. Choose a monomial cutoff,
use finite ratio sets, expand only powers capable of reaching that cutoff, and
verify by residual. The actual dominance order, not a guessed total degree, must
guide truncation.
\end{summarybox}


\part{Worked transseries calculations}

\chapter{An implicit equation with an exponentially small correction}
\label{ch:implicit-exp}

\section{The problem}

Solve, for large positive $x$,
\begin{equation}
  y=x+\e^{-y}.
  \label{eq:implicit-exp-main}
\end{equation}
A power-series ansatz in $x^{-1}$ finds only $y\asympseries x$, because the
correction is beyond all powers. We therefore set
\[
  y=x+\delta,
  \qquad \delta\precdom1.
\]
Then
\begin{equation}
  \delta=\e^{-x}\e^{-\delta}.
  \label{eq:delta-equation}
\end{equation}
The natural small ratio is $z=\e^{-x}$.

\section{Direct coefficient recursion}

Assume
\[
  \delta=c_1z+c_2z^2+c_3z^3+\cdots.
\]
Expand the right side:
\[
  z\e^{-\delta}
  =z\left(1-\delta+\frac{\delta^2}{2}
  -\frac{\delta^3}{6}+\cdots\right).
\]
Matching powers of $z$ gives
\[
 c_1=1,
 \qquad c_2=-c_1=-1,
\]
\[
 c_3=-c_2+\frac12c_1^2=\frac32,
\]
\[
 c_4=-c_3+c_1c_2-\frac16c_1^3=-\frac83.
\]
Continuing,
\begin{equation}
  y=x+\e^{-x}-\e^{-2x}+\frac32\e^{-3x}
  -\frac83\e^{-4x}+\frac{125}{24}\e^{-5x}
  -\frac{54}{5}\e^{-6x}+\cdots.
  \label{eq:implicit-expansion}
\end{equation}

\section{Closed coefficient formula from Lambert \texorpdfstring{$W$}{W}}

Equation~\eqref{eq:delta-equation} is equivalent to
\[
  \delta\e^{\delta}=\e^{-x}.
\]
By definition of the Lambert $W$ function,
\[
  \delta=W(\e^{-x}).
\]
The Lagrange inversion formula gives
\begin{equation}
  W(z)=\sum_{n=1}^{\infty}
  \frac{(-n)^{n-1}}{n!}z^n,
  \qquad |z|<\e^{-1}.
  \label{eq:lambert-small-series}
\end{equation}
Substituting $z=\e^{-x}$ reproduces
\eqref{eq:implicit-expansion}. For $x>1$, this particular transseries actually
converges.

\begin{proof}[Derivation by Lagrange inversion]
Write $\delta=z\phi(\delta)$ with $\phi(u)=\e^{-u}$. Lagrange inversion says
\[
  [z^n]\delta(z)=\frac1n[u^{n-1}]\phi(u)^n
  =\frac1n[u^{n-1}]\e^{-nu}.
\]
Since
\[
  [u^{n-1}]\e^{-nu}=\frac{(-n)^{n-1}}{(n-1)!},
\]
the coefficient is $(-n)^{n-1}/n!$.
\end{proof}

\section{Why a contraction proves uniqueness}

On the affine set $x+\e^{-x}\R[[\e^{-x}]]$, define
\[
  \mathcal F(Y)=x+\e^{-Y}.
\]
If $Y-Z$ begins at $\e^{-kx}$, then
\[
  \mathcal F(Y)-\mathcal F(Z)
  =-\e^{-Z}(Y-Z)+\text{smaller terms}
\]
begins at $\e^{-(k+1)x}$. Thus the map pushes every error one full sector down.
The fixed point is unique in this support class.

\section{Residual check}

Let
\[
  y_N=x+\sum_{n=1}^{N}c_n\e^{-nx}.
\]
If the coefficients are correct through $N$, then
\[
  y_N-x-\e^{-y_N}=\bigO(\e^{-(N+1)x}).
\]
This check is easy to automate and is more reliable than inspecting individual
coefficient formulas.

\begin{bigidea}
A correction can be invisible to every inverse power and still possess a
perfectly ordinary convergent series in a different small monomial. The right
question is not ``does a power series exist?'' but ``which asymptotic ratio
generates the correction?''
\end{bigidea}

\chapter{Lambert \texorpdfstring{$W$}{W} at infinity: logarithms inside powers}
\label{ch:lambert-large}

\section{The equation}

Consider
\begin{equation}
  y+\log y=x,
  \qquad y>0.
  \label{eq:lambert-large-equation}
\end{equation}
Exponentiating gives
\[
  y\e^y=\e^x,
\]
so
\[
  y=W(\e^x).
\]
The dominant balance is $y\sim x$, but the logarithm contributes a correction of
size $\log x$.

\section{Successive approximation}

Start with $y_0=x$. Then
\[
  y=x-\log y
\]
suggests
\[
  y_1=x-\log x.
\]
Substitute again:
\begin{align*}
  y_2
  &=x-\log(x-\log x)\\
  &=x-\log x-
    \log\left(1-\frac{\log x}{x}\right)\\
  &=x-\log x+\frac{\log x}{x}
    +\frac{(\log x)^2}{2x^2}+\cdots.
\end{align*}
This is not yet the final coefficient at $x^{-2}$ because the $x^{-1}$
correction feeds back into the logarithm. Systematic recursion resolves it.

\section{Polynomial recursion}

Let $L=\log x$ and seek
\begin{equation}
  y=x-L+\sum_{n=1}^{\infty}\frac{P_n(L)}{x^n},
  \label{eq:W-large-ansatz}
\end{equation}
where $P_n$ is a polynomial. Put $t=x^{-1}$ and
\[
  U(t)=-Lt+\sum_{n\ge1}P_n(L)t^{n+1}.
\]
Since $y=x(1+U)$, equation~\eqref{eq:lambert-large-equation} becomes
\[
  \sum_{n\ge1}P_n(L)t^n+\log(1+U(t))=0.
\]
Therefore
\begin{equation}
  P_n(L)=-[t^n]\log\left(
    1-Lt+\sum_{j=1}^{n-1}P_j(L)t^{j+1}
  \right).
  \label{eq:Pn-recursion}
\end{equation}
The coefficient $P_n$ depends only on earlier polynomials.

The first six are
\begin{align*}
 P_1(L)&=L,\\
 P_2(L)&=\frac{L^2}{2}-L,\\
 P_3(L)&=\frac{L^3}{3}-\frac{3L^2}{2}+L,\\
 P_4(L)&=\frac{L^4}{4}-\frac{11L^3}{6}+3L^2-L,\\
 P_5(L)&=\frac{L^5}{5}-\frac{25L^4}{12}
          +\frac{35L^3}{6}-5L^2+L,\\
 P_6(L)&=\frac{L^6}{6}-\frac{137L^5}{60}
          +\frac{75L^4}{8}-\frac{85L^3}{6}
          +\frac{15L^2}{2}-L.
\end{align*}
Hence
\begin{align}
 W(\e^x)
 &=x-L+\frac{L}{x}
 +\frac{L(L-2)}{2x^2}
 +\frac{L(2L^2-9L+6)}{6x^3}
 \notag\\
 &\quad
 +\frac{L(3L^3-22L^2+36L-12)}{12x^4}
 +\bigO\left(\frac{L^5}{x^5}\right).
 \label{eq:W-large-expansion}
\end{align}

\section{Closed formula with Stirling numbers}

Let $\genfrac{[}{]}{0pt}{}{n}{k}$ denote the unsigned Stirling number of the first kind. The
classical large-argument expansion can be written
\begin{equation}
  W(z)=L_1-L_2+
  \sum_{k=0}^{\infty}\sum_{m=1}^{\infty}
   \frac{(-1)^k}{m!}
   \genfrac{[}{]}{0pt}{}{k+m}{k+1}
   \frac{L_2^m}{L_1^{k+m}},
  \label{eq:W-stirling-expansion}
\end{equation}
where $L_1=\log z$ and $L_2=\log\log z$. Setting $z=\e^x$ gives
$L_1=x$ and $L_2=\log x$.

Formula~\eqref{eq:W-stirling-expansion} is an excellent example of a log-power
transseries: the coefficient at each power $x^{-n}$ is a polynomial in
$\log x$.

\section{Compositional interpretation}

The map
\[
  T(y)=y+\log y
\]
is positive infinite and has a compositional inverse in $\T$. The expansion
above is precisely $T^{[-1]}(x)$. Thus the Lambert $W$ asymptotic is not an
isolated trick; it is an instance of the general composition-group theorem.

\section{A useful accuracy estimate}

Since $L/x\to0$, each successive displayed term is smaller than the preceding
one in the transseries order:
\[
  x\succdom L\succdom \frac{L}{x}
  \succdom\frac{L^2}{x^2}\succdom\cdots.
\]
For fixed truncation order $N$, substitution into
$y+\log y-x$ leaves a residual of order
\[
  \bigO\left(\frac{(\log x)^{N+1}}{x^{N+1}}\right).
\]

\begin{summarybox}
The large Lambert $W$ expansion illustrates how logarithms become polynomial
coefficients of inverse powers. A coefficient-extraction recursion determines
each polynomial, while Stirling numbers give a closed double sum. Formally, the
calculation is compositional inversion in $\T$.
\end{summarybox}

\chapter{The Euler differential equation and a hidden constant}
\label{ch:euler-ode}

\section{Formal solution}

Consider
\begin{equation}
  y'+y=\frac1x.
  \label{eq:euler-ode}
\end{equation}
Seek
\[
  \phihat_0(x)=\sum_{n=0}^{\infty}a_nx^{-n-1}.
\]
Substitution gives
\[
  a_0=1,
  \qquad a_n=na_{n-1},
\]
so
\begin{equation}
  \phihat_0(x)=
  \sum_{n=0}^{\infty}\frac{n!}{x^{n+1}}.
  \label{eq:euler-formal-series}
\end{equation}
The series diverges for every finite $x$ because the ratio of consecutive terms
is $(n+1)/x$.

\section{The missing homogeneous solution}

The homogeneous equation
\[
  y'+y=0
\]
has solution $C\e^{-x}$. Therefore the complete formal transseries solution is
\begin{equation}
  \widehat y(x;C)=
  \sum_{n=0}^{\infty}\frac{n!}{x^{n+1}}+C\e^{-x}.
  \label{eq:euler-complete-transseries}
\end{equation}
Every value of $C$ has the same inverse-power asymptotic expansion. The
transseries reveals the integration constant that the power sector hides.

\section{Exact analytic solution}

Multiplying \eqref{eq:euler-ode} by $\e^x$ gives
\[
  (\e^xy)'=\frac{\e^x}{x}.
\]
Hence, on a chosen branch,
\begin{equation}
  y(x)=\e^{-x}\bigl(\Ei(x)+C\bigr),
  \label{eq:euler-exact-solution}
\end{equation}
where $\Ei'(x)=\e^x/x$. Repeated integration by parts yields
\[
  \Ei(x)\asympseries
  \frac{\e^x}{x}
  \left(1+\frac1x+\frac{2!}{x^2}
  +\frac{3!}{x^3}+\cdots\right),
\]
which reproduces \eqref{eq:euler-formal-series}.

\section{Optimal truncation}

The $n$th term has size $n!/x^{n+1}$. It is smallest near $n\approx x$.
Stirling's formula gives the minimum scale
\[
  \frac{n!}{x^{n+1}}\bigg|_{n\approx x}
  \asymp \sqrt{\frac{2\pi}{x}}\,\e^{-x}.
\]
Thus the smallest achievable error from the power sector alone is on the same
exponential scale as the homogeneous solution.

\begin{bigidea}
The ``missing'' $C\e^{-x}$ term lives at precisely the scale where the divergent
power expansion stops improving. Optimal truncation points toward the next
transseries sector.
\end{bigidea}

\section{Borel preview}

Under the order-one Borel transform
\[
  \Borel\left(x^{-n-1}\right)=\frac{\xi^n}{n!},
\]
series~\eqref{eq:euler-formal-series} becomes
\begin{equation}
  \Borel\phihat_0(\xi)=\sum_{n\ge0}\xi^n=\frac1{1-\xi}.
  \label{eq:euler-borel}
\end{equation}
The factorial divergence has become a simple pole at $\xi=1$. Laplace
integration above or below this pole gives two lateral sums whose difference is
proportional to $\e^{-x}$. Part~VI develops this in detail.

\begin{summarybox}
The simplest linear ODE already displays the whole transseries story: a
factorially divergent perturbative sector, an exponentially small homogeneous
sector carrying a free constant, optimal truncation at the exponential scale,
and a Borel singularity encoding the same exponential.
\end{summarybox}

\chapter{Nonlinearity creates a ladder of sectors}
\label{ch:nonlinear-ode}

\section{The logistic equation}

Consider
\begin{equation}
  y'=y-y^2.
  \label{eq:logistic}
\end{equation}
The equilibrium $y=1$ has linearized perturbation $\eta'=-\eta$, so the natural
small monomial near $x=+\infty$ is $\e^{-x}$.

The exact one-parameter solution tending to $1$ is
\begin{equation}
  y(x)=\frac1{1+A\e^{-x}}.
  \label{eq:logistic-exact}
\end{equation}
Expanding geometrically,
\begin{equation}
  y=1-A\e^{-x}+A^2\e^{-2x}
    -A^3\e^{-3x}+\cdots.
  \label{eq:logistic-transseries}
\end{equation}

\section{Sector recursion without using the exact solution}

Write
\[
  y=1+\sum_{k\ge1}c_k\e^{-kx}.
\]
Then
\[
  y'=-\sum_{k\ge1}k c_k\e^{-kx},
\]
and
\[
  y-y^2=-\sum_{k\ge1}c_k\e^{-kx}
  -\left(\sum_{k\ge1}c_k\e^{-kx}\right)^2.
\]
At sector $k\ge2$,
\begin{equation}
  (k-1)c_k=\sum_{j=1}^{k-1}c_jc_{k-j}.
  \label{eq:logistic-sector-recurrence}
\end{equation}
The first coefficient $c_1$ is free. The recurrence gives
\[
  c_k=c_1^k.
\]
Taking $c_1=-A$ recovers \eqref{eq:logistic-transseries}.

This is the basic nonlinear mechanism: a first exponential perturbation is free,
but its products force all higher sectors.

\section{A more realistic Riccati example}

Now consider
\begin{equation}
  y'+y=\frac1x+y^2.
  \label{eq:riccati-transseries}
\end{equation}
The base power sector
\[
  \Phi_0(x)=\sum_{n\ge0}a_nx^{-n-1}
\]
satisfies
\begin{equation}
  a_0=1,
  \qquad
  a_n=na_{n-1}+\sum_{j=0}^{n-1}a_ja_{n-1-j}.
  \label{eq:riccati-a-recurrence}
\end{equation}
Thus
\[
  \Phi_0=\frac1x+\frac2{x^2}+\frac8{x^3}
  +\frac{44}{x^4}+\frac{296}{x^5}
  +\frac{2312}{x^6}+\cdots.
\]
The coefficients again grow rapidly.

Linearize by writing $y=\Phi_0+\eta$. To first order,
\[
  \eta'+\eta=2\Phi_0\eta,
\]
so
\[
  \eta'=(-1+2\Phi_0)\eta.
\]
Therefore
\begin{align*}
  \eta
  &=C\exp\left(-x+2\int\Phi_0\dd x\right)\\
  &=C\e^{-x}x^2
   \left(1-\frac4x-\frac8{x^3}
   -\frac{52}{x^4}+\cdots\right).
\end{align*}
The $x^2$ factor comes from integrating the leading $2/x$ term.

\section{Full sector equation}

Use the ansatz
\[
  y=\sum_{k=0}^{\infty}
  \sigma^k\e^{-kx}\Phi_k(x).
\]
Substitution into \eqref{eq:riccati-transseries} gives
\begin{equation}
  \Phi_k'+(1-k)\Phi_k
  =\sum_{j=0}^{k}\Phi_j\Phi_{k-j}.
  \label{eq:riccati-sector-equation}
\end{equation}
For $k\ge1$,
\begin{equation}
  \Phi_k'+(1-k-2\Phi_0)\Phi_k
  =\sum_{j=1}^{k-1}\Phi_j\Phi_{k-j}.
  \label{eq:riccati-sector-recursive}
\end{equation}
The $k=1$ equation is homogeneous and carries the free parameter. Every
$k\ge2$ equation is forced by lower sectors. Dominant balance suggests
\[
  \Phi_k(x)\asymp x^{2k}\times
  \text{an inverse-power series}.
\]

\section{Resonance warning}

In multi-parameter problems, different combinations of actions may produce the
same exponential monomial. Then the right side of a sector equation can resonate
with the homogeneous solution, forcing powers of $\log x$ or polynomial factors
in the transseries parameters. The simple one-action Riccati example avoids this
complication.

\begin{summarybox}
Linearization identifies the first exponential scale and its free coefficient.
Nonlinear products then generate a ladder of higher sectors. Sector equations
are typically linear in the unknown sector and forced by lower ones, making the
transseries recursively computable.
\end{summarybox}

\chapter{A singular perturbation and a boundary-layer exponential}
\label{ch:singular-perturbation}

\section{A small parameter changes the asymptotic variable}

Transseries need not use $x\to\infty$. Consider $\eps\to0^+$ in
\begin{equation}
  \eps\,\frac{\dd y}{\dd x}+y=\frac1{1+x},
  \qquad x>0.
  \label{eq:singular-perturbation}
\end{equation}
Formally invert the operator $1+\eps D$:
\[
  y_{\mathrm{out}}
  =(1+\eps D)^{-1}\frac1{1+x}
  =\sum_{n=0}^{\infty}(-\eps D)^n\frac1{1+x}.
\]
Since
\[
  D^n\frac1{1+x}=\frac{(-1)^n n!}{(1+x)^{n+1}},
\]
we obtain
\begin{equation}
  \widehat y_{\mathrm{out}}
  =\sum_{n=0}^{\infty}
    \frac{n!\,\eps^n}{(1+x)^{n+1}}.
  \label{eq:outer-divergent}
\end{equation}
For fixed $x$, this series diverges for every $\eps\ne0$.

\section{The homogeneous exponential}

The homogeneous equation
\[
  \eps y'+y=0
\]
has solution
\[
  C\e^{-x/\eps}.
\]
Thus the complete formal transseries is
\begin{equation}
  \widehat y(x,\eps;C)
  =\sum_{n=0}^{\infty}
    \frac{n!\,\eps^n}{(1+x)^{n+1}}
   +C\e^{-x/\eps}.
  \label{eq:singular-complete-transseries}
\end{equation}
The exponential is smaller than every power of $\eps$ for fixed $x>0$.

\section{Exact solution}

An integrating factor gives
\begin{equation}
  y(x)=\e^{-x/\eps}
  \left[
    C+\frac1\eps\int^x
      \frac{\e^{t/\eps}}{1+t}\dd t
  \right].
  \label{eq:singular-exact}
\end{equation}
Repeated integration by parts produces the outer series
\eqref{eq:outer-divergent}. The lower limit of the integral and the boundary
condition determine $C$.

\section{Boundary layers}

Suppose a boundary condition at $x=0$ disagrees with the value predicted by the
outer expansion. The discrepancy is repaired by $C\e^{-x/\eps}$, which changes
rapidly on the stretched scale
\[
  X=\frac{x}{\eps}.
\]
In the boundary layer, the exponential is order one; away from the boundary, it
is beyond all algebraic orders in $\eps$.

This is a concrete reason transseries appear in singular perturbation theory:
no finite power series in $\eps$ can satisfy both the differential equation and
a mismatched boundary condition.

\section{The same model in different clothes}

Set
\[
  X=\frac{1+x}{\eps}.
\]
After a simple rescaling, the factorial series in
\eqref{eq:outer-divergent} reduces to the Euler series of
Chapter~\ref{ch:euler-ode}. Its Borel singularity therefore lies at the action
corresponding to $\e^{-x/\eps}$.

\begin{summarybox}
Singular perturbations naturally produce power series in a small parameter plus
exponentials such as $\e^{-x/\eps}$. The power series may diverge factorially,
while the exponential carries boundary data invisible to all algebraic orders.
\end{summarybox}


% ============================================================================
\part{Divergence, Borel summation, and resurgence}
\label{part:resurgence}

\chapter{Factorial divergence is information, not failure}
\label{ch:factorial-divergence}

\section{A formal expansion need not converge}

Return to a formal inverse-power series
\begin{equation}
  \phihat(z)=\sum_{n=0}^{\infty}\frac{a_n}{z^{n+1}},
  \qquad z\to\infty.
  \label{eq:generic-inverse-power-series}
\end{equation}
The hat reminds us that the series is formal: at first we make no claim that the
sum converges. This is not an embarrassment. In asymptotic analysis, the useful
question is whether finite truncations approximate a function in a controlled
way.

\begin{definition}[Poincar\'e asymptotic expansion]
Let $S$ be an unbounded region in the complex plane. We write
\[
  f(z)\asympseries \sum_{n=0}^{\infty}\frac{a_n}{z^{n+1}}
  \qquad (z\to\infty\text{ in }S)
\]
if, for every fixed $N\ge0$,
\[
  f(z)-\sum_{n=0}^{N-1}\frac{a_n}{z^{n+1}}
  =\bigO\!\left(z^{-N-1}\right).
\]
Thus, with the normalization in \eqref{eq:generic-inverse-power-series}, the
remainder after terms $0,\ldots,N-1$ is $\bigO(z^{-N-1})$.
\end{definition}

The definition asks us to fix $N$ first and then let $z$ grow. It does
\emph{not} ask us to send $N$ to infinity for a fixed $z$. Those two limiting
processes need not agree.

\begin{example}[The Euler series again]
The equation
\[
  y'(z)+y(z)=\frac1z
\]
led in Chapter~\ref{ch:euler-ode} to
\begin{equation}
  \widehat E(z)=\sum_{n=0}^{\infty}\frac{n!}{z^{n+1}}.
  \label{eq:euler-series-resurgence}
\end{equation}
For every fixed nonzero $z$, the ratio of consecutive terms is
\[
  \frac{(n+1)!/z^{n+2}}{n!/z^{n+1}}=\frac{n+1}{z},
\]
whose magnitude eventually exceeds $1$. The series therefore diverges for every
finite $z$. Nevertheless, its first several terms can approximate an exact
solution extremely well when $z$ is large.
\end{example}

\section{Why the terms first decrease and then increase}

Suppose, schematically,
\begin{equation}
  a_n\sim K\,\frac{\Gamma(n+\beta)}{A^{n+\beta}}
  \qquad(n\to\infty),
  \label{eq:factorial-over-power-growth}
\end{equation}
where $A\ne0$. Then the $n$th term of \eqref{eq:generic-inverse-power-series}
has size approximately
\[
  \abs{K}\,
  \frac{\Gamma(n+\beta)}{\abs{A}^{n+\Re\beta}\abs z^{n+1}}.
\]
The ratio of successive magnitudes is approximately
\[
  \frac{n+\Re\beta}{\abs{Az}}.
\]
Thus the terms decrease while $n\lesssim\abs{Az}$ and increase afterward. The
smallest term occurs near
\begin{equation}
  N_*(z)\approx\abs{Az}.
  \label{eq:least-term-index}
\end{equation}
Using Stirling's formula, its size is roughly an algebraic factor times
\begin{equation}
  \e^{-\abs{Az}}.
  \label{eq:least-term-exponential}
\end{equation}

\begin{bigidea}
Factorial divergence and exponentially small scales are two views of the same
phenomenon. If coefficients grow like $n!/A^n$, the optimally truncated error is
usually of size $\e^{-Az}$. The divergent tail is pointing toward the next
transseries sector.
\end{bigidea}

\section{Optimal truncation}

For a convergent series, adding more terms eventually improves the answer. For a
divergent asymptotic series, the sensible rule is different.

\begin{recipebox}
To use a factorially divergent expansion numerically:
\begin{enumerate}
\item Estimate the magnitude of successive terms.
\item Stop near the least term, usually where their ratio crosses $1$.
\item Treat the least-term size as the natural scale of the unresolved
remainder.
\item If greater accuracy is needed, describe that remainder by another
exponential sector rather than by blindly adding more terms.
\end{enumerate}
\end{recipebox}

For the Euler series at real $z>0$, the $n$th term is $n!/z^{n+1}$, so the least
term lies near $n=z$. Stirling's formula gives
\[
  \frac{n!}{z^{n+1}}\bigg|_{n\approx z}
  \asymp \sqrt{\frac{2\pi}{z}}\,\e^{-z}.
\]
The homogeneous solution of the differential equation is exactly $C\e^{-z}$.
The agreement of exponential scales is not an accident.

\section{Gevrey order}

Factorial growth can be measured systematically.

\begin{definition}[Gevrey order]
A formal series $\sum_{n\ge0}a_n z^{-n-1}$ is \emph{Gevrey of order $s$} in a
coefficientwise sense if there are constants $C,A>0$ such that
\begin{equation}
  \abs{a_n}\le C A^n\Gamma(1+sn)
  \qquad(n\ge0).
  \label{eq:gevrey-s-bound}
\end{equation}
The most common case in ordinary differential equations is $s=1$, where the
bound is essentially $CA^n n!$.
\end{definition}

A convergent inverse-power series is Gevrey-$0$ in this informal terminology.
Gevrey-$1$ series are the natural domain of the ordinary Borel transform. Higher
Gevrey orders call for Borel--Leroy transforms or repeated changes of scale,
which we discuss in Chapter~\ref{ch:beyond-one-summability}.

\section{Flat functions explain nonuniqueness}

A function $r(z)$ is \emph{flat} relative to the inverse-power scale in a
sector if
\[
  r(z)=\bigO(z^{-N})
\]
for every $N$. For example, $\e^{-z}$ is flat as $z\to+\infty$. Consequently,
if $f$ has a power-series asymptotic expansion, then so does
$f+C\e^{-z}$, with exactly the same coefficients.

This proves that a Poincar\'e expansion by itself generally does not determine a
function. One must also provide exponentially small data, boundary conditions,
analyticity in a sufficiently wide sector, or a summation prescription.

There is an important refinement. Strong Gevrey bounds in sectors wider than a
critical angle can force uniqueness; this is the content of Watson-type
uniqueness theorems. At the elementary level, the lesson is simply that angular
analytic information can rule out flat additions that are invisible on one ray.

\section{Where factorials come from}

Several common mechanisms create factorial coefficients.

\begin{enumerate}
\item \textbf{Repeated differentiation.} Since
\[
  \frac{\dd^n}{\dd z^n}\frac1z=(-1)^n\frac{n!}{z^{n+1}},
\]
iterating a differential operator naturally produces $n!$.

\item \textbf{Repeated integration by parts.} Laplace-type integrals generate
successive derivatives of an amplitude, again producing factorial growth.

\item \textbf{Nearby singularities.} Taylor coefficients of a function are
controlled by its nearest complex singularities. After the Borel transform,
this elementary principle becomes a direct explanation of large-order growth.

\item \textbf{Combinatorial proliferation.} Nonlinear recursion may create a
number of diagrams, trees, or partitions growing like $n!$.
\end{enumerate}

\section{Large order as a diagnostic tool}

Suppose numerical coefficients satisfy \eqref{eq:factorial-over-power-growth}.
Then
\begin{equation}
  \frac{a_n}{a_{n-1}}
  \sim \frac{n+\beta-1}{A}.
  \label{eq:ratio-large-order}
\end{equation}
A plot of $a_n/a_{n-1}$ against $n$ is therefore approximately a straight line.
Its slope estimates $1/A$, while the intercept estimates $(\beta-1)/A$.
After removing these leading factors, subleading $1/n$ corrections reveal
coefficients of a neighboring transseries sector.

\begin{warningbox}
The nearest action need not be real or unique. A conjugate pair of actions can
make the coefficients oscillate, and several actions of equal modulus can
interfere. Ratio tests then require refinement, such as even/odd subsequences,
Richardson extrapolation, or Borel--Pad\'e analysis.
\end{warningbox}

\begin{summarybox}
A divergent asymptotic series can be highly predictive. Factorial-over-power
coefficient growth determines a least-term index, an exponentially small error
scale, and often the location of a Borel singularity. Divergence is the first
visible trace of the larger transseries.
\end{summarybox}


\chapter{The Borel--Laplace microscope}
\label{ch:borel-laplace}

\section{The basic trick}

The ordinary Borel transform divides out the factorial growth that causes a
Gevrey-$1$ series to diverge. For
\[
  \phihat(z)=\sum_{n=0}^{\infty}\frac{a_n}{z^{n+1}},
\]
define
\begin{equation}
  (\Borel\phihat)(\xi)
  =\widehat\varphi(\xi)
  \defeq \sum_{n=0}^{\infty}\frac{a_n}{n!}\xi^n.
  \label{eq:borel-transform-definition}
\end{equation}
If $a_n$ grows no faster than $CA^n n!$, the new series has a positive radius
of convergence, at least $1/A$.

The directional Laplace transform reverses the operation:
\begin{equation}
  (\Laplace_\theta\widehat\varphi)(z)
  \defeq
  \int_{0}^{\e^{i\theta}\infty}
     \e^{-z\xi}\widehat\varphi(\xi)\dd\xi,
  \label{eq:directional-laplace}
\end{equation}
provided the ray avoids singularities, the continued Borel transform grows
slowly enough, and $\Re(z\e^{i\theta})>0$.

Indeed,
\begin{equation}
  \int_{0}^{\e^{i\theta}\infty}
      \e^{-z\xi}\xi^n\dd\xi
  =\frac{n!}{z^{n+1}},
  \label{eq:laplace-monomial}
\end{equation}
with compatible choices of direction. Formally applying
$\Laplace_\theta\Borel$ therefore restores the original series term by term.

\begin{definition}[Directional Borel sum]
If the Borel transform converges near $0$, can be analytically continued along
the ray of angle $\theta$, and has at most exponential growth there, then
\begin{equation}
  \mathcal S_\theta\phihat
  \defeq \Laplace_\theta\Borel\phihat
  \label{eq:directional-borel-sum}
\end{equation}
is the directional Borel sum of $\phihat$.
\end{definition}

A constant term $c$ is simply carried along:
$\mathcal S_\theta(c+\phihat)=c+\mathcal S_\theta\phihat$.

\section{A three-stage journey}

It helps to keep the domains separate.

\begin{center}
\begin{tikzpicture}[node distance=18mm,>=Latex]
  \node[draw,rounded corners,fill=softblue,minimum width=34mm,minimum height=13mm] (formal)
    {formal $z^{-1}$ series};
  \node[draw,rounded corners,fill=softgreen,right=of formal,minimum width=34mm,minimum height=13mm] (borel)
    {Borel germ at $\xi=0$};
  \node[draw,rounded corners,fill=softorange,right=of borel,minimum width=34mm,minimum height=13mm] (function)
    {analytic function of $z$};
  \draw[->,thick] (formal) -- node[above]{divide by $n!$} node[below]{$\Borel$} (borel);
  \draw[->,thick] (borel) -- node[above]{continue and integrate} node[below]{$\Laplace_\theta$} (function);
\end{tikzpicture}
\end{center}

The first arrow is algebraic. The difficult information lies in the middle:
how can the germ be analytically continued, where are its singularities, and
how does it grow? The second arrow converts this Borel-plane geometry into
exponentially small behavior in the original variable.

\section{The Euler series becomes a rational function}

For \eqref{eq:euler-series-resurgence},
\begin{equation}
  \Borel\widehat E(\xi)
   =\sum_{n=0}^{\infty}\xi^n
   =\frac1{1-\xi}.
  \label{eq:euler-borel-transform}
\end{equation}
A violently divergent factorial series has become the simplest possible
meromorphic function. Its first singularity is a pole at $\xi=1$.

If the integration ray avoids the positive real axis, then
\begin{equation}
  \mathcal S_\theta\widehat E(z)
  =\int_{0}^{\e^{i\theta}\infty}
       \frac{\e^{-z\xi}}{1-\xi}\dd\xi.
  \label{eq:euler-directional-sum}
\end{equation}
Differentiating under the integral sign verifies that this is an actual solution
of
\[
  y'+y=\frac1z.
\]
The Borel transform has therefore converted a formal solution into analytic
solutions.

\section{Why Borel singularities create exponentials}

Suppose the Borel transform has a singularity at $\xi=A$. In a Laplace integral,
translate $\xi=A+u$ near that point:
\[
  \e^{-z\xi}=\e^{-Az}\e^{-zu}.
\]
Every contribution localized near $A$ carries the factor $\e^{-Az}$. Thus
\begin{equation}
  \boxed{\text{Borel singularity at }A
  \quad\longleftrightarrow\quad
  \text{exponential scale }\e^{-Az}.}
  \label{eq:borel-singularity-dictionary}
\end{equation}
This is the basic Borel-plane dictionary behind resurgence.

\section{Watson's lemma}

The Laplace transform automatically has the original formal expansion.

\begin{theorem}[Watson's lemma, elementary form]
Suppose $\varphi$ is analytic near $0$ with
\[
  \varphi(\xi)=\sum_{n=0}^{N-1}b_n\xi^n+\bigO(\xi^N)
\]
and has sufficient growth control on a Laplace ray. Then
\begin{equation}
  \int_0^{\e^{i\theta}\infty}\e^{-z\xi}\varphi(\xi)\dd\xi
  =\sum_{n=0}^{N-1}\frac{b_n n!}{z^{n+1}}
   +\bigO(z^{-N-1})
  \label{eq:watson-lemma}
\end{equation}
inside every closed subsector where the integral converges uniformly.
\end{theorem}

\begin{proof}[Proof idea]
Split the ray at a small fixed distance $\rho$. On $0\le\abs\xi\le\rho$,
insert the Taylor polynomial and integrate each monomial using
\eqref{eq:laplace-monomial}. The Taylor remainder contributes
$\bigO(z^{-N-1})$ after the substitution $u=z\xi$. On the remaining part of the
ray, the exponential kernel supplies a bound of order $\e^{-c\abs z}$, which is
smaller than every inverse power. The details are estimates, not new ideas.
\end{proof}

\section{Algebra in the Borel plane}

Multiplication of formal series becomes convolution. If
\[
  \phihat(z)=\sum_{n\ge0}a_n z^{-n-1},
  \qquad
  \widehat\psi(z)=\sum_{n\ge0}b_n z^{-n-1},
\]
then
\begin{equation}
  \Borel(\phihat\widehat\psi)
  =(\Borel\phihat)*(\Borel\widehat\psi),
  \qquad
  (f*g)(\xi)=\int_0^\xi f(\eta)g(\xi-\eta)\dd\eta.
  \label{eq:borel-convolution}
\end{equation}
The beta integral
\[
  \int_0^\xi \eta^m(\xi-\eta)^n\dd\eta
  =\frac{m!n!}{(m+n+1)!}\xi^{m+n+1}
\]
proves the identity coefficient by coefficient.

Convolution explains why nonlinear equations produce new singularities at sums
of old singularities. If one factor is singular near $A$ and another near $B$,
their convolution can be singular near $A+B$. In the original variable, this is
just multiplication of exponentials:
\[
  \e^{-Az}\e^{-Bz}=\e^{-(A+B)z}.
\]

\section{What can go wrong}

Borel summation is powerful but not automatic.

\begin{enumerate}
\item The Borel series may have zero radius of convergence because growth is
faster than $n!$.
\item A singularity may lie on the desired integration ray.
\item Analytic continuation may exist only along restricted paths.
\item The continued Borel transform may grow too fast for an ordinary Laplace
integral.
\item Different lateral paths may give different answers.
\end{enumerate}

The next chapters turn the fourth and fifth items from obstacles into useful
structure.

\section{Borel--Pad\'e summation in practice}

Often only finitely many coefficients $a_0,\ldots,a_N$ are known. One can still
probe the Borel plane.

\begin{recipebox}
A basic Borel--Pad\'e calculation proceeds as follows.
\begin{enumerate}
\item Form the truncated Borel polynomial
$\sum_{n=0}^N a_n\xi^n/n!$.
\item Replace it by a near-diagonal Pad\'e approximant, a rational function whose
Taylor series matches as many coefficients as possible.
\item Inspect stable strings or clusters of poles as approximations to branch
cuts and isolated singularities.
\item Integrate the rational approximation along a safe ray.
\item Vary the Pad\'e order and ray slightly to test numerical stability.
\end{enumerate}
\end{recipebox}

Spurious pole-zero pairs, called Froissart doublets, can appear. A numerical plot
is evidence, not a proof of analytic continuation. Nevertheless, Borel--Pad\'e
analysis is often an excellent microscope for discovering the actions and
Stokes directions of an unknown problem.

\begin{summarybox}
The Borel transform removes a factorial, turning a divergent Gevrey-$1$ series
into a locally convergent germ. Analytic continuation exposes singularities;
the Laplace transform returns to the original variable. A singularity at $A$
produces the scale $\e^{-Az}$, and multiplication becomes convolution.
\end{summarybox}


\chapter{Stokes phenomena: when an invisible term switches on}
\label{ch:stokes}

\section{A singular direction}

For the Euler Borel transform $1/(1-\xi)$, the positive real Laplace ray runs
through the pole at $\xi=1$. There are two natural detours:
\begin{align}
  \mathcal S_{0^+}\widehat E(z)
   &=\int_{0}^{\infty\e^{+i0}}
      \frac{\e^{-z\xi}}{1-\xi}\dd\xi,\\
  \mathcal S_{0^-}\widehat E(z)
   &=\int_{0}^{\infty\e^{-i0}}
      \frac{\e^{-z\xi}}{1-\xi}\dd\xi.
  \label{eq:euler-lateral-sums}
\end{align}
The notation means that the contour passes just above or just below the positive
axis.

\begin{center}
\begin{tikzpicture}[x=1.1cm,y=1.1cm,>=Latex]
  \draw[->] (-0.2,0) -- (5.6,0) node[right] {$\Re\xi$};
  \draw[->] (0,-1.0) -- (0,1.4) node[above] {$\Im\xi$};
  \fill[deepred] (2.6,0) circle (2.2pt) node[below=5pt] {$A=1$};
  \draw[deepblue,thick,->] (0,0.18) -- (2.15,0.18)
      .. controls (2.35,0.18) and (2.38,0.55) .. (2.6,0.55)
      .. controls (2.82,0.55) and (2.86,0.18) .. (3.05,0.18) -- (5.2,0.18);
  \draw[deepgreen,thick,->] (0,-0.18) -- (2.15,-0.18)
      .. controls (2.35,-0.18) and (2.38,-0.55) .. (2.6,-0.55)
      .. controls (2.82,-0.55) and (2.86,-0.18) .. (3.05,-0.18) -- (5.2,-0.18);
  \node[deepblue] at (4.5,0.52) {$0^+$};
  \node[deepgreen] at (4.5,-0.52) {$0^-$};
\end{tikzpicture}
\end{center}

\section{The Euler jump by one residue}

The difference of the two contours encloses the pole. With the orientation in
\eqref{eq:euler-lateral-sums}, the residue theorem gives
\begin{equation}
  \mathcal S_{0^+}\widehat E(z)
  -\mathcal S_{0^-}\widehat E(z)
  =2\pi i\,\e^{-z}.
  \label{eq:euler-stokes-jump}
\end{equation}
Indeed, the residue of $\e^{-z\xi}/(1-\xi)$ at $\xi=1$ is $-\e^{-z}$, and the
small contour resulting from ``upper minus lower'' is clockwise, supplying the
second minus sign.

The two sums have the same inverse-power expansion. Their difference is flat,
exactly as the discussion in Chapter~\ref{ch:factorial-divergence} predicted.

\section{The transseries parameter absorbs the jump}

The complete formal solution of the Euler equation is
\begin{equation}
  \widehat y(z;C)=\widehat E(z)+C\e^{-z}.
  \label{eq:euler-full-transseries-stokes}
\end{equation}
To represent the same analytic solution above and below the Stokes ray, require
\[
  \mathcal S_{0^+}\widehat E+C_+\e^{-z}
  =\mathcal S_{0^-}\widehat E+C_-\e^{-z}.
\]
Using \eqref{eq:euler-stokes-jump}, we obtain
\begin{equation}
  C_+=C_--2\pi i.
  \label{eq:euler-parameter-jump}
\end{equation}
The function itself need not jump. What jumps is the asymptotic coordinate $C$
used to describe it.

\begin{bigidea}
A Stokes phenomenon is a change of asymptotic coordinates. A contribution that
was exponentially subdominant can acquire a new coefficient when a singular
direction is crossed. The full transseries remains consistent because its
parameters transform at the same time.
\end{bigidea}

\section{Stokes and anti-Stokes terminology}

For an exponential $\e^{-Az}$, two geometrically important sets of rays occur:
\begin{align*}
  \Im(Az)&=0, & &\text{the exponential has constant phase},\\
  \Re(Az)&=0, & &\text{the exponential changes between decay and growth}.
\end{align*}
Different communities interchange the names \emph{Stokes line} and
\emph{anti-Stokes line} for these sets. To avoid ambiguity, we will say
\emph{Borel singular direction} for a Laplace direction hitting a singularity,
and \emph{equal-magnitude line} when two saddle exponentials have equal
magnitude.

\section{Dominance can change without discontinuity}

Consider
\[
  F(z)=1+\e^{-z}.
\]
For $\Re z>0$, the exponential is small; for $\Re z<0$, it is large. The
function is entire and does not jump at $\Re z=0$. Only the useful asymptotic
organization changes. A genuine Stokes jump is subtler: it concerns how analytic
continuation and asymptotic representation interact when an exponentially small
sector is switched into the expansion.

\section{The Stokes automorphism}

Let $\mathcal S_{\theta^-}$ and $\mathcal S_{\theta^+}$ denote lateral
resummations of a whole formal transseries just below and above a singular
direction. We define the formal Stokes automorphism $\mathfrak S_\theta$ by the
convention
\begin{equation}
  \mathcal S_{\theta^+}
  =\mathcal S_{\theta^-}\compose\mathfrak S_\theta.
  \label{eq:stokes-automorphism-convention}
\end{equation}
It is an automorphism because it must respect addition, multiplication, and the
formal equations satisfied by the transseries. In the Euler example it acts by
a translation of the transseries parameter.

The exact sign of a quoted ``Stokes constant'' depends on conventions: direction
of the Laplace kernel, orientation of lateral contours, normalization of the
exponential sector, and whether the automorphism maps lower data to upper data
or conversely. A reliable calculation always states these choices.

\section{Smooth switching at finite size}

The discontinuous Stokes jump belongs to an ideal asymptotic description as
$\abs z\to\infty$. At large but finite $\abs z$, optimal truncation reveals a
rapid yet smooth transition, typically on an angular scale of order
$\abs z^{-1/2}$. In many classical problems the multiplier is approximated by
an error function. This is called Berry smoothing.

There is no contradiction. A step function is the limiting profile of a smooth
transition whose width shrinks as the asymptotic parameter grows.

\begin{summarybox}
When a Laplace ray crosses a Borel singularity, its two lateral sums can differ
by an exponentially small term. The corresponding transseries parameter changes
so that one analytic solution has compatible descriptions on both sides. Stokes
jumps are changes of asymptotic coordinates, not mysterious discontinuities of
the underlying function.
\end{summarybox}


\chapter{Airy functions: two saddles and a Stokes geometry}
\label{ch:airy}

\section{Why Airy's equation is the standard laboratory}

Airy's equation
\begin{equation}
  y''(x)-x y(x)=0
  \label{eq:airy-equation}
\end{equation}
is linear, exactly solvable, and already contains the geometry of competing
exponentials. It describes the universal local behavior near a simple turning
point in many differential equations.

Seek a WKB-type solution
\[
  y(x)=\exp(S(x)).
\]
Then
\[
  S''+(S')^2=x.
\]
At leading order, $(S')^2\sim x$, so
\begin{equation}
  S(x)\sim\pm\frac23 x^{3/2}.
  \label{eq:airy-actions}
\end{equation}
The two fundamental exponential scales are therefore
\[
  \e^{-\zeta},\qquad \e^{+\zeta},
  \qquad \zeta=\frac23 x^{3/2}.
\]

\section{The algebraic prefactor}

Write
\[
  y(x)=\e^{\lambda\zeta}x^\alpha u(x),
  \qquad \lambda=\pm1,
  \qquad u(x)\sim1+\cdots.
\]
Substitution into \eqref{eq:airy-equation} shows that the next balance requires
\[
  \alpha=-\frac14.
\]
Thus the two formal WKB sectors begin as
\begin{equation}
  y_\pm(x)
  \sim x^{-1/4}\e^{\pm\frac23x^{3/2}}
  \left(1+\sum_{n\ge1}\frac{c_n^{(\pm)}}{x^{3n/2}}\right).
  \label{eq:airy-wkb-sectors}
\end{equation}
The coefficient series are divergent of Gevrey order $1$ in the variable
$\zeta^{-1}$.

For real $x>0$, conventional normalizations give
\begin{align}
  \Ai(x)
  &\asympseries
  \frac{\e^{-\zeta}}{2\sqrt\pi\,x^{1/4}}
  \left(1-\frac{5}{48x^{3/2}}
    +\frac{385}{4608x^3}-\cdots\right),
  \label{eq:airy-ai-positive}\\
  \Bi(x)
  &\asympseries
  \frac{\e^{+\zeta}}{\sqrt\pi\,x^{1/4}}
  \left(1+\frac{5}{48x^{3/2}}
    +\frac{385}{4608x^3}+\cdots\right).
  \label{eq:airy-bi-positive}
\end{align}

\section{The contour integral and its saddles}

One representation is
\begin{equation}
  \Ai(x)=\frac{1}{2\pi i}
  \int_{\mathcal C}
    \exp\!\left(\frac{t^3}{3}-xt\right)\dd t,
  \label{eq:airy-contour-integral}
\end{equation}
where $\mathcal C$ runs between sectors in which the cubic exponential decays.
The phase
\[
  \Phi(t)=\frac{t^3}{3}-xt
\]
has critical points
\[
  \Phi'(t)=t^2-x=0,
  \qquad t_\pm=\pm\sqrt x.
\]
Their phase values are
\[
  \Phi(t_+)=-\frac23x^{3/2}=-\zeta,
  \qquad
  \Phi(t_-)=+\frac23x^{3/2}=+\zeta.
\]
These are precisely the two WKB exponentials.

A steepest-descent contour can be decomposed into elementary contours, called
Lefschetz thimbles, attached to the saddles. As $\arg x$ changes, the valid
thimble decomposition changes by integer combinations. This geometric jump is
the integral version of a Stokes automorphism.

\section{Action differences matter}

The difference between the saddle actions is
\begin{equation}
  2\zeta=\frac43x^{3/2}.
  \label{eq:airy-action-difference}
\end{equation}
Consequently, the late coefficients of the expansion around one saddle know
about the other saddle through an exponential scale
$\e^{-2\zeta}$. Relative to the dominant factor already extracted from a sector,
this is the size of the neighboring saddle contribution.

\section{Rotating the argument}

Let $x=re^{i\phi}$. Then
\[
  \zeta=\frac23r^{3/2}e^{3i\phi/2}.
\]
The two exponentials have equal magnitude when $\Re\zeta=0$, that is, on rays
\[
  \phi=\frac{\pi}{3},\ \pi,\ -\frac{\pi}{3}
  \pmod{2\pi}.
\]
Between these rays, one saddle is exponentially dominant and the other
subdominant. Analytic continuation of $\Ai$ is smooth, but its most economical
asymptotic representation changes from sector to sector.

\begin{intuitionbox}
Imagine describing a landscape at night with two lamps. In one region the first
lamp overwhelms the second, so the second is invisible to a power-series camera.
After rotating the landscape, the dim lamp becomes comparable or dominant. The
lamp did not suddenly appear; only the hierarchy of brightness changed.
\end{intuitionbox}

\section{What the Airy example teaches}

Airy's equation packages several universal ideas:
\begin{enumerate}
\item exponential actions arise by integrating a leading momentum,
$\int\sqrt{x}\dd x$;
\item algebraic prefactors arise at the next asymptotic order;
\item each saddle carries a divergent fluctuation series;
\item differences of saddle actions govern large-order behavior;
\item Stokes jumps are changes in contour or saddle decomposition.
\end{enumerate}
The same pattern survives in far more complicated ordinary differential
equations, integrals, quantum-mechanical problems, and path integrals.

\begin{summarybox}
Airy's equation has two formal sectors, proportional to
$x^{-1/4}\e^{\pm2x^{3/2}/3}$ times divergent series. Its contour integral has two
matching saddles. Rotating $x$ changes which saddle is dominant and how the
integration contour decomposes, giving a concrete geometric model of Stokes
phenomena.
\end{summarybox}


\chapter{Resurgence: sectors talking to one another}
\label{ch:resurgence}

\section{From isolated summation to a network}

Borel summation treats one divergent series. Resurgence adds a stronger claim:
the singularities encountered while continuing its Borel transform are not
arbitrary obstacles. Their local expansions reproduce the Borel transforms of
other sectors in the same problem. Information seems to ``re-emerge'' at those
singular points, which motivated \'Ecaille's term \emph{resurgence}.

A typical one-parameter transseries has the form
\begin{equation}
  \widehat\Phi(z,\sigma)
  =\sum_{k=0}^{\infty}
    \sigma^k\e^{-kAz}z^{k\beta}\phihat_k(z),
  \qquad
  \phihat_k(z)=\sum_{n=0}^{\infty}\frac{a_n^{(k)}}{z^{n+1}}.
  \label{eq:one-parameter-resurgent-transseries}
\end{equation}
The sector $k=0$ is often called perturbative. The sectors $k\ge1$ are
nonperturbative or multi-instanton sectors. The terminology comes from physics,
but the algebra applies to ordinary differential equations and integrals just
as well.

\section{Endless continuability, informally}

A convergent germ $\widehat\varphi(\xi)$ at the Borel origin is called resurgent
when it can be analytically continued along every finite-length path after one
avoids a locally finite collection of singular points, with suitable control on
the resulting singularities. Precise definitions differ according to the class
of singularities and the amount of uniformity required
\cite{EcalleResurgent,SauzinIntro,DorigoniIntro}.

The adjective \emph{endless} means that continuation can be repeated: reaching
one singularity does not terminate the story. One can go around it, examine the
new germ, and continue farther.

\section{The simplest large-order calculation}

Suppose the Borel transform of $\phihat_0$ has its nearest singularity at
$\xi=A$ and, near that point, behaves like a simple pole
\begin{equation}
  \widehat\varphi_0(\xi)
  \sim \frac{r}{\xi-A}.
  \label{eq:simple-borel-pole}
\end{equation}
Expanding the pole about the origin gives
\[
  \frac{r}{\xi-A}
  =-\frac{r}{A}\frac1{1-\xi/A}
  =-r\sum_{n=0}^{\infty}\frac{\xi^n}{A^{n+1}}.
\]
Since the coefficient of $\xi^n$ in the Borel transform is
$a_n^{(0)}/n!$, we expect
\begin{equation}
  a_n^{(0)}\sim-r\frac{n!}{A^{n+1}}.
  \label{eq:simple-pole-large-order}
\end{equation}
For the Euler series, $A=1$ and $r=-1$, so this gives $a_n=n!$ exactly.

This elementary argument is Darboux's principle in Borel clothing: the nearest
singularity controls the late Taylor coefficients.

\section{A neighboring sector inside a singularity}

A more characteristic resurgent singularity is not merely a pole. Schematically,
one may find
\begin{equation}
  \widehat\varphi_0(A+\xi)
  \sim
  \frac{S_{01}}{2\pi i}
  \widehat\varphi_1(\xi)\log\xi
  +\text{less singular terms},
  \label{eq:resurgent-local-model}
\end{equation}
where $S_{01}$ is a Stokes constant and
$\widehat\varphi_1=\Borel\phihat_1$. Expanding the right-hand side about the
origin and applying coefficient asymptotics yields a full large-order relation
of the form
\begin{equation}
  a_n^{(0)}
  \sim \frac{S_{01}}{2\pi i}
  \sum_{m\ge0}a_m^{(1)}
  \frac{\Gamma(n-m+\gamma)}{A^{n-m+\gamma}}
  +\text{contributions from other singularities}.
  \label{eq:generic-large-order-resurgence}
\end{equation}
The shift $\gamma$ depends on the powers of $z$ used to normalize the two
sectors. Formula \eqref{eq:generic-large-order-resurgence} is schematic but its
message is exact: \emph{late coefficients of one sector are determined by early
coefficients of another}.

\section{Reading the neighboring sector from data}

Suppose $A$, $\gamma$, and the leading Stokes constant have been estimated. Then
one can peel off the leading factorial growth:
\[
  \frac{2\pi i}{S_{01}}
  \frac{A^{n+\gamma}a_n^{(0)}}{\Gamma(n+\gamma)}
  \sim a_0^{(1)}
  +\frac{A a_1^{(1)}}{n+\gamma-1}
  +\frac{A^2a_2^{(1)}}{(n+\gamma-1)(n+\gamma-2)}+
eedspace{1em}\cdots.
\]
Successive $1/n$ corrections recover $a_1^{(1)},a_2^{(1)},\ldots$. In good
examples, dozens of low-order perturbative coefficients can numerically predict
the first few nonperturbative coefficients with striking accuracy.

\section{Why nonlinear equations generate a lattice of actions}

Suppose an equation contains products of sectors with actions $A$ and $B$.
Their product has action $A+B$. Iterating nonlinearities creates an additive
semigroup
\[
  \Omega=\{n_1A_1+\cdots+n_rA_r:n_j\in\N\}.
\]
Correspondingly, Borel singularities tend to occur at points of $\Omega$ and
its signed or analytically continued variants. This is the analytic counterpart
of the transmonomial lattice generated by
$\e^{-A_1z},\ldots,\e^{-A_rz}$.

\begin{warningbox}
The action lattice can have coincidences. If two different integer combinations
give the same action, the problem is resonant. Logarithmic sectors and polynomial
factors in $z$ may then be forced. A naive ansatz containing only
$\e^{-kAz}$ times ordinary power series may be incomplete.
\end{warningbox}

\section{Resurgence does not mean convergence}

A resurgent series is usually still divergent. Resurgence says that its Borel
transform has a controlled and highly organized continuation, not that the
original series converges. Nor is every divergent series resurgent. Some Borel
transforms have natural boundaries or singular sets too dense for endless
continuation.

\section{The conceptual triangle}

Three structures determine one another:
\begin{center}
\begin{tikzpicture}[>=Latex,node distance=18mm]
  \node[draw,rounded corners,fill=softblue,align=center] (late)
    {late coefficients\\of sector $0$};
  \node[draw,rounded corners,fill=softgreen,align=center,right=36mm of late] (sing)
    {singularity near\\$\xi=A$};
  \node[draw,rounded corners,fill=softorange,align=center,below=24mm of $(late)!0.5!(sing)$] (next)
    {early coefficients\\of sector $1$};
  \draw[<->,thick] (late) -- node[above]{Darboux/large order} (sing);
  \draw[<->,thick] (sing) -- node[right]{local resurgence} (next);
  \draw[<->,thick] (next) -- node[left]{large-order relation} (late);
\end{tikzpicture}
\end{center}
This triangle is the practical heart of resurgence.

\begin{summarybox}
A resurgent transseries is a network of sectors. Borel singularities of one
sector encode other sectors, and the resulting local singular expansions imply
large-order relations among their coefficients. Nonlinearity adds actions and
therefore creates whole lattices of exponential sectors and Borel singularities.
\end{summarybox}


\chapter{Alien derivatives and bridge equations}
\label{ch:alien-calculus}

\section{Why invent another derivative?}

An ordinary derivative measures local variation with respect to $z$. To study
resurgence, we need an operator that measures the singular behavior reached by
analytic continuation in the Borel plane. \'Ecaille's \emph{alien derivative}
$\Delta_\omega$ extracts the new resurgent information located at the Borel
singularity $\omega$.

The word \emph{alien} emphasizes that the operator is nonlocal in the original
variable. It can detect an exponentially small sector that no finite number of
ordinary derivatives of a power-series asymptotic expansion would reveal.

A full construction requires the algebra of singular germs and careful choices
of paths. Here we focus on the rules that make the operator useful.

\section{A schematic extraction rule}

If the continued Borel transform of $\phihat_0$ near $\omega$ contains
\[
  \frac{S_{01}}{2\pi i}\widehat\varphi_1(\xi-\omega)
  \log(\xi-\omega),
\]
then, under a common normalization,
\begin{equation}
  \Delta_\omega\phihat_0=S_{01}\phihat_1.
  \label{eq:alien-extracts-sector}
\end{equation}
More complicated singularities produce combinations of sectors. The numerical
factor assigned to $S_{01}$ varies among conventions, but the invariant content
is the coupling between sectors.

Alien derivatives satisfy a Leibniz rule in their standard logarithmic
normalization:
\begin{equation}
  \Delta_\omega(FG)
  =(\Delta_\omega F)G+F(\Delta_\omega G).
  \label{eq:alien-leibniz}
\end{equation}
This makes them derivations of the resurgent algebra.

\section{Pointed alien derivatives}

It is often convenient to attach the exponential weight to the operator:
\begin{equation}
  \dot\Delta_\omega
  \defeq \e^{-\omega z}\Delta_\omega.
  \label{eq:pointed-alien-derivative}
\end{equation}
With a standard choice of alien normalization, the pointed derivative is chosen
to commute with ordinary differentiation,
\begin{equation}
  [\partial_z,\dot\Delta_\omega]=0.
  \label{eq:pointed-commutator}
\end{equation}
Other references place the exponential weight or sign elsewhere. The safe
practice is to check the stated commutator before importing a formula.

\section{The Stokes automorphism as an exponential}

Along a singular direction $\theta$, let $\Gamma_\theta$ be the set of Borel
singularities on the ray. Formally, the Stokes automorphism can be written
\begin{equation}
  \mathfrak S_\theta
  =\exp\!\left(
      \sum_{\omega\in\Gamma_\theta}
      \e^{-\omega z}\Delta_\omega
    \right)
  =\exp\!\left(
      \sum_{\omega\in\Gamma_\theta}
      \dot\Delta_\omega
    \right),
  \label{eq:stokes-auto-alien}
\end{equation}
with path-ordering or refined ``lateral'' alien operators inserted when the
singularity geometry demands it. This resembles the familiar identity
``finite translation equals the exponential of a derivative.'' The alien
derivatives are infinitesimal generators of a finite Stokes jump.

\section{The bridge equation}

Now suppose $\widehat\Phi(z,\sigma)$ is a formal solution of a differential
equation. Apply $\dot\Delta_\omega$ to that equation. Because the pointed alien
derivative commutes with $\partial_z$ and obeys a Leibniz rule, the result solves
the linearized equation around $\widehat\Phi$. But differentiating the original
solution with respect to its free parameter $\sigma$ also solves that linearized
equation. Under suitable one-dimensionality assumptions, the two must be
proportional:
\begin{equation}
  \dot\Delta_\omega\widehat\Phi(z,\sigma)
  =S_\omega(\sigma)\,
    \partial_\sigma\widehat\Phi(z,\sigma).
  \label{eq:bridge-equation-general}
\end{equation}
This is a \emph{bridge equation}: it bridges analytic continuation in the Borel
plane and ordinary differentiation in transseries-parameter space.

For the simplest one-action problem, the first relation may reduce to
\begin{equation}
  \dot\Delta_A\widehat\Phi
  =S_1\partial_\sigma\widehat\Phi,
  \label{eq:simple-bridge-equation}
\end{equation}
where $S_1$ is a Stokes constant.

\section{Parameter translation}

Insert \eqref{eq:simple-bridge-equation} into
\eqref{eq:stokes-auto-alien}. If no other alien derivatives contribute along
the ray, then
\[
  \mathfrak S_0
  =\exp(S_1\partial_\sigma).
\]
But
\[
  \exp(S_1\partial_\sigma)F(\sigma)=F(\sigma+S_1),
\]
so
\begin{equation}
  \mathfrak S_0\widehat\Phi(z,\sigma)
  =\widehat\Phi(z,\sigma+S_1).
  \label{eq:stokes-parameter-translation}
\end{equation}
The Euler parameter jump in \eqref{eq:euler-parameter-jump} is the elementary
prototype of this general statement, after conventions and normalizations are
matched.

\section{Median summation}

Assume real coefficients and a Stokes ray along the positive real direction.
The two lateral Borel sums are often complex conjugates, even though the desired
answer is real. With the convention
\eqref{eq:stokes-automorphism-convention}, define the median resummation by
\begin{equation}
  \mathcal S_{\mathrm{med}}
  =\mathcal S_{0^-}\compose\mathfrak S_0^{1/2}
  =\mathcal S_{0^+}\compose\mathfrak S_0^{-1/2}.
  \label{eq:median-resummation}
\end{equation}
It divides the Stokes jump symmetrically between the two sides. Under appropriate
reality conditions, the result is real and ambiguity-free.

For a pure parameter translation $\sigma\mapsto\sigma+S_1$, the half-power
translates by $S_1/2$. Thus median summation can be understood as choosing a
half-shifted transseries parameter before taking a lateral sum.

\section{Several parameters and resonance}

With actions $A_1,\ldots,A_r$, introduce parameters
$\bm\sigma=(\sigma_1,\ldots,\sigma_r)$ and sectors
\[
  \bm\sigma^{\bm n}
  \e^{-(\bm n\cdot\bm A)z}\phihat_{\bm n}(z).
\]
Bridge equations become vector fields on parameter space:
\begin{equation}
  \dot\Delta_\omega\widehat\Phi
  =\sum_{j=1}^r
    S_{\omega,j}(\bm\sigma)
    \frac{\partial\widehat\Phi}{\partial\sigma_j}.
  \label{eq:multivariate-bridge}
\end{equation}
The Stokes automorphism is then the time-one flow of a generally nonlinear
formal vector field.

When action combinations coincide, logarithms and additional parameters may be
necessary. Alien calculus still works, but the parameter-space vector field can
mix sectors in a triangular or resonant fashion.

\begin{warningbox}
Alien calculus is convention-sensitive notation for invariant analytic facts.
Never compare bare Stokes constants from two sources until the sector
normalizations, Borel transform, Laplace kernel, lateral orientation, and bridge
equation convention have all been aligned.
\end{warningbox}

\begin{summarybox}
Alien derivatives extract singular information at specified Borel locations.
Their exponential generates Stokes automorphisms. Bridge equations identify
alien differentiation with vector fields in transseries-parameter space, often
turning a Stokes jump into an ordinary parameter translation. Median summation
uses half of this transformation to cancel lateral ambiguity symmetrically.
\end{summarybox}


\chapter{Beyond one Borel sum}
\label{ch:beyond-one-summability}

\section{Borel--Leroy transforms}

If coefficients grow like $\Gamma(n+\beta)$ rather than exactly $n!$, it is
convenient to divide by that gamma factor. One version of the Borel--Leroy
transform is
\begin{equation}
  \Borel_\beta
  \left(\sum_{n\ge0}a_n z^{-n-\beta}\right)
  =\sum_{n\ge0}\frac{a_n}{\Gamma(n+\beta)}\xi^{n+\beta-1}.
  \label{eq:borel-leroy}
\end{equation}
The inverse is still a Laplace integral because
\[
  \int_0^\infty\e^{-z\xi}\xi^{n+\beta-1}\dd\xi
  =\frac{\Gamma(n+\beta)}{z^{n+\beta}}.
\]
This formulation keeps algebraic prefactors attached to an exponential sector
without awkward index shifts.

\section{\texorpdfstring{$k$}{k}-summability and higher Gevrey order}

For growth comparable to $\Gamma(1+n/k)$ or, after a reciprocal convention,
$\Gamma(1+kn)$, one uses a Borel transform of order $k$ together with a Laplace
kernel such as $\exp[-(z\xi)^k]$. The exact formulas vary with normalization.
The principle is stable:
\begin{enumerate}
\item divide out the characteristic gamma growth;
\item analytically continue in the transformed plane;
\item use a matching integral kernel to restore the original monomials.
\end{enumerate}

Different levels of an equation can possess different Gevrey orders. A single
ordinary Borel transform then cannot regularize every level at once.

\section{Multisummability}

A formal series is \emph{multisummable} when it can be decomposed into finitely
many pieces, each summable at its own Gevrey level, in a compatible sequence of
directions. Roughly,
\[
  \phihat=\phihat^{(k_1)}+\cdots+\phihat^{(k_r)},
  \qquad k_1<\cdots<k_r,
\]
and each component is $k_j$-summable after the faster levels have been handled.
Multisummability is a fundamental theorem for broad classes of formal solutions
of analytic differential equations.

\section{Acceleration}

Sometimes the analytically continued Borel transform grows faster than any
ordinary exponential, so the Laplace integral still diverges. \'Ecaille's
\emph{acceleration} changes the Borel variable and kernel to slow this growth,
then continues the summation process. One can think of it as replacing an
inadequate microscope by one adapted to a more rapidly separated asymptotic
scale.

Acceleration is one reason the full theory of transseries goes beyond the
simple slogan ``Borel transform and Laplace transform.'' Yet the slogan remains
the correct first model.

\section{Hyperasymptotics}

Optimal truncation leaves an exponentially small remainder. Hyperasymptotics
expands that remainder using contributions from adjacent saddles or Borel
singularities, then optimally truncates those expansions in turn.

For a Laplace-type integral with saddles $s_0,s_1,\ldots$, the first expansion
around $s_0$ has a remainder controlled by the nearest action difference
$A_{01}$. One rewrites the remainder as an integral associated with $s_1$,
expands it, and obtains an error of a still smaller exponential order. Repeating
this process creates a hierarchy:
\[
  \text{power accuracy}
  \quad\longrightarrow\quad
  \e^{-A_{01}z}\text{ accuracy}
  \quad\longrightarrow\quad
  \e^{-(A_{01}+A_{12})z}\text{ accuracy}
  \quad\longrightarrow\cdots.
\]
The resulting special functions that smooth truncated saddle contributions are
often called terminants.

\section{Exponential asymptotics}

Exponential asymptotics is a practical methodology, especially common in
singular perturbation theory:
\begin{enumerate}
\item compute the late-term ansatz of a divergent regular expansion;
\item infer the singulant $A(x)$ and prefactor from its recurrence;
\item truncate near the least term;
\item derive a differential equation for the exponentially small remainder;
\item resolve its rapid change across a Stokes curve, often obtaining an error
function profile.
\end{enumerate}
This can reveal beyond-all-orders waves, capillary ripples, boundary-layer
leakage, and other effects without constructing the complete alien calculus.

\section{Transasymptotic rearrangement}

A sector expansion
\[
  \sum_{k\ge0}\sigma^k\e^{-kAz}\phihat_k(z)
\]
is ordered by exponentials when $\sigma$ is fixed and $\Re(Az)\gg1$. Near a
region where
\[
  \tau=\sigma\e^{-Az}z^\beta
\]
is no longer tiny, it may be better to hold $\tau$ fixed and sum over $k$ first:
\begin{equation}
  \widehat\Phi(z,\sigma)
  =\sum_{n\ge0}\frac{F_n(\tau)}{z^n}.
  \label{eq:transasymptotic-rearrangement}
\end{equation}
The coefficient functions $F_n$ resum infinitely many exponential sectors. This
is called a transasymptotic rearrangement. It can extend an approximation through
regions where the original sector-by-sector ordering breaks down.

\section{A practical hierarchy of methods}

\begin{longtable}{@{}p{0.23\textwidth}p{0.32\textwidth}p{0.37\textwidth}@{}}
\toprule
Method & What it adds & Typical use\\
\midrule
Poincar\'e asymptotics & finitely many algebraic terms & coarse large-parameter approximation\\
Optimal truncation & best use of a divergent series & exponential estimate of the remainder\\
Borel summation & analytic value for one Gevrey series & reconstructing a canonical sector\\
Resurgent transseries & coupled exponential sectors & Stokes jumps and ambiguity cancellation\\
Hyperasymptotics & re-expansion of remainders & exponentially improved numerical accuracy\\
Multisummability & several Gevrey levels & singular differential equations with multiple scales\\
Transasymptotics & sectors resummed into coefficient functions & transition regions and global matching\\
\bottomrule
\end{longtable}

\begin{summarybox}
Ordinary Borel summation handles one Gevrey-$1$ level. Borel--Leroy transforms,
$k$-summability, multisummability, and acceleration handle other growth scales.
Hyperasymptotics repeatedly expands optimally truncated remainders, while
transasymptotic rearrangement sums sectors in a new order when the exponential
parameter becomes order one.
\end{summarybox}


\chapter{A reusable analytic workflow}
\label{ch:analytic-workflow}

\section{Step 1: find the formal sectors}

Start from the equation, recurrence, or integral and identify all leading
balances. For an ODE, substitute
\[
  y\sim\e^{-A(z)}z^\beta
\]
into the linearized equation. For an integral, locate all relevant saddles and
endpoint singularities. Every independent action is a candidate exponential
scale.

\section{Step 2: derive coefficient recurrences}

For each sector, write
\[
  \Phi_k(z)=\e^{-A_k z}z^{\beta_k}
  \sum_{n\ge0}a_n^{(k)}z^{-n}
\]
or the appropriate nonconstant-action version. Substitute and solve in dominance
order. Keep the normalization $a_0^{(k)}$ explicit; changing it inversely changes
the corresponding Stokes constant.

\section{Step 3: inspect late terms}

Compute enough coefficients to test
\[
  a_n^{(0)}\sim
  K\frac{\Gamma(n+\gamma)}{A^{n+\gamma}}.
\]
Use ratios, Richardson extrapolation, and known symmetries. Complex-conjugate or
equidistant singularities may require separating subsequences or fitting a sum
of oscillatory contributions.

\section{Step 4: enter the Borel plane}

Divide by the appropriate gamma factor. Determine the radius of convergence,
then seek analytic continuation. With exact equations, the Borel transform may
satisfy an integral or differential equation. With finite data, use conformal
maps and Pad\'e approximants.

\section{Step 5: locate singular directions}

Each Borel singularity $A$ determines directions for which the Laplace ray hits
it. Each action also determines curves where $\Re(Az)=0$, across which
exponential dominance changes. Draw both geometries; they answer related but not
identical questions.

\section{Step 6: compute discontinuities}

For a pole, use residues. For logarithmic or algebraic branch points, compute the
discontinuity across a cut. Match that discontinuity to the Borel transform of a
neighboring sector. This determines a Stokes constant once sector normalizations
are fixed.

\section{Step 7: choose the physical or analytic parameter}

Boundary conditions, contour choices, regularity, reality, or quantization
conditions determine transseries parameters. Across a Stokes direction, transform
them according to the Stokes automorphism rather than keeping their numerical
coordinates artificially fixed.

\section{Step 8: validate at three levels}

A strong analysis performs three independent checks:
\begin{enumerate}
\item \textbf{Formal residual:} substituting a truncation leaves an error at the
predicted next transmonomial.
\item \textbf{Large-order check:} coefficients of one sector reproduce the
predicted action and early data of another.
\item \textbf{Numerical check:} lateral or median resummation agrees with direct
integration or an exact solution in overlapping domains.
\end{enumerate}
Agreement at all three levels makes normalization or sign errors much less
likely.

\begin{summarybox}
The analytic workflow runs from dominant balance to sectors, from sectors to
late terms, from late terms to Borel singularities, and from singularities to
Stokes transformations. Formal, large-order, and numerical residual checks form
a useful triangle of verification.
\end{summarybox}

% ============================================================================
\part{Where transseries live and where they are used}

\chapter{Hardy fields, \texorpdfstring{$H$}{H}-fields, and asymptotic differential algebra}
\label{ch:h-fields}

\section{Functions modulo eventual agreement}

Two real functions defined for all sufficiently large $x$ determine the same
\emph{germ at $+\infty$} if they eventually agree. A germ forgets behavior on
every bounded interval and retains only the tail. This is exactly the viewpoint
needed for asymptotics.

A \emph{Hardy field} is a field of such germs that is closed under
differentiation. Rational functions form a Hardy field. So do many fields built
from algebraic functions, exponentials, and logarithms. Every nonzero element of
a Hardy field has eventually constant sign; otherwise its infinitely many zeros
would conflict with field and derivative closure through Rolle-type arguments.
Consequently, Hardy fields carry a natural eventual order.

Transseries are formal rather than initially analytic, but their order and
derivative behave as though they were germs of eventually well-behaved
functions. $H$-fields axiomatize this behavior.

\section{The axioms of an \texorpdfstring{$H$}{H}-field}

Let $K$ be an ordered differential field with constant field
\[
  C=\{f\in K:f'=0\}.
\]
Write $f>C$ when $f$ is larger than every constant. Let
\[
  \mathcal O=\{f\in K:\abs f\le c\text{ for some }c\in C_{>0}\}
\]
be the convex hull of the constants, and let
\[
  \mathfrak o=\{f\in K:\abs f<c\text{ for every }c\in C_{>0}\}
\]
be its infinitesimal ideal.

\begin{definition}[$H$-field]
An ordered differential field $K$ is an $H$-field if:
\begin{enumerate}
\item whenever $f>C$, one has $f'>0$;
\item every bounded element has a unique constant part up to an infinitesimal:
\begin{equation}
  \mathcal O=C+\mathfrak o.
  \label{eq:h-field-bounded-decomposition}
\end{equation}
\end{enumerate}
\end{definition}

The first axiom says that a function infinitely larger than constants is
increasing. The second says that every bounded element looks like
``constant plus something tending to zero.'' Both are familiar facts for tame
function germs.

Suitable fields of logarithmic-exponential transseries are $H$-fields. Their
constant field is $\R$, their infinitesimals are precisely the series dominated
by $1$, and the leading monomial decides eventual sign.

\section{The natural valuation}

The dominance relation defines a valuation. Declare
\[
  f\asymp g
  \quad\Longleftrightarrow\quad
  f\precsimdom g\text{ and }g\precsimdom f.
\]
Two nonzero elements have the same value when they have the same asymptotic
magnitude up to constant factors. We may write
\[
  \val(f)>\val(g)
  \quad\Longleftrightarrow\quad f\precdom g,
\]
so larger valuation means smaller magnitude.

For a transseries, the value is determined by the leading monomial. The familiar
non-Archimedean law
\[
  \val(f+g)\ge\min\{\val(f),\val(g)\}
\]
is merely the statement that adding two quantities cannot produce something
larger than the larger summand, though cancellation can make it smaller.

\section{The asymptotic couple}

Differentiation interacts with valuation through the logarithmic derivative.
If $\gamma=\val(f)$, define schematically
\begin{equation}
  \psi(\gamma)=\val\left(\frac{f'}{f}\right).
  \label{eq:asymptotic-couple-psi}
\end{equation}
Under the usual hypotheses, this depends only on $\gamma$, not on the chosen
$f$. The ordered abelian value group $\Gamma$ together with
$\psi:\Gamma^{\ne0}\to\Gamma$ is an \emph{asymptotic couple}.

For simple monomials:
\begin{align*}
  (x^a)^\dagger&=\frac{a}{x},\\
  (\e^{x^2})^\dagger&=2x,\\
  ((\log x)^b)^\dagger&=\frac{b}{x\log x}.
\end{align*}
The map $\psi$ records these derivative scales abstractly. It is a compact
combinatorial skeleton of asymptotic differentiation.

\section{Liouville closure}

A differential field is \emph{Liouville closed} if it contains antiderivatives
and exponential integrals in the differential-algebraic sense: for every $f$,
there are $g\ne0$ and $h$ with
\[
  \frac{g'}g=f,
  \qquad h'=f,
\]
and the underlying ordered field is real closed. A sufficiently rich standard
field of transseries has these closure properties. This formalizes our earlier
claims that one can integrate transseries and solve first-order linear
differential equations inside the field.

The word ``Liouville'' here is differential-algebraic and should not be confused
with Liouville's theorem in complex analysis.

\section{Newtonianity}

For a polynomial equation, Newton polygons compare powers and valuations to find
dominant balances. Differential polynomials require a differential analogue.
An $H$-field is called \emph{newtonian}, roughly, when every differential
polynomial of Newton degree one has a zero in the valuation ring. This is an
asymptotic version of algebraic henselianity.

The technical definition is beyond our scope, but the idea is concrete:
formal Newton iteration should not fail merely because the unknown occurs with
derivatives. Transseries possess enough monomials and closure operations to
solve broad classes of asymptotic differential equations recursively.

\section{Model theory}

Aschenbrenner, van den Dries, and van der Hoeven developed a model theory of
$H$-fields centered on transseries \cite{ADHBook}. In an appropriate language,
properties such as Liouville closure, newtonianity, and the absence of certain
asymptotic gaps give an axiomatic description of the first-order behavior of the
standard transseries field. This connects asymptotic expansion to questions
normally asked about algebraically closed fields or real closed fields:
Which statements are consequences of the axioms? Which extensions are
existentially closed? Which decision procedures are possible?

The surprise is that an object built from expressions such as
$\e^{x+\e^{-x}}$ has a coherent global algebraic theory, not merely a collection
of ad hoc expansion rules.

\section{From formal series to actual germs}

A formal transseries may be realized by an analytic germ through summation. For
special classes, this realization respects order, algebra, composition, and
differentiation. One seeks an injective map
\[
  \mathfrak s:\T_{\mathrm{summable}}\longrightarrow\mathcal H
\]
into a Hardy field or field of analytic germs such that
\begin{align*}
  \mathfrak s(T+U)&=\mathfrak s(T)+\mathfrak s(U),\\
  \mathfrak s(TU)&=\mathfrak s(T)\mathfrak s(U),\\
  \mathfrak s(T')&=(\mathfrak sT)',\\
  \mathfrak s(T\compose U)&=(\mathfrak sT)\compose(\mathfrak sU)
\end{align*}
whenever the operations are defined. Results of this kind connect formal
logarithmic-exponential series with analyzable germs
\cite{vanDenDriesMacintyreMarker1997,vanDenDriesMacintyreMarker2001,CostinBook}.

Not every formal transseries has a unique analytic realization without a chosen
summation structure. Stokes data is precisely part of what must be supplied.

\begin{summarybox}
Hardy fields describe germs of tame functions at infinity; $H$-fields abstract
their ordered differential behavior. Transseries form a central $H$-field whose
valuation comes from dominance and whose derivative is encoded by an asymptotic
couple. Differential-algebraic closure and model theory explain why transseries
solve asymptotic equations so systematically.
\end{summarybox}


\chapter{Surreal numbers and transseries}
\label{ch:surreals}

\section{A family resemblance}

Conway's surreal numbers form a proper class $\mathbf{No}$ containing the real
numbers, all ordinal numbers, infinitesimals, and infinitely large quantities.
Every surreal has a normal form
\begin{equation}
  a=\sum_{i<\alpha} r_i\omega^{a_i},
  \label{eq:conway-normal-form}
\end{equation}
where $\alpha$ is an ordinal, $r_i\in\R\setminus\{0\}$, and the exponents $a_i$
strictly decrease. This looks remarkably like a Hahn series whose monomial group
is generated recursively by powers of the simplest infinite surreal $\omega$.

Compare it with a transseries normal form
\[
  T=\sum_{\m\in S}c_{\m}\m.
\]
Both use ordered supports, leading terms, and transfinite sums. The main
difference is emphasis: surreal numbers begin as a universal ordered field,
while transseries begin as a universal language for asymptotic functions and
therefore carry composition and differentiation as essential operations.

\section{Exponentiation on the surreals}

Gonshor defined an exponential map
\[
  \exp:\mathbf{No}\to\mathbf{No}_{>0}
\]
extending the real exponential and compatible with the ordered-field structure.
Its inverse is a logarithm. This allows surreal counterparts of iterated
exponentials, iterated logarithms, and logarithmic-exponential monomials.

One may therefore send the transseries variable $x$ to an infinite surreal such
as $\omega$ and interpret many transseries as surreal numbers. The leading
transmonomial becomes the leading Conway monomial.

\section{Embedding transseries}

There are structure-preserving embeddings of major transseries fields into the
surreals, compatible with order and exponential structure. In suitable
frameworks, differential structure can also be made compatible
\cite{BerarducciMantovaSurreal,ADHBook}. Symbolically,
\[
  \T\hookrightarrow\mathbf{No},
  \qquad x\longmapsto\omega.
\]
For example,
\[
  x+\log x+\e^{-x}
  \quad\longmapsto\quad
  \omega+\log\omega+\e^{-\omega}.
\]
The infinitesimal $\e^{-\omega}$ is smaller than every power $\omega^{-n}$,
exactly mirroring the transseries dominance relation.

\section{Derivations on the surreals}

A derivation on $\mathbf{No}$ should satisfy the ordinary algebraic rules,
annihilate real constants, respect exponentiation,
\[
  D(\e^a)=\e^aD(a),
\]
and behave well with summable families. Constructing such a derivation on a
proper class with enormous supports is nontrivial. Berarducci and Mantova built
surreal derivations with strong closure properties, making the analogy with
transseries precise \cite{BerarducciMantovaSurreal}.

There is no requirement that an abstract infinite number possess a derivative.
The derivative is additional structure selected to make $\omega$ play the role
of a variable tending to infinity.

\section{Why the surreals are larger}

The ordinary logarithmic-exponential transseries field permits only finite
exponential height and logarithmic depth for each individual element. The
surreals contain vastly longer transfinite patterns. They also contain many
objects with no immediate interpretation as ordinary asymptotic germs.

Conversely, function composition $T\compose U$ has no automatic counterpart in
an arbitrary ordered exponential field: it depends on distinguishing a variable
from a value. Thus neither slogan ``transseries are just surreal numbers'' nor
``surreal numbers are just transseries'' is accurate. There is a deep embedding
and shared Hahn-series geometry, but the intended structures differ.

\section{Germs, transseries, and surreals}

A useful three-way picture is
\[
  \boxed{\text{analytic germs}}
  \quad\longleftrightarrow\quad
  \boxed{\text{formal transseries}}
  \quad\longrightarrow\quad
  \boxed{\text{surreal values at }\omega}.
\]
The first arrow requires summation or asymptotic expansion and is generally not
one-to-one without analytic restrictions. The second is algebraic and can be
made into a rich embedding. Research relating germs, transseries, and surreal
numbers asks how much of this triangle can simultaneously preserve order,
exponential, composition, and differentiation
\cite{BerarducciMantovaGerms,NumbersGermsTransseries}.

\begin{summarybox}
Surreal normal forms and transseries are both Hahn-like sums of ordered
monomials. Exponentiation and suitable derivations permit embeddings of major
transseries fields into the surreals, often by sending $x$ to $\omega$. The
surreals are much larger, while transseries carry a more explicit function-like
structure through composition and asymptotic differentiation.
\end{summarybox}


\chapter{Dynamical systems, normal forms, and iteration}
\label{ch:dynamics}

\section{Normal forms remove terms scale by scale}

Near a fixed point, a differential equation or iteration can often be simplified
by a change of coordinates. Ordinary power series remove nonresonant polynomial
terms. Exponentially small or logarithmic effects remain after every algebraic
order and require transseries.

Suppose that, after a coordinate change near infinity,
\[
  x_{n+1}=F(x_n),
  \qquad F(x)=x+1+\frac{a}{x}+\bigO(x^{-2}).
\]
An Abel coordinate $V$ satisfies
\begin{equation}
  V(F(x))=V(x)+1.
  \label{eq:abel-equation-dynamics}
\end{equation}
One expects
\[
  V(x)=x-a\log x+\frac{b_1}{x}+\cdots,
\]
together with exponentially small corrections that encode global analytic
data. Once $V$ is known, fractional iterates are
\[
  F^{\compose t}(x)=V^{-1}(V(x)+t).
\]
This is the construction introduced formally in
Chapter~\ref{ch:composition}; resurgence describes how different analytic
normalizations of $V$ are related.

\section{Dulac series}

Near a singular point of a planar analytic vector field, transition maps can
involve powers and logarithms such as
\begin{equation}
  \sum_{\alpha,k}c_{\alpha,k}x^\alpha(\log x)^k.
  \label{eq:dulac-series}
\end{equation}
These are called Dulac series or generalized Dulac expansions. Iterated
transition maps can require exponential refinements as well. Such expansions
play an important role in the study of limit cycles and local return maps.

The support conditions are again essential. Without a well-ordered dominance
structure, composition of two return-map expansions could assign infinitely many
contributions to one coefficient.

\section{Quasianalyticity}

A class of functions is quasianalytic if a function is determined by its full
asymptotic expansion: no nonzero flat function belongs to the class. Ordinary
smooth functions are not quasianalytic because $\e^{-1/x^2}$ is flat at $0$.
But carefully restricted classes of Dulac or analyzable functions can be
quasianalytic.

Quasianalyticity matters in dynamics because a return map with zero formal
expansion should then be the identity, ruling out hidden cycles or degeneracies.
Transseries plus a summation theorem can provide exactly the additional
structure needed to turn a formal expansion into a uniqueness statement.

\section{Separatrix splitting}

In nearly integrable systems, stable and unstable manifolds may coincide to all
orders in a perturbation parameter yet differ by
\[
  \Delta(\eps)\sim C\eps^\beta\e^{-A/\eps}.
\]
Every coefficient of the ordinary power series for $\Delta$ is zero. The actual
splitting is beyond all algebraic orders and can create chaotic intersections.
Transseries are therefore not a decorative refinement: the leading dynamical
effect may live entirely in an exponential sector.

\section{Difference equations}

Difference equations naturally combine composition and asymptotics. A simple
example is
\begin{equation}
  y(x+1)-y(x)=\frac1x.
  \label{eq:difference-harmonic}
\end{equation}
A formal solution begins
\[
  y(x)\sim\log x-\frac1{2x}-\frac1{12x^2}
       +\frac1{120x^4}-\cdots,
\]
a version of the asymptotic expansion of the digamma function. Bernoulli
numbers make the series divergent, and exponentially small periodic information
is not captured by the inverse-power series alone.

Difference equations also lead beyond finite logarithmic depth. Solving
\[
  F(x+1)=\exp(F(x))
\]
or constructing regular iteration of the exponential motivates
transexponential and hyperseries extensions
\cite{SchmelingThesis,PadgettTransexponential,vanDenDriesHoevenKaplan}.

\section{Parabolic germs and \'Ecaille--Voronin invariants}

A holomorphic germ tangent to the identity,
\[
  f(z)=z+az^{k+1}+\cdots,
\]
has formal normal forms that do not completely determine its analytic conjugacy
class. Sectorial Fatou coordinates differ on overlaps by exponentially small
maps. These differences, organized into \'Ecaille--Voronin invariants, are a
canonical example of analytic data invisible to ordinary formal power series.

The moral is the same as for the Euler equation: formal normalization leaves
flat ambiguities, and Stokes-type transition data classifies the analytic
objects.

\begin{summarybox}
Transseries enter dynamics through normal forms, fractional iteration, Dulac
maps, separatrix splitting, and difference equations. They record effects that
vanish to every power of a perturbation parameter and organize sectorial
coordinate changes whose mismatches are analytic invariants.
\end{summarybox}


\chapter{Saddles, instantons, and mathematical physics}
\label{ch:physics}

\section{The universal saddle expansion}

Consider an integral
\begin{equation}
  I(g)=\int_{\mathcal C}\e^{-S(x)/g}a(x)\dd x,
  \qquad g\to0.
  \label{eq:semiclassical-integral}
\end{equation}
Critical points $x_j$ of $S$ satisfy $S'(x_j)=0$. Expanding near one saddle
gives
\[
  I_j(g)\sim
  \e^{-S(x_j)/g}g^{\beta_j}
  \sum_{n=0}^{\infty}a_{j,n}g^n.
\]
The full integral is a linear combination
\begin{equation}
  I(g)\sim
  \sum_j\sigma_j
  \e^{-S_j/g}g^{\beta_j}\phihat_j(g).
  \label{eq:saddle-transseries}
\end{equation}
This is a transseries with actions $S_j$ and parameters determined by how the
contour decomposes into steepest-descent cycles.

The Airy integral of Chapter~\ref{ch:airy} is the simplest nontrivial example.
In more dimensions, saddles can be isolated critical points, critical manifolds,
complex solutions, or singular endpoints.

\section{Instantons}

In quantum mechanics or field theory, a Euclidean classical solution connecting
distinct vacua is called an instanton. Its action $A$ produces a factor
\[
  \e^{-A/g}
  \quad\text{or}\quad
  \e^{-A/\hbar}.
\]
Fluctuations around the instanton produce a divergent power series. Configurations
with $k$ instantons produce sectors $\e^{-kA/g}$, often with powers and logarithms
of $g$ arising from zero modes and interactions.

The term \emph{instanton sector} is therefore a physical realization of the
formal sector index $k$ in \eqref{eq:one-parameter-resurgent-transseries}.

\section{Ambiguity cancellation}

In many problems, the perturbative Borel transform has a singularity on the
physical positive axis. Its two lateral sums differ by an imaginary quantity of
order $\e^{-A/g}$. A nonperturbative saddle sector has its own ambiguous
continuation, with the opposite imaginary jump. The sum is real and unambiguous
once all required sectors are included.

Symbolically,
\[
  \operatorname{Im}\mathcal S_{0^\pm}\phihat_0
  +\operatorname{Im}\mathcal S_{0^\pm}
     \left(\e^{-A/g}\phihat_1\right)=0.
\]
This cancellation is sometimes called the Bogomolny--Zinn-Justin mechanism in
quantum-mechanical examples. Resurgence explains it as one Stokes automorphism
acting consistently on the whole transseries.

\section{Perturbation theory predicts nonperturbative physics}

Large-order behavior of perturbative coefficients can reveal:
\begin{enumerate}
\item the smallest instanton or complex-saddle action;
\item the power-law prefactor from fluctuation zero modes;
\item Stokes constants;
\item low-order coefficients around a nonperturbative saddle.
\end{enumerate}
Conversely, a semiclassical saddle calculation predicts the divergence pattern
of high perturbative orders. This two-way relation is one of the most useful
consequences of resurgence.

\section{Renormalons and other singularities}

Not every Borel singularity in quantum field theory is associated with a finite-
action classical solution. Renormalization-group effects can produce
\emph{renormalon} singularities. Their status and interpretation depend on the
theory, observable, and regularization. A complete transseries may therefore
need sectors motivated by Borel-plane analysis rather than by obvious real
instantons.

This is a useful warning against identifying ``nonperturbative'' exclusively
with ``real classical solution.'' Complex saddles, singular configurations, and
collective effects can all contribute.

\section{Exact WKB}

For a Schr\"odinger-type equation
\begin{equation}
  -\hbar^2\psi''(x)+Q(x)\psi(x)=0,
  \label{eq:exact-wkb-equation}
\end{equation}
set
\[
  \psi(x)=\exp\!\left(\frac1\hbar\int^x S(t,\hbar)\dd t\right).
\]
The Riccati equation
\[
  S^2+\hbar S'=Q
\]
produces a divergent WKB series. Its Borel sums change across Stokes curves
emanating from turning points. Period integrals
\[
  A_\gamma=\oint_\gamma\sqrt{Q(x)}\dd x
\]
act as instanton actions. Exact quantization conditions combine perturbative
periods and exponentially small tunneling sectors into resurgent relations.

\section{Many modern arenas}

Resurgent transseries appear in matrix models, topological strings, gauge
theories, hydrodynamic gradient expansions, nonlinear wave equations, and
semiclassical spectral problems. The details differ, but the reusable pattern is
stable:
\[
  \text{saddles/actions}
  \longleftrightarrow
  \text{Borel singularities}
  \longleftrightarrow
  \text{large-order growth}
  \longleftrightarrow
  \text{Stokes data}.
\]
The primer by Aniceto, Ba\c{s}ar, and Schiappa develops this viewpoint across
many physical examples \cite{AnicetoBasarSchiappa}; Dorigoni gives an accessible
entry to alien calculus \cite{DorigoniIntro}.

\begin{warningbox}
In an infinite-dimensional path integral, the ``integral'' may itself be formal,
and existence of all required saddles or resummations can be difficult. The
finite-dimensional saddle picture is a guide, not an automatic theorem for every
quantum field theory.
\end{warningbox}

\begin{summarybox}
Small-coupling integrals and semiclassical problems decompose into saddle sectors
$\e^{-S_j/g}$ times divergent fluctuation series. Instantons are particular
saddles; large perturbative orders encode their data. Stokes ambiguities cancel
only after all coupled sectors are transformed and summed consistently.
\end{summarybox}


\chapter{Computer algebra and numerical exploration}
\label{ch:computer-algebra}

\section{Why representation matters}

A computer cannot store an infinite support. It stores a finite truncation plus
a description of the admissible monomial grid. A practical transseries object
might contain:
\begin{enumerate}
\item a ratio set $\bm\mu=(\mu_1,\ldots,\mu_r)$;
\item an exponent lower bound $\bm k_0\in\Z^r$;
\item a finite map from exponent vectors to coefficients;
\item a truncation cut expressed in dominance order;
\item recursive expression nodes for exponentials, logarithms, and composition.
\end{enumerate}
The mathematical support theorem guarantees that omitted terms cannot contribute
above the cut in infinitely many hidden ways.

\section{Canonicalization}

Different generator expressions can denote the same monomial:
\[
  (\e^{-x})^2=\e^{-2x},
  \qquad
  x^{-1}\e^{-\log x}=x^{-2}.
\]
A computer algebra system must normalize logarithmic exponents, combine
exponentials, and detect relations among ratio generators. Otherwise equal terms
will fail to collect and leading-term comparisons will be wrong.

One robust approach represents a transmonomial recursively by a normalized
exponent $L$ in $\e^L$, with $L$ itself a purely large transseries, together with
explicit powers from lower logarithmic levels. Comparison then descends through
the recursive structure.

\section{Truncated multiplication}

Given finite maps for $T$ and $U$, naive multiplication forms every pair of
stored terms. Better algorithms exploit order:
\begin{enumerate}
\item sort terms by decreasing dominance;
\item use a priority queue of pairwise products;
\item emit the largest unseen product;
\item merge duplicates;
\item stop as soon as the next possible product lies below the truncation cut.
\end{enumerate}
This is analogous to multiplying sparse multivariate series with an admissible
monomial order.

\section{Recursive expansion engines}

Operations such as inverse, logarithm, exponential, and composition reduce to
universal templates after leading-term factorization:
\begin{align*}
  (1+U)^{-1}&=\sum_{n\ge0}(-U)^n,\\
  \log(1+U)&=\sum_{n\ge1}\frac{(-1)^{n+1}}nU^n,\\
  \e^U&=\sum_{n\ge0}\frac{U^n}{n!},\\
  T(x+U)&=\sum_{n\ge0}\frac{T^{(n)}(x)}{n!}U^n.
\end{align*}
The engine needs a proof that $U\precdom1$ in the relevant asymptotic scale and
a termination test showing that $U^n$ eventually lies below the requested cut.

\section{Equality is harder than expansion}

For a normalized finite truncation, equality is easy. Exact zero testing for rich
transseries expressions is subtler because cancellation may occur after deep
functional identities or normalization across exponential levels. Differential
algebra and model theory supply decision procedures for important classes, while
recent work develops explicit algorithms for larger logarithmic-exponential
transseries fields \cite{ChenFangHoeven2026}.

No symbolic system should silently infer exact zero merely because every term
computed so far vanished. That proves only that the expression lies below the
current cut.

\section{Numerical Borel analysis}

A practical numerical pipeline is:
\begin{enumerate}
\item generate $N$ coefficients at high precision;
\item divide by the expected gamma growth;
\item build Pad\'e or conformal-Borel approximants;
\item locate stable singularity patterns;
\item integrate laterally with adaptive quadrature;
\item compare the jump with the predicted neighboring sector;
\item fit transseries parameters to boundary or initial data.
\end{enumerate}

Conformal mapping can greatly improve Pad\'e resolution. If the nearest Borel
singularity is known to be at $A>0$ with a cut $[A,\infty)$, a map from the cut
plane to the unit disk places the entire first sheet inside a convergent Taylor
region. The precise map should reflect the known singularity geometry rather
than being chosen mechanically.

\section{Richardson acceleration}

Suppose a sequence has
\[
  q_n=q+\frac{c_1}{n}+\frac{c_2}{n^2}+\cdots.
\]
Richardson transforms eliminate successive inverse powers of $n$, accelerating
convergence to $q$. Large-order tests routinely use this after dividing out the
leading $\Gamma(n+\beta)A^{-n}$ growth.

A simple first transform is
\[
  R_n=(n+1)q_{n+1}-nq_n=q+\bigO(n^{-2}).
\]
Repeated or generalized versions can expose Stokes constants hidden behind slow
$1/n$ corrections.

\section{Verification by exact arithmetic when possible}

Many coefficient recurrences use rational numbers and symbolic parameters.
Compute them exactly before converting to floating point. Exact arithmetic can
reveal a missed cancellation, a hidden factorization, or an integer sequence.
Use high-precision numerical arithmetic only for analytic continuation and
resummation steps that genuinely require it.

\section{A minimal data contract}

A reliable transseries library should distinguish:
\begin{description}[style=nextline]
\item[Formal equality]
proved after complete normalization;
\item[Equality modulo a cut]
agreement of all terms above a specified monomial;
\item[Asymptotic equivalence]
agreement of leading terms;
\item[Numerical agreement]
closeness after a specified summation and parameter choice.
\end{description}
Conflating these notions is a common source of plausible but incorrect symbolic
output.

\begin{summarybox}
Computer algebra stores finite truncations on certified monomial grids.
Canonicalization, dominance-aware multiplication, and cut-aware recursive
expansion are essential. Exact equality is deeper than matching a truncation.
For divergent sectors, high-precision Borel--Pad\'e and large-order analysis
connect symbolic coefficient generation to numerical resummation.
\end{summarybox}


\chapter{Extensions, boundaries, and several meanings of ``transseries''}
\label{ch:boundaries}

\section{Two traditions}

The word \emph{transseries} is used in two overlapping traditions.

\begin{enumerate}
\item In ordered algebra and model theory, a transseries is a Hahn-like formal
sum of logarithmic-exponential monomials, equipped with order, derivative,
exponential, logarithm, and composition.
\item In asymptotic analysis and physics, a transseries is often written as a
sum of exponential sectors, each carrying a divergent power series and one or
more free parameters.
\end{enumerate}

The first tradition emphasizes legal infinite algebra and universal closure. The
second emphasizes Borel summation, Stokes data, and analytic reconstruction.
They meet when the sector prefactors and fluctuations are embedded into a common
ordered monomial field, but neither vocabulary should be forced into the other's
narrowest notation.

\section{Finite height and depth}

An ordinary logarithmic-exponential transseries has finite exponential height
and finite logarithmic depth for each element. It may contain
\[
  \exp\!\bigl(\exp(x)+x^2\bigr)
\]
but not automatically an object requiring infinitely many nested exponentials
in one element. Likewise, each transseries uses only finitely many iterated
logarithms as basic scales.

This elementwise finiteness still allows a field containing objects of every
finite height and depth. It is analogous to the set of all finite polynomials:
there is no uniform degree bound, but each polynomial has finite degree.

\section{Transexponentials and hyperseries}

Functional equations involving iteration can demand new scales. A
transexponential grows faster than every finite iterate of the ordinary
exponential while satisfying an iteration relation adapted to exponentiation.
Logarithmic hyperseries and related constructions systematically adjoin such
transfinite levels
\cite{SchmelingThesis,vanDenDriesHoevenKaplan,PadgettTransexponential}.

These extensions preserve the central design principles:
ordered monomials, well-based supports, recursive closure, and a derivative
compatible with asymptotic growth. What changes is the ordinal length of the
scale hierarchy.

\section{Oscillation}

Real transseries ordered by eventual magnitude do not directly encode persistent
oscillation such as $\sin x$, because it has no eventual sign and neither tends
to zero nor infinity. Complex exponential monomials
\[
  \e^{i\phi(x)}
\]
can be adjoined, but their magnitude may be $1$, so phase order must be tracked
separately from growth order. WKB and Fourier-type transseries therefore use a
richer combination of exponential actions and oscillatory phases.

For $z$ in a complex sector, $\e^{-Az}$ can itself switch between decay, growth,
and oscillation as the argument rotates. Complex transseries are naturally
sectorial.

\section{Branching and logarithms}

The complex logarithm is multivalued. A formal expression involving $\log z$
or $z^\alpha$ must be tied to a branch and a sector. Analytic continuation around
the origin can add $2\pi i$ to a logarithm and change sector parameters. Formal
algebra alone does not choose the branch.

Similarly, a Borel transform can live on a ramified Riemann surface rather than
the punctured plane. Alien calculus records continuation among sheets and around
singularities.

\section{Not every function has a useful transseries}

Some functions have dense oscillations, natural boundaries, or asymptotic
behavior not captured by a selected monomial scale. A function can possess a
Poincar\'e expansion yet fail to be resurgent. It can also require generalized
monomials not present in the basic field.

The correct conclusion is not that transseries failed, but that every formal
language has a domain. One should ask:
\begin{enumerate}
\item Which scales are generated by the problem?
\item Is the support well based?
\item What analytic continuation is available?
\item Which summation method matches the coefficient growth?
\end{enumerate}

\section{Formal existence versus analytic existence}

A formal transseries solution can often be constructed recursively even when no
analytic solution has yet been proved. Turning it into a function requires
additional theorems: summability, fixed-point estimates, sectorial analyticity,
or integral representations. Conversely, an analytic solution may exist but
have no expansion in a chosen transseries class.

Keep the logical levels separate:
\[
  \text{formal construction}
  \quad\ne\quad
  \text{summability}
  \quad\ne\quad
  \text{global analytic continuation}.
\]
A rigorous argument states which level has been established.

\section{A compact map of the subject}

\begin{center}
\begin{tikzpicture}[>=Latex,node distance=12mm and 17mm,
  every node/.style={draw,rounded corners,align=center,inner sep=5pt}]
  \node[fill=softblue] (scales) {asymptotic\\scales};
  \node[fill=softblue,right=of scales] (hahn) {Hahn fields and\\well-based support};
  \node[fill=softgreen,right=of hahn] (calc) {exp, log, derivative,\\composition};
  \node[fill=softorange,below=of hahn] (formal) {formal solutions\\and sectors};
  \node[fill=softorange,left=of formal] (late) {factorial divergence\\and late terms};
  \node[fill=softred,right=of formal] (borel) {Borel singularities\\and Stokes data};
  \node[fill=softgreen,below=of formal] (analytic) {summed analytic\\solutions};
  \draw[->] (scales)--(hahn);
  \draw[->] (hahn)--(calc);
  \draw[->] (calc)--(formal);
  \draw[<->] (formal)--(late);
  \draw[<->] (formal)--(borel);
  \draw[->] (late)--(borel);
  \draw[->] (borel)--(analytic);
  \draw[->] (formal)--(analytic);
\end{tikzpicture}
\end{center}

\begin{summarybox}
``Transseries'' names both a formal ordered field and a sectorial language for
resurgent asymptotics. Basic logarithmic-exponential fields have elementwise
finite height and depth; hyperseries extend them. Oscillation, branching, and
nonresurgent singularity patterns require extra structure. Formal existence,
summability, and global analytic existence are distinct claims.
\end{summarybox}

\chapter{A brief historical and bibliographic guide}
\label{ch:history}

\section{Long before the name}

Asymptotic expansions are older than transseries. Euler freely manipulated
divergent series in the eighteenth century. Poincar\'e clarified asymptotic
expansions in the nineteenth century. Stokes discovered that asymptotic
coefficients of special functions change their effective form across rays in
the complex plane. Hardy's \emph{Orders of Infinity} organized logarithmic and
exponential growth scales \cite{HardyOrders}.

These developments supplied the analytic motivation: ordinary power series were
too small a language for many limiting problems.

\section{Generalized power series}

Hahn's 1907 construction placed formal sums over ordered abelian groups on a
firm algebraic foundation \cite{Hahn1907}. Neumann's lemma later supplied a
powerful support theorem ensuring that products and geometric-series inverses
remain well defined \cite{Neumann1949}. Modern transseries inherit this Hahn-field
backbone.

The lesson of this algebraic line is that ``infinite formal expression'' is not
enough. One must choose an ordered monomial group and a support condition that
makes every coefficient a finite sum.

\section{Dahn--G\"oring, \'Ecaille, and independent origins}

Logarithmic-exponential transseries were introduced in closely related forms by
Dahn and G\"oring in model-theoretic work and by Jean \'Ecaille in his theory of
analyzable functions and resurgence \cite{DahnGoring,EcalleAnalyzable}. The two
traditions emphasized different questions:
\begin{itemize}
\item ordered fields, exponentiation, and logical structure on one side;
\item divergent asymptotics, alien calculus, and Stokes phenomena on the other.
\end{itemize}
The modern subject benefits from both.

\section{Computer algebra and van der Hoeven's synthesis}

Joris van der Hoeven developed transseries as a systematic computational
language, including algorithms for arithmetic, composition, differentiation,
integration, differential equations, and generalized summation
\cite{vanDerHoevenBook,vanDerHoevenOperators}. His monograph is a central source
for the formal and algorithmic viewpoint. Edgar's survey gives a shorter,
non-specialist construction with particular attention to grids and witnesses
\cite{EdgarBeginners}.

\section{Model theory and \texorpdfstring{$H$}{H}-fields}

The papers of van den Dries, Macintyre, and Marker connected logarithmic-
exponential series with o-minimal structures and Hardy fields
\cite{vanDenDriesMacintyreMarker1997,vanDenDriesMacintyreMarker2001}. The later
monograph of Aschenbrenner, van den Dries, and van der Hoeven developed the broad
asymptotic differential algebra and model theory summarized in
Chapter~\ref{ch:h-fields} \cite{ADHBook}.

\section{Resurgence in analysis and physics}

Costin developed rigorous generalized Borel summation and transseries for broad
classes of nonlinear differential systems
\cite{CostinTopological,CostinBook,CostinRankOne}. Sauzin and Dorigoni provide
entry points to \'Ecaille's resurgent framework
\cite{SauzinIntro,DorigoniIntro}. The extensive primer of Aniceto, Ba\c{s}ar,
and Schiappa explains large-order relations, nonlinear resonance, alien calculus,
and physical applications \cite{AnicetoBasarSchiappa}.

\section{Surreals and larger scales}

Berarducci and Mantova constructed compatible transserial derivations on the
surreal numbers and interpreted transseries as germs of surreal functions
\cite{BerarducciMantovaSurreal,BerarducciMantovaGerms}. Logarithmic hyperseries
and transexponential fields extend the scale beyond ordinary finite-height
logarithmic-exponential series
\cite{vanDenDriesHoevenKaplan,PadgettTransexponential}.

\section{Suggested next reading}

After this tutorial, choose according to purpose:
\begin{description}[style=nextline]
\item[Formal calculations]
Edgar's \emph{Transseries for Beginners}, then van der Hoeven's
\emph{Transseries and Real Differential Algebra}.
\item[Ordered differential algebra]
Aschenbrenner--van den Dries--van der Hoeven,
\emph{Asymptotic Differential Algebra and Model Theory of Transseries}.
\item[Borel summation and rigorous ODEs]
Costin's \emph{Asymptotics and Borel Summability}.
\item[Alien calculus]
Dorigoni's introduction, then Sauzin's notes and \'Ecaille's original works.
\item[Physics and detailed examples]
Aniceto--Ba\c{s}ar--Schiappa's primer.
\item[Surreal connections]
Berarducci--Mantova and Aschenbrenner's essay on numbers, germs, and transseries.
\end{description}

\begin{summarybox}
Transseries grew from several lines: classical asymptotics and Stokes phenomena,
Hahn's generalized series, model theory, \'Ecaille's resurgence, and computer
algebra. The modern subject is not owned by one notation. Its unity lies in
ordered asymptotic scales, legal infinite algebra, and the recovery of
exponentially small analytic information.
\end{summarybox}

% ============================================================================
\part{Problems, solutions, and reference material}

\chapter{Guided exercises}
\label{ch:exercises}

The exercises are grouped roughly by topic, not strictly by difficulty. Most
require only algebra and calculus. A star marks a problem that introduces a
slightly more advanced idea. Complete solutions begin in
Chapter~\ref{ch:solutions}.

\section{Asymptotic scales and first transseries}

\begin{exercise}[Ordering a mixed scale]
\label{ex:order-mixed}
Arrange the following monomials from largest to smallest as $x\to+\infty$:
\[
  \e^{\e^x},\quad \e^{x^2},\quad x^x,\quad
  x^{-100}\e^{3x},\quad x^{100},\quad(\log x)^7,
  \quad1,\quad x^{-2},\quad\e^{-\sqrt x},\quad\e^{-x}.
\]
Justify every neighboring comparison by taking a quotient or logarithm.
\end{exercise}

\begin{exercise}[A rational expansion]
\label{ex:rational-expansion}
Expand
\[
  \frac{x^2+1}{x^2-x}
\]
through the term of order $x^{-4}$, and verify the result by multiplying by the
denominator.
\end{exercise}

\begin{exercise}[Beyond every power]
\label{ex:beyond-power}
Prove that
\[
  \e^{-\sqrt x}\precdom x^{-N}
\]
for every fixed $N>0$. Then compare $\e^{-x}$ and $\e^{-\sqrt x}$.
\end{exercise}

\begin{exercise}[Flat ambiguity]
\label{ex:flat-ambiguity}
Suppose
\[
  f(x)\asympseries\sum_{n\ge0}a_nx^{-n}
  \qquad(x\to+\infty).
\]
Show directly from the definition that $f(x)+C\e^{-x}$ has the same asymptotic
power series for every constant $C$. Why does this not imply that all the
functions are equal?
\end{exercise}

\begin{exercise}[Additive and multiplicative normal form]
\label{ex:normal-form}
Let
\[
  T=3\e^{x^2+x}-5x^{100}\e^{x^2}+7x^{-2}.
\]
Find its dominant monomial, leading coefficient, additive decomposition into
giant/constant/small parts, and multiplicative form $T=a\m(1+U)$ with
$U\precdom1$.
\end{exercise}

\section{Supports and Hahn-series algebra}

\begin{exercise}[Which supports are legal?]
\label{ex:legal-supports}
For each formal sum below, decide whether its support is well based in the
large-$x$ dominance order used in this book:
\[
  \sum_{n\ge0}x^{-n},\qquad
  \sum_{n\ge0}x^n,\qquad
  \sum_{n\ge1}x^{-1/n},\qquad
  \sum_{n\ge0}\e^{-nx}.
\]
Explain the obstruction in each illegal case.
\end{exercise}

\begin{exercise}[Why a product coefficient is finite]
\label{ex:finite-product}
Let
\[
  A=\sum_{n\ge0}x^{-n},
  \qquad B=\sum_{n\ge0}(n+1)x^{-n}.
\]
Compute the coefficient of $x^{-k}$ in $AB$. Why are only finitely many pairs of
terms involved for each fixed $k$? Sum the answer in closed form as an ordinary
rational function and check it.
\end{exercise}

\begin{exercise}[A two-scale inverse]
\label{ex:two-scale-inverse}
Expand
\[
  \frac1{1+x^{-1}+\e^{-x}}
\]
through $x^{-3}$ in the purely algebraic sector and through
$x^{-2}\e^{-x}$ in the sector linear in $\e^{-x}$. Ignore terms containing
$\e^{-2x}$. Explain why this ``rectangular'' truncation is not an initial segment
of the full dominance order.
\end{exercise}

\begin{exercise}[Logarithm of a perturbed monomial]
\label{ex:log-perturbed}
Find the expansion through $x^{-3}$ of
\[
  \log\!\left[x^3\left(1+\frac2x-\frac1{x^2}\right)\right].
\]
\end{exercise}

\begin{exercise}[Asymptotic integration]
\label{ex:integrate-exp-x2}
Seek a formal antiderivative of $\e^{x^2}$ in the form
\[
  I(x)=\frac{\e^{x^2}}{2x}
  \sum_{n=0}^{\infty}c_nx^{-2n}.
\]
Derive a recurrence for $c_n$ and compute $c_0,c_1,c_2,c_3$. Verify the result by
differentiation through the displayed orders.
\end{exercise}

\section{Composition, inverse, and recursive equations}

\begin{exercise}[A tiny inner perturbation]
\label{ex:tiny-composition}
Expand $\e^{x+\e^{-x}}$ through the term proportional to $\e^{-2x}$ after
collecting absolute transmonomials. Check the same result using Taylor's formula
for $T(x+U)$.
\end{exercise}

\begin{exercise}[Inverse of $t+\log t$]
\label{ex:inverse-t-log-t}
The inverse of $T(t)=t+\log t$ is $S(x)=W(\e^x)$. Put $L=\log x$ and verify by
composition that
\[
  S(x)=x-L+\frac{L}{x}
   +\frac{L^2/2-L}{x^2}
   +\frac{L^3/3-3L^2/2+L}{x^3}
   +\bigO\!\left(\frac{L^4}{x^4}\right).
\]
\end{exercise}

\begin{exercise}[An exponential fixed point]
\label{ex:exp-fixed-point}
Solve
\[
  Y=x+\e^{-Y}
\]
through the fifth power of $q=\e^{-x}$. First derive the coefficient recursion by
substitution. Then recognize the closed form involving Lambert $W$.
\end{exercise}

\begin{exercise}[A logarithmic fixed point]
\label{ex:log-fixed-point}
Find a transseries solution through $x^{-3}$ of
\[
  U=\log(x+U).
\]
Write the coefficients as polynomials in $L=\log x$ and verify the residual.
\end{exercise}

\begin{exercise}[Riccati coefficient recursion]
\label{ex:riccati-recursion}
For
\[
  y'+y=\frac1x+y^2,
\]
insert $y=\sum_{n\ge1}a_nx^{-n}$. Derive the recurrence for $a_n$ and compute
$a_1,\ldots,a_6$.
\end{exercise}

\begin{exercise}[A boundary-layer parameter]
\label{ex:boundary-layer-parameter}
For
\[
  \eps y'+y=\frac1{1+x},\qquad y(0)=y_0,
\]
write the exact solution with the integral beginning at $0$. Show formally that,
near the outer region, the coefficient of $\e^{-x/\eps}$ begins
\[
  C(\eps)\asympseries
  y_0-1-\eps-2!\eps^2-3!\eps^3-\cdots.
\]
Interpret this result without pretending that the factorial series converges.
\end{exercise}

\section{Borel summation and Stokes jumps}

\begin{exercise}[An alternating Euler series]
\label{ex:alternating-euler}
Let
\[
  \widehat F(z)=\sum_{n=0}^{\infty}\frac{(-1)^n n!}{z^{n+1}}.
\]
Compute its Borel transform. Show that its positive-direction Borel sum is
unambiguous for $z>0$, and verify directly that the summed function satisfies
\[
  F'-F=-\frac1z.
\]
\end{exercise}

\begin{exercise}[Convolution in action]
\label{ex:borel-convolution}
Using Exercise~\ref{ex:alternating-euler}, compute
\[
  \Borel(\widehat F^2)(\xi)
\]
by convolution and simplify the integral to an elementary function.
\end{exercise}

\begin{exercise}[The sign of the Euler jump]
\label{ex:euler-jump}
Starting from
\[
  \Borel\widehat E(\xi)=\frac1{1-\xi},
\]
use a small contour around $\xi=1$ to derive
\[
  \mathcal S_{0^+}\widehat E-
  \mathcal S_{0^-}\widehat E=2\pi i\e^{-z}
\]
with the lateral convention of Chapter~\ref{ch:stokes}. Draw the orientation and
account for both minus signs.
\end{exercise}

\begin{exercise}[A pole predicts late coefficients]
\label{ex:pole-late-coefficients}
Suppose
\[
  \Borel\phihat(\xi)=h(\xi)+\frac{r}{\xi-A},
\]
where $h$ is analytic in a disk larger than $\abs A$. Derive the leading
large-$n$ behavior of the coefficients in
$\phihat(z)=\sum a_nz^{-n-1}$. What changes if two poles have the same modulus?
\end{exercise}

\begin{exercise}[Least-term estimate]
\label{ex:least-term}
For $z=100$, estimate the best truncation index and least-term magnitude of
\[
  \sum_{n\ge0}\frac{n!}{z^{n+1}}.
\]
Use Stirling's formula. Compare the answer with the exponential scale $\e^{-100}$.
\end{exercise}

\begin{exercise}[Higher Gevrey growth]
\label{ex:higher-gevrey}
Consider $a_n=(2n)!$. Determine its Gevrey order under
the bound \eqref{eq:gevrey-s-bound}. Show that the ordinary Borel series
$\sum a_n\xi^n/n!$ still has radius zero. What kind of generalized Borel
transform is suggested?
\end{exercise}

\begin{exercise}[Oscillation from conjugate singularities]
\label{ex:conjugate-singularities}
Let
\[
  \widehat\varphi(\xi)
  =\frac1{1-2(\cos\theta)\xi+\xi^2},
  \qquad0<\theta<\pi.
\]
Find its Taylor coefficients and hence the coefficients $a_n$ of the formal
inverse-power series whose Borel transform is $\widehat\varphi$. Explain the
oscillation in terms of the two singularities $\e^{\pm i\theta}$.
\end{exercise}

\begin{exercise}[Median value of the Euler integral]
\label{ex:euler-median}
For real $z>0$, show that
\[
  \frac12\left(
    \mathcal S_{0^+}\widehat E+
    \mathcal S_{0^-}\widehat E
  \right)
\]
equals the Cauchy principal value of
\[
  \int_0^\infty\frac{\e^{-z\xi}}{1-\xi}\dd\xi.
\]
Why is this average real? How does it relate to median resummation in this linear
example?
\end{exercise}

\section{Resurgence, saddles, and Stokes data}

\begin{exercise}[Bridge equation gives a translation]
\label{ex:bridge-translation}
Assume
\[
  \dot\Delta_A\widehat\Phi=S_1\partial_\sigma\widehat\Phi
\]
and that no other singularity lies on the chosen Stokes ray. Derive the action
of the Stokes automorphism on $\sigma$. Then derive its half-power used in median
summation.
\end{exercise}

\begin{exercise}[Airy prefactor]
\label{ex:airy-prefactor}
Insert
\[
  y(x)=\e^{\lambda(2/3)x^{3/2}}x^\alpha,
  \qquad\lambda=\pm1,
\]
into $y''-xy=0$. After the leading $x$ terms cancel, determine $\alpha$ from the
next largest power. What is the size of the residual after this choice?
\end{exercise}

\begin{exercise}[Complex saddles of a quartic integral]
\label{ex:quartic-saddles}
Consider
\[
  I(g)=\int_{-\infty}^{\infty}
  \exp\!\left[-\frac1g
    \left(\frac{x^2}{2}+\frac{x^4}{4}\right)
  \right]\dd x,
  \qquad g>0.
\]
Find all critical points of the action and their action values. Rescale
$x=\sqrt g\,t$ and derive the formal perturbative coefficients about $x=0$.
Explain why they alternate and grow factorially, and relate the sign to the
nearest complex saddles.
\end{exercise}

\begin{exercise}[Adding actions]
\label{ex:adding-actions}
Suppose a nonlinear equation has sectors with actions $A$ and $B$. Show
algebraically and in the Borel plane why a sector with action $A+B$ should be
expected. What special issue occurs when $B=2A$ and the equation already has an
independent sector of action $2A$?
\end{exercise}

\begin{exercise}[A transasymptotic logistic solution]
\label{ex:logistic-transasymptotic}
For the exact logistic solution
\[
  y(x)=\frac1{1+A\e^{-x}},
\]
set $\tau=A\e^{-x}$. Show that the sector expansion is a geometric series in
$\tau$, while the transasymptotic rearrangement sums it to one coefficient
function $F_0(\tau)$. Explain which representation is preferable when
$\abs\tau\ll1$ and when $\tau=O(1)$.
\end{exercise}

\section{Structural and computational questions}

\begin{exercise}[$H$-field intuition]
\label{ex:h-field-intuition}
For each of the germs
\[
  x^3+\log x,\qquad 7+x^{-1},\qquad \e^{-x},
\]
identify whether it is larger than every constant, bounded with a constant part,
or infinitesimal. Check the two $H$-field axioms on these examples.
\end{exercise}

\begin{exercise}[Logarithmic derivatives and scale]
\label{ex:log-derivative-scale}
Compute the logarithmic derivatives of
\[
  x^a(\log x)^b,
  \qquad \e^{x^p},
  \qquad \e^{\e^x}.
\]
Order the resulting logarithmic derivatives for fixed nonzero $a,b$ and $p>1$.
What do they say about sensitivity of the original monomials to a small change in
$x$?
\end{exercise}

\begin{exercise}[A surreal evaluation]
\label{ex:surreal-evaluation}
Formally evaluate
\[
  T(x)=x^2+3x+\log x+5+\e^{-x}
\]
at $x=\omega$. Identify the leading Conway monomial and order every displayed
term. Which term is infinitesimal? Why does this evaluation not by itself define
function composition on all surreal numbers?
\end{exercise}

\begin{exercise}[A difference-equation check]
\label{ex:digamma-check}
Verify by expansion through $x^{-5}$ that
\[
  y(x)=\log x-\frac1{2x}-\frac1{12x^2}
       +\frac1{120x^4}+\bigO(x^{-6})
\]
satisfies
\[
  y(x+1)-y(x)=\frac1x+\bigO(x^{-6}).
\]
Why is an additive constant not determined by the difference equation?
\end{exercise}

\begin{exercise}[A finite ratio grid]
\label{ex:ratio-grid}
Let
\[
  \mu_1=x^{-1},\qquad\mu_2=\e^{-x},
\]
and consider
\[
  T=1+2x^{-1}+3x^{-2}+\e^{-x}(1-4x^{-1}+7x^{-3})
    +5\e^{-2x}x^{-2}.
\]
Represent every monomial by an exponent vector $(k_1,k_2)$ with
$\mu_1^{k_1}\mu_2^{k_2}$. Draw the occupied lattice points. Give a rectangular
cut that retains all displayed terms and state why it is not the same as a
single dominance cut.
\end{exercise}

\begin{exercise}[Exact equality or equality to a cut?]
\label{ex:equality-cut}
A symbolic program expands
\[
  \log(\e^x(1+\e^{-x}))-x-\log(1+\e^{-x})
\]
through $\e^{-20x}$ and obtains zero. What can it conclude? Now simplify the
expression using exact formal identities. Contrast the two arguments.
\end{exercise}

\begin{exercise}[Classify the claim]
\label{ex:classify-claim}
Classify each statement as formal, asymptotic, summability, or global analytic.
Explain why no statement automatically implies all the others.
\begin{enumerate}[label=(\alph*)]
\item A recursive transseries solves an ODE coefficient by coefficient.
\item Its Borel transform continues along a ray with exponential growth bounds.
\item The Borel sum solves the ODE in a sector.
\item That solution extends meromorphically to the whole plane.
\item The residual after $N$ terms is $\bigO(x^{-N-1})$.
\end{enumerate}
\end{exercise}

\begin{exercise}[Capstone: the Euler equation]
\label{ex:euler-capstone}
For $y'+y=1/z$, carry out the complete workflow:
\begin{enumerate}[label=(\alph*)]
\item derive the perturbative recurrence;
\item identify the homogeneous sector;
\item compute the Borel transform;
\item calculate the lateral jump;
\item derive the parameter transformation;
\item estimate optimal truncation;
\item verify that the Borel sums satisfy the ODE.
\end{enumerate}
State clearly which steps are formal and which are analytic.
\end{exercise}

\chapter{Complete solutions}
\label{ch:solutions}

\section{Solutions to asymptotic-scale exercises}

\subsection*{Solution to Exercise~\ref{ex:order-mixed}}

The order is
\begin{align*}
  \e^{\e^x}
  &\succdom \e^{x^2}
  \succdom x^x
  \succdom x^{-100}\e^{3x}
  \succdom x^{100}
  \succdom(\log x)^7\\
  &\succdom1
  \succdom x^{-2}
  \succdom\e^{-\sqrt x}
  \succdom\e^{-x}.
\end{align*}
For the first comparison, take logarithms of the quotient:
\[
  \log\frac{\e^{\e^x}}{\e^{x^2}}=\e^x-x^2\to+\infty.
\]
Next, $x^x=\e^{x\log x}$ and
$x^2-x\log x=x(x-\log x)\to+\infty$. Also,
\[
  \log\frac{x^x}{x^{-100}\e^{3x}}
  =x\log x-3x+100\log x\to+\infty,
\]
and
\[
  \log\frac{x^{-100}\e^{3x}}{x^{100}}
  =3x-200\log x\to+\infty.
\]
Every positive power of $x$ dominates every power of $\log x$, while
$(\log x)^7\to\infty$. The comparisons around $1$ are immediate. Finally,
$x^{-2}/\e^{-\sqrt x}=x^{-2}\e^{\sqrt x}\to\infty$, and
\[
  \frac{\e^{-\sqrt x}}{\e^{-x}}=\e^{x-\sqrt x}\to\infty.
\]

\subsection*{Solution to Exercise~\ref{ex:rational-expansion}}

Divide numerator and denominator by $x^2$:
\[
  \frac{x^2+1}{x^2-x}
  =\frac{1+x^{-2}}{1-x^{-1}}.
\]
Using $(1-u)^{-1}=1+u+u^2+\cdots$ with $u=x^{-1}$,
\begin{align*}
  (1+x^{-2})(1+x^{-1}+x^{-2}+x^{-3}+x^{-4}+\cdots)
  &=1+x^{-1}+2x^{-2}+2x^{-3}+2x^{-4}
    +\bigO(x^{-5}).
\end{align*}
Multiplying by $1-x^{-1}$ gives
\[
  1+x^{-2}+\bigO(x^{-5}),
\]
which is the normalized numerator to the required order.

\subsection*{Solution to Exercise~\ref{ex:beyond-power}}

Set $t=\sqrt x$. Then
\[
  x^N\e^{-\sqrt x}=t^{2N}\e^{-t}.
\]
Repeated use of l'H\^opital's rule, or the fact that the exponential series
contains a term larger than any fixed power, gives
$t^{2N}\e^{-t}\to0$. Hence $\e^{-\sqrt x}\precdom x^{-N}$.
Moreover,
\[
  \frac{\e^{-x}}{\e^{-\sqrt x}}
  =\e^{-x+\sqrt x}\to0,
\]
so $\e^{-x}\precdom\e^{-\sqrt x}$.

\subsection*{Solution to Exercise~\ref{ex:flat-ambiguity}}

Fix $N$. Since $\e^{-x}=\smallo(x^{-N})$,
\[
  \left[f(x)+C\e^{-x}\right]
  -\sum_{n=0}^{N-1}a_nx^{-n}
  =\left[f(x)-\sum_{n=0}^{N-1}a_nx^{-n}\right]
   +C\e^{-x}
  =\bigO(x^{-N}).
\]
This holds for every $N$, so all the functions have the same Poincar\'e
expansion. They are not equal because their pairwise differences are
$(C_1-C_2)\e^{-x}$, nonzero when $C_1\ne C_2$. The power scale is simply too
coarse to see that difference.

\subsection*{Solution to Exercise~\ref{ex:normal-form}}

The first two terms share the large factor $\e^{x^2}$, but their quotient is
\[
  \frac{x^{100}\e^{x^2}}{\e^{x^2+x}}
  =x^{100}\e^{-x}\to0.
\]
Thus
\[
  \magop T=\e^{x^2+x},
  \qquad\lc(T)=3.
\]
Both exponential terms are giants and $7x^{-2}$ is small, so
\[
  T_{\mathrm{giant}}
   =3\e^{x^2+x}-5x^{100}\e^{x^2},
  \qquad T_{\mathrm{constant}}=0,
  \qquad T_{\mathrm{small}}=7x^{-2}.
\]
Factoring the leading term gives
\[
  T=3\e^{x^2+x}
  \left(
    1-\frac53x^{100}\e^{-x}
      +\frac73x^{-2}\e^{-x^2-x}
  \right).
\]
Both correction monomials tend to zero, so the parenthesized correction is
$1+U$ with $U\precdom1$.

\section{Solutions to support and algebra exercises}

\subsection*{Solution to Exercise~\ref{ex:legal-supports}}

The supports of
\[
  \sum_{n\ge0}x^{-n}
  \quad\text{and}\quad
  \sum_{n\ge0}\e^{-nx}
\]
are well based: each has a largest monomial $1$, and moving through the index
$n$ gives a strictly decreasing chain.

The support $\{x^n:n\ge0\}$ is illegal for a large-$x$ transseries because it has
no largest monomial and contains the infinite increasing chain
$1\precdom x\precdom x^2\precdom\cdots$.

The support $\{x^{-1/n}:n\ge1\}$ is also illegal. Its monomials increase toward
$1$:
\[
  x^{-1}\precdom x^{-1/2}\precdom x^{-1/3}\precdom\cdots,
\]
and no member is largest. The fact that the exponents converge does not rescue
the support; the order condition is combinatorial, not topological.

\subsection*{Solution to Exercise~\ref{ex:finite-product}}

To obtain $x^{-k}$, choose $x^{-j}$ from $A$ and $x^{-(k-j)}$ from $B$. The
allowed values are $j=0,\ldots,k$, so the coefficient is the finite sum
\[
  \sum_{j=0}^k(k-j+1)
  =1+2+\cdots+(k+1)
  =\frac{(k+1)(k+2)}2.
\]
Writing $r=x^{-1}$,
\[
  A=\frac1{1-r},
  \qquad B=\frac1{(1-r)^2},
  \qquad AB=\frac1{(1-r)^3}.
\]
The binomial expansion
\[
  (1-r)^{-3}=\sum_{k\ge0}\binom{k+2}{2}r^k
\]
confirms the coefficient.

\subsection*{Solution to Exercise~\ref{ex:two-scale-inverse}}

Put $a=x^{-1}$ and $q=\e^{-x}$. Expand first in $q$:
\[
  \frac1{1+a+q}
  =\frac1{1+a}-\frac{q}{(1+a)^2}+\bigO(q^2).
\]
Now
\begin{align*}
  (1+a)^{-1}&=1-a+a^2-a^3+\bigO(a^4),\\
  (1+a)^{-2}&=1-2a+3a^2-4a^3+\cdots.
\end{align*}
Therefore the requested two-scale truncation is
\begin{equation}
  1-x^{-1}+x^{-2}-x^{-3}
  -\e^{-x}+2x^{-1}\e^{-x}-3x^{-2}\e^{-x}
  +\text{omitted terms}.
  \label{eq:solution-two-scale-inverse}
\end{equation}
Every algebraic monomial $x^{-N}$ dominates $\e^{-x}$. Hence a genuine initial
segment reaching the term $\e^{-x}$ would have to retain infinitely many
algebraic powers. The finite selection in
\eqref{eq:solution-two-scale-inverse} is a useful rectangular cut in the
$(a,q)$ exponent lattice, not a single cut in total dominance order.

\subsection*{Solution to Exercise~\ref{ex:log-perturbed}}

Let $a=x^{-1}$ and $u=2a-a^2$. Then
\[
  \log\left[x^3(1+u)\right]=3\log x+\log(1+u).
\]
Through $a^3$,
\begin{align*}
  u&=2a-a^2,\\
  u^2&=4a^2-4a^3+\bigO(a^4),\\
  u^3&=8a^3+\bigO(a^4).
\end{align*}
Thus
\begin{align*}
  \log(1+u)
  &=u-\frac{u^2}{2}+\frac{u^3}{3}+\bigO(a^4)\\
  &=2a-3a^2+\frac{14}{3}a^3+\bigO(a^4).
\end{align*}
Therefore
\[
  \boxed{
  3\log x+\frac2x-\frac3{x^2}+\frac{14}{3x^3}
  +\bigO(x^{-4})}.
\]

\subsection*{Solution to Exercise~\ref{ex:integrate-exp-x2}}

Write
\[
  I=\e^{x^2}F,
  \qquad
  F=\frac1{2x}\sum_{n\ge0}c_nx^{-2n}.
\]
The equation $I'=\e^{x^2}$ becomes
\[
  F'+2xF=1.
\]
Equivalently, put
$F=\sum_{n\ge0}a_nx^{-2n-1}$, so $a_n=c_n/2$. Then
\[
  2a_0=1,
  \qquad 2a_n=(2n-1)a_{n-1}\quad(n\ge1).
\]
Hence
\[
  c_0=1,
  \qquad c_n=\frac{2n-1}{2}c_{n-1},
\]
and
\[
  c_0=1,
  \quad c_1=\frac12,
  \quad c_2=\frac34,
  \quad c_3=\frac{15}{8}.
\]
Thus
\[
  I(x)\asympseries
  \frac{\e^{x^2}}{2x}
  \left(1+\frac1{2x^2}+\frac3{4x^4}
    +\frac{15}{8x^6}+\cdots\right).
\]
Differentiating and using the recurrence cancels every negative even power in
$F'+2xF$ successively, leaving $1$ to the computed order.

\section{Solutions to composition and recursion exercises}

\subsection*{Solution to Exercise~\ref{ex:tiny-composition}}

With $q=\e^{-x}$,
\[
  \e^{x+q}=\e^x\e^q
  =\e^x\left(1+q+\frac{q^2}{2}+\frac{q^3}{6}+\cdots\right).
\]
After multiplying by $\e^x$,
\[
  \boxed{
  \e^{x+\e^{-x}}
  =\e^x+1+\frac12\e^{-x}+\frac16\e^{-2x}
   +\bigO(\e^{-3x})}.
\]
Taylor's formula gives the same result because every derivative of
$T(x)=\e^x$ equals $\e^x$:
\[
  T(x+q)=\sum_{n\ge0}\frac{T^{(n)}(x)}{n!}q^n
  =\e^x\sum_{n\ge0}\frac{q^n}{n!}.
\]

\subsection*{Solution to Exercise~\ref{ex:inverse-t-log-t}}

Write
\[
  S=x+s_0+\frac{s_1}{x}+\frac{s_2}{x^2}+\frac{s_3}{x^3},
\]
where
\[
  s_0=-L,
  \quad s_1=L,
  \quad s_2=\frac{L^2}{2}-L,
  \quad s_3=\frac{L^3}{3}-\frac32L^2+L.
\]
Since
\[
  \log S=L+
  \log\left(1+\frac{s_0}{x}+\frac{s_1}{x^2}
                    +\frac{s_2}{x^3}+\cdots\right),
\]
the logarithm through $x^{-3}$ is
\[
  \log S
  =L+\frac{s_0}{x}
   +\frac{s_1-s_0^2/2}{x^2}
   +\frac{s_2-s_0s_1+s_0^3/3}{x^3}
   +\bigO(L^4x^{-4}).
\]
Substitution gives
\[
  \log S
  =L-\frac{L}{x}
   +\frac{L-L^2/2}{x^2}
   +\frac{3L^2/2-L-L^3/3}{x^3}
   +\bigO(L^4x^{-4}).
\]
Adding $S$ cancels every displayed correction and yields
$S+\log S=x+\bigO(L^4x^{-4})$. Thus $T(S(x))=x$ to the stated order.

\subsection*{Solution to Exercise~\ref{ex:exp-fixed-point}}

Set $Y=x+\delta$ and $q=\e^{-x}$. The equation becomes
\begin{equation}
  \delta=q\e^{-\delta}.
  \label{eq:exercise-delta-equation}
\end{equation}
Let $\delta=c_1q+c_2q^2+\cdots$. Expanding
$\e^{-\delta}=1-\delta+\delta^2/2-\cdots$ and comparing powers gives
\[
  c_1=1,
  \quad c_2=-1,
  \quad c_3=\frac32,
  \quad c_4=-\frac83,
  \quad c_5=\frac{125}{24}.
\]
For the closed form, multiply \eqref{eq:exercise-delta-equation} by
$\e^\delta$:
\[
  \delta\e^\delta=q,
\]
so $\delta=W(q)$. Lagrange inversion gives
\[
  W(q)=\sum_{n\ge1}\frac{(-n)^{n-1}}{n!}q^n.
\]
Hence
\[
  \boxed{
  Y=x+q-q^2+\frac32q^3-\frac83q^4
      +\frac{125}{24}q^5+\bigO(q^6)}.
\]

\subsection*{Solution to Exercise~\ref{ex:log-fixed-point}}

Let $L=\log x$ and seek
\[
  U=L+\frac{a_1}{x}+\frac{a_2}{x^2}+\frac{a_3}{x^3}+\cdots,
\]
where the $a_j$ are polynomials in $L$. Since
\[
  \log(x+U)=L+\log\left(1+\frac{U}{x}\right)
  =L+\frac{U}{x}-\frac{U^2}{2x^2}
      +\frac{U^3}{3x^3}+\cdots,
\]
comparison gives
\begin{align*}
  a_1&=L,\\
  a_2&=a_1-\frac{L^2}{2}=L-\frac{L^2}{2},\\
  a_3&=a_2-La_1+\frac{L^3}{3}
      =L-\frac32L^2+\frac13L^3.
\end{align*}
Thus
\[
  \boxed{
  U=L+\frac{L}{x}
  +\frac{L-L^2/2}{x^2}
  +\frac{L-3L^2/2+L^3/3}{x^3}
  +\bigO(L^4x^{-4})}.
\]
Substituting this expression into the logarithm reproduces it through $x^{-3}$,
so the residual is $\bigO(L^4x^{-4})$.

\subsection*{Solution to Exercise~\ref{ex:riccati-recursion}}

The coefficient of $x^{-1}$ gives $a_1=1$. For $m\ge2$, the derivative
contributes $-(m-1)a_{m-1}$ to the coefficient of $x^{-m}$, while
$y^2$ contributes
$\sum_{j=1}^{m-1}a_ja_{m-j}$. Hence
\begin{equation}
  a_m=(m-1)a_{m-1}
      +\sum_{j=1}^{m-1}a_ja_{m-j}.
  \label{eq:solution-riccati-recurrence}
\end{equation}
Successive values are
\begin{align*}
  a_1&=1,\\
  a_2&=1+1=2,\\
  a_3&=2\cdot2+(1\cdot2+2\cdot1)=8,\\
  a_4&=3\cdot8+(1\cdot8+2\cdot2+8\cdot1)=44,\\
  a_5&=4\cdot44+(44+16+16+44)=296,\\
  a_6&=5\cdot296+(296+88+64+88+296)=2312.
\end{align*}
The rapidly growing numbers signal factorial divergence.

\subsection*{Solution to Exercise~\ref{ex:boundary-layer-parameter}}

Multiplying the ODE by $\e^{x/\eps}$ and integrating from $0$ to $x$ gives
\begin{equation}
  y(x)=\e^{-x/\eps}
  \left[y_0+\frac1\eps\int_0^x
    \frac{\e^{t/\eps}}{1+t}\dd t\right].
  \label{eq:boundary-exact-solution}
\end{equation}
A formal outer particular solution is
\[
  \widehat y_{\mathrm{out}}(x,\eps)
  =\sum_{n\ge0}\frac{n!\eps^n}{(1+x)^{n+1}}.
\]
The complete formal solution is
$\widehat y_{\mathrm{out}}+C(\eps)\e^{-x/\eps}$. Imposing the boundary condition
at $x=0$ therefore gives, formally,
\[
  C(\eps)=y_0-\widehat y_{\mathrm{out}}(0,\eps)
  \asympseries y_0-1-\eps-2!\eps^2-3!\eps^3-\cdots.
\]
This is an asymptotic description of the boundary-layer amplitude, not a
convergent formula. It may be optimally truncated or Borel summed in an
appropriate direction. Its role is to match the boundary data while the factor
$\e^{-x/\eps}$ makes the correction disappear in the outer region.

\section{Solutions to Borel and Stokes exercises}

\subsection*{Solution to Exercise~\ref{ex:alternating-euler}}

Dividing each coefficient by $n!$ gives
\[
  \Borel\widehat F(\xi)
  =\sum_{n\ge0}(-1)^n\xi^n
  =\frac1{1+\xi}.
\]
Its only finite singularity is at $\xi=-1$, so the positive Laplace ray is clear.
For $z>0$,
\begin{equation}
  F(z)=\int_0^\infty\frac{\e^{-z\xi}}{1+\xi}\dd\xi
  \label{eq:alternating-euler-sum}
\end{equation}
converges absolutely. Differentiate under the integral sign:
\begin{align*}
  F'(z)
  &=-\int_0^\infty\frac{\xi\e^{-z\xi}}{1+\xi}\dd\xi\\
  &=-\int_0^\infty\e^{-z\xi}\dd\xi
    +\int_0^\infty\frac{\e^{-z\xi}}{1+\xi}\dd\xi\\
  &=-\frac1z+F(z).
\end{align*}
Thus $F'-F=-1/z$. Watson's lemma recovers the alternating factorial expansion.

\subsection*{Solution to Exercise~\ref{ex:borel-convolution}}

By the convolution rule,
\[
  \Borel(\widehat F^2)(\xi)
  =\int_0^\xi
    \frac{\dd\eta}{(1+\eta)(1+\xi-\eta)}.
\]
Use
\[
  \frac1{(1+\eta)(1+\xi-\eta)}
  =\frac1{\xi+2}
   \left(\frac1{1+\eta}+\frac1{1+\xi-\eta}\right).
\]
Both terms integrate to $\log(1+\xi)$, so
\[
  \boxed{
  \Borel(\widehat F^2)(\xi)
  =\frac{2\log(1+\xi)}{\xi+2}}.
\]
The apparent singularity at $\xi=-2$ is entangled with the chosen branch of the
logarithm; convolution has propagated the original singularity at $-1$ and
created the summed location $-2$.

\subsection*{Solution to Exercise~\ref{ex:euler-jump}}

The upper contour runs from left to right above the pole; subtracting the lower
integral means traversing the lower contour from right to left. Together they
form a small clockwise loop around $\xi=1$. Therefore
\[
  \mathcal S_{0^+}\widehat E-
  \mathcal S_{0^-}\widehat E
  =-2\pi i\operatorname*{Res}_{\xi=1}
    \frac{\e^{-z\xi}}{1-\xi}.
\]
The residue is
\[
  \lim_{\xi\to1}
  (\xi-1)\frac{\e^{-z\xi}}{1-\xi}
  =-\e^{-z}.
\]
The clockwise orientation supplies one minus sign and the denominator
$1-\xi$ supplies the other:
\[
  -2\pi i(-\e^{-z})=2\pi i\e^{-z}.
\]

\subsection*{Solution to Exercise~\ref{ex:pole-late-coefficients}}

Expand the pole at the origin:
\[
  \frac{r}{\xi-A}
  =-r\sum_{n\ge0}\frac{\xi^n}{A^{n+1}}.
\]
Since
$\Borel\phihat=\sum a_n\xi^n/n!$, the pole contributes
\[
  a_n=-r\frac{n!}{A^{n+1}}.
\]
The coefficients of $h$ are exponentially smaller relative to this term because
$h$ is analytic in a larger disk. Thus
\[
  a_n\sim-r\frac{n!}{A^{n+1}}.
\]
If several singularities have the same modulus, their contributions must be
added. For a conjugate pair, this sum often looks like
$n!\abs A^{-n}\cos(n\theta+\phi)$ and oscillates. A simple ratio test then need
not converge.

\subsection*{Solution to Exercise~\ref{ex:least-term}}

For
$t_n=n!/100^{n+1}$,
\[
  \frac{t_{n+1}}{t_n}=\frac{n+1}{100}.
\]
The terms decrease until $n+1$ reaches $100$. The two terms $t_{99}$ and
$t_{100}$ are equal and minimal, so one truncates around index $99$ or $100$.
Stirling's formula gives
\[
  t_{100}
  =\frac{100!}{100^{101}}
  \sim\sqrt{\frac{2\pi}{100}}\e^{-100}
  \approx0.2507\e^{-100}
  \approx9.3\times10^{-45}.
\]
Thus the least term has exactly the predicted exponential scale $\e^{-100}$,
modulo an algebraic prefactor.

\subsection*{Solution to Exercise~\ref{ex:higher-gevrey}}

The coefficients satisfy the Gevrey-$2$ bound with $C=A=1$:
\[
  a_n=\Gamma(1+2n).
\]
The ordinary Borel coefficients are
$b_n=(2n)!/n!$. Their ratio is
\[
  \frac{b_{n+1}}{b_n}
  =\frac{(2n+2)(2n+1)}{n+1}
  =2(2n+1)\to\infty.
\]
Hence $\sum b_n\xi^n$ has radius zero. One should divide by
$\Gamma(1+2n)$, or use the equivalent Borel transform of the corresponding
higher order. Terminology alternates between ``Gevrey-$2$ summation'' and a
reciprocal $k$-summability convention, so the gamma factor is safer than the
bare name.

\subsection*{Solution to Exercise~\ref{ex:conjugate-singularities}}

The generating function for Chebyshev polynomials of the second kind is
\[
  \frac1{1-2t\xi+\xi^2}
  =\sum_{n\ge0}U_n(t)\xi^n.
\]
For $t=\cos\theta$,
\[
  U_n(\cos\theta)
  =\frac{\sin((n+1)\theta)}{\sin\theta}.
\]
Therefore
\[
  a_n=n!\frac{\sin((n+1)\theta)}{\sin\theta}.
\]
The denominator factors as
\[
  (1-\e^{i\theta}\xi)(1-\e^{-i\theta}\xi),
\]
so the Borel singularities are at $\xi=\e^{\pm i\theta}$. Their two
factorial-over-power contributions have conjugate phases and combine into the
sine. This is the simplest model of oscillatory large-order behavior.

\subsection*{Solution to Exercise~\ref{ex:euler-median}}

Remove a symmetric interval $(1-\delta,1+\delta)$ around the pole. The upper and
lower indentation integrals have opposite imaginary parts. Their average tends,
as $\delta\downarrow0$, to
\[
  \operatorname{PV}\int_0^\infty
  \frac{\e^{-z\xi}}{1-\xi}\dd\xi.
\]
For real $z>0$, complex conjugation maps the upper contour to the lower contour,
so the two lateral sums are conjugates and their average is real.

In this one-pole linear example, the symmetric average is the scalar shadow of
median resummation. In the full transseries, one more precisely applies half of
the Stokes automorphism to the parameter and then takes a lateral sum. The two
descriptions agree after the chosen homogeneous-sector normalization is taken
into account.

\section{Solutions to resurgence and saddle exercises}

\subsection*{Solution to Exercise~\ref{ex:bridge-translation}}

The Stokes automorphism is
\[
  \mathfrak S
  =\exp(\dot\Delta_A)
  =\exp(S_1\partial_\sigma).
\]
For any formal function of the parameter,
\[
  \exp(S_1\partial_\sigma)F(\sigma)
  =\sum_{n\ge0}\frac{S_1^n}{n!}F^{(n)}(\sigma)
  =F(\sigma+S_1).
\]
Hence
\[
  \mathfrak S\widehat\Phi(z,\sigma)
  =\widehat\Phi(z,\sigma+S_1).
\]
Its formal half-power is
\[
  \mathfrak S^{1/2}
  =\exp\left(\frac{S_1}{2}\partial_\sigma\right),
\]
which translates $\sigma$ by $S_1/2$. Signs reverse if the opposite convention
for relating upper and lower lateral sums is adopted.

\subsection*{Solution to Exercise~\ref{ex:airy-prefactor}}

Let
\[
  \phi(x)=\lambda\frac23x^{3/2}.
\]
Then
\[
  \frac{y'}y=\phi'+\frac\alpha x
  =\lambda\sqrt x+\frac\alpha x,
\]
and
\begin{align*}
  \frac{y''}{y}
  &=\left(\lambda\sqrt x+\frac\alpha x\right)^2
    +\frac{\lambda}{2\sqrt x}-\frac\alpha{x^2}\\
  &=x+\lambda\left(2\alpha+\frac12\right)x^{-1/2}
    +(\alpha^2-\alpha)x^{-2}.
\end{align*}
The leading $x$ cancels against the $xy$ term in Airy's equation. The next term
vanishes when
\[
  2\alpha+\frac12=0,
  \qquad\alpha=-\frac14.
\]
The remaining relative residual is
\[
  (\alpha^2-\alpha)x^{-2}
  =\frac5{16}x^{-2}.
\]
The first fluctuation correction in the WKB series cancels this residual.

\subsection*{Solution to Exercise~\ref{ex:quartic-saddles}}

The action is
\[
  S(x)=\frac{x^2}{2}+\frac{x^4}{4},
  \qquad S'(x)=x(1+x^2).
\]
Thus the saddles are
\[
  x_0=0,
  \qquad x_\pm=\pm i.
\]
Their actions are
\[
  S(0)=0,
  \qquad S(\pm i)=-\frac12+\frac14=-\frac14.
\]
Now put $x=\sqrt g\,t$:
\begin{align*}
  I(g)
  &=\sqrt g\int_{-\infty}^{\infty}
    \e^{-t^2/2}\e^{-gt^4/4}\dd t\\
  &\asympseries
    \sqrt g\sum_{n\ge0}\frac{(-g/4)^n}{n!}
    \int_{-\infty}^{\infty}t^{4n}\e^{-t^2/2}\dd t.
\end{align*}
The Gaussian moment is
\[
  \sqrt{2\pi}(4n-1)!!,
\]
so
\begin{equation}
  I(g)\asympseries\sqrt{2\pi g}
  \sum_{n\ge0}c_ng^n,
  \qquad
  c_n=(-1)^n\frac{(4n-1)!!}{4^n n!}
      =(-1)^n\frac{(4n)!}{2^{4n}(2n)!n!}.
  \label{eq:quartic-coefficients}
\end{equation}
Their ratio is
\[
  \frac{c_{n+1}}{c_n}
  =-\frac{(4n+1)(4n+3)}{4(n+1)}
  \sim-4n.
\]
Thus the coefficients grow factorially and alternate. The law
$a_{n+1}/a_n\sim n/A$ points to $A=-1/4$, exactly the action difference from the
origin to either complex saddle. A negative Borel action produces alternating
signs and a singularity away from the positive summation ray.

\subsection*{Solution to Exercise~\ref{ex:adding-actions}}

In the original variable,
\[
  \left(\e^{-Az}\phihat_A\right)
  \left(\e^{-Bz}\phihat_B\right)
  =\e^{-(A+B)z}(\phihat_A\phihat_B).
\]
Therefore a nonlinear product forces a sector of action $A+B$. In the Borel
plane, multiplication becomes convolution. Singular behavior near $A$ convolved
with singular behavior near $B$ can pinch the convolution integral at
$A+B$, creating a singularity there.

If $B=2A$, the action $2A$ may arise both as an independent elementary sector and
as the product of two $A$ sectors. Their recursions mix. This is resonance: the
same transmonomial receives contributions from different mechanisms. Additional
powers of $z$, logarithms, or a larger parameter space may be needed to solve the
resulting compatibility equations.

\subsection*{Solution to Exercise~\ref{ex:logistic-transasymptotic}}

With $\tau=A\e^{-x}$,
\[
  y=\frac1{1+\tau}
  =1-\tau+\tau^2-\tau^3+\cdots
  \qquad(\abs\tau<1).
\]
The sector expansion is therefore the geometric transseries
$\sum_{k\ge0}(-1)^kA^k\e^{-kx}$. The transasymptotic rearrangement sums all those
sectors into
\[
  F_0(\tau)=\frac1{1+\tau},
\]
with no inverse-$x$ corrections in this exact example. When $\abs\tau\ll1$, the
sector expansion displays the hierarchy transparently. When $\tau=O(1)$, no
finite number of sectors is uniformly accurate, while $F_0(\tau)$ remains the
right compact description except near its pole $\tau=-1$.

\section{Solutions to structural and computational exercises}

\subsection*{Solution to Exercise~\ref{ex:h-field-intuition}}

The germ $x^3+\log x$ exceeds every real constant for large $x$, and
\[
  (x^3+\log x)'=3x^2+\frac1x>0.
\]
This illustrates the first $H$-field axiom.

The germ $7+x^{-1}$ is bounded and decomposes as the constant $7$ plus the
infinitesimal $x^{-1}$. The germ $\e^{-x}$ is itself infinitesimal, with constant
part $0$. These illustrate
$\mathcal O=C+\mathfrak o$. The examples do not prove the axioms for an entire
field, but they show exactly what the axioms are abstracting.

\subsection*{Solution to Exercise~\ref{ex:log-derivative-scale}}

Direct calculation gives
\begin{align*}
  \left(x^a(\log x)^b\right)^\dagger
  &=\frac{a}{x}+\frac{b}{x\log x},\\
  \left(\e^{x^p}\right)^\dagger
  &=p x^{p-1},\\
  \left(\e^{\e^x}\right)^\dagger
  &=\e^x.
\end{align*}
For fixed nonzero $a,b$ and $p>1$,
\[
  \e^x\succdom p x^{p-1}
  \succdom \frac{a}{x}+\frac{b}{x\log x}.
\]
The logarithmic derivative is the fractional rate of change: for small $h$,
\[
  \frac{f(x+h)}{f(x)}
  \approx\exp\!\left(h\frac{f'}f\right).
\]
Thus $\e^{\e^x}$ is extraordinarily sensitive to a small shift in $x$, while a
power-log monomial changes only by a relative amount of order $h/x$.

\subsection*{Solution to Exercise~\ref{ex:surreal-evaluation}}

The formal evaluation is
\[
  T(\omega)=\omega^2+3\omega+\log\omega+5+\e^{-\omega}.
\]
Its displayed terms have order
\[
  \omega^2\succdom\omega\succdom\log\omega
  \succdom1\succdom\e^{-\omega}.
\]
The leading Conway monomial is $\omega^2$, and $\e^{-\omega}$ is a positive
infinitesimal smaller than every power $\omega^{-n}$.

Evaluation at one infinite argument treats the transseries as a value in an
ordered exponential field. Composition requires more: it must assign values
coherently as the argument varies through positive infinite surreals and obey the
chain rule. Arbitrary surreal numbers do not automatically carry such a
compatible function-composition structure.

\subsection*{Solution to Exercise~\ref{ex:digamma-check}}

Use
\begin{align*}
  \log(x+1)-\log x
  &=x^{-1}-\frac12x^{-2}+\frac13x^{-3}
    -\frac14x^{-4}+\frac15x^{-5}-\frac16x^{-6}+\cdots,\\
  (x+1)^{-1}-x^{-1}
  &=-x^{-2}+x^{-3}-x^{-4}+x^{-5}-x^{-6}+\cdots,\\
  (x+1)^{-2}-x^{-2}
  &=-2x^{-3}+3x^{-4}-4x^{-5}+5x^{-6}+\cdots,\\
  (x+1)^{-4}-x^{-4}
  &=-4x^{-5}+10x^{-6}+\cdots.
\end{align*}
After multiplying by $-1/2$, $-1/12$, and $1/120$, respectively, the
coefficients of $x^{-2},\ldots,x^{-6}$ all cancel. Therefore
\[
  y(x+1)-y(x)=\frac1x+\bigO(x^{-7}),
\]
which is stronger than requested. If $y$ is a solution, so is $y+C$, because the
forward difference annihilates constants. A boundary condition or normalization
is needed to choose $C$.

\subsection*{Solution to Exercise~\ref{ex:ratio-grid}}

The occupied exponent vectors are
\[
  (0,0),(1,0),(2,0),
  (0,1),(1,1),(3,1),(2,2).
\]
Here $(k_1,k_2)$ represents
$x^{-k_1}\e^{-k_2x}$. A rectangle such as
\[
  0\le k_1\le3,
  \qquad0\le k_2\le2
\]
contains every displayed point, though it also contains unoccupied points. It is
not a dominance initial segment. For every $N$, the omitted algebraic monomial
$x^{-N}$ with vector $(N,0)$ dominates every monomial with $k_2\ge1$, including
terms retained in the rectangle.

\subsection*{Solution to Exercise~\ref{ex:equality-cut}}

The computation through $\e^{-20x}$ proves only that the expression is
$\bigO(\e^{-21x})$ relative to that chosen grid, or more carefully that all
coefficients above the cut vanish. It does not by itself prove exact zero.

Exact formal identities do prove zero:
\begin{align*}
  \log\left(\e^x(1+\e^{-x})\right)
  &=\log(\e^x)+\log(1+\e^{-x})\\
  &=x+\log(1+\e^{-x}),
\end{align*}
where positivity and the formal logarithm rules are valid. Subtracting the last
two terms gives exactly zero. The first argument is finite evidence; the second
is an identity in the transseries field.

\subsection*{Solution to Exercise~\ref{ex:classify-claim}}

\begin{enumerate}[label=(\alph*)]
\item This is a \emph{formal} statement: coefficients cancel in the formal
algebra.
\item This is a \emph{summability hypothesis}: it supplies the continuation and
growth needed for a Laplace integral.
\item This is a \emph{sectorial analytic} conclusion obtained from summation.
\item This is a \emph{global analytic} continuation statement, much stronger
than sectorial existence.
\item This is an \emph{asymptotic} estimate for finite truncations.
\end{enumerate}
A formal solution can fail to be summable. A summable sector can have singularities
preventing global continuation. A function can have an asymptotic series without
that series being Borel summable. Each arrow requires its own hypotheses.

\subsection*{Solution to Exercise~\ref{ex:euler-capstone}}

\paragraph{(a) Formal recurrence.}
Put $\widehat y=\sum_{n\ge0}a_nz^{-n-1}$. Then
\[
  \widehat y'+\widehat y=\frac1z
\]
gives $a_0=1$ and $a_n=n a_{n-1}$, hence $a_n=n!$.

\paragraph{(b) Homogeneous sector.}
The homogeneous equation $y'+y=0$ gives $C\e^{-z}$. Thus
\[
  \widehat y(z;C)=\sum_{n\ge0}\frac{n!}{z^{n+1}}+C\e^{-z}.
\]
This step is formal algebra plus an exact elementary ODE calculation.

\paragraph{(c) Borel transform.}
Dividing by $n!$ gives
\[
  \Borel\widehat E(\xi)=\sum_{n\ge0}\xi^n=\frac1{1-\xi}.
\]
The equality begins near $\xi=0$ and its rational continuation is analytic away
from $1$.

\paragraph{(d) Lateral jump.}
The residue calculation of Exercise~\ref{ex:euler-jump} yields
\[
  \mathcal S_{0^+}\widehat E-
  \mathcal S_{0^-}\widehat E=2\pi i\e^{-z}.
\]
This is analytic: it uses contour integrals and continuation.

\paragraph{(e) Parameter transformation.}
Representing one analytic solution on both sides requires
\[
  C_+=C_--2\pi i.
\]
This is the Stokes change of asymptotic coordinate.

\paragraph{(f) Optimal truncation.}
The ratio of consecutive terms is $(n+1)/\abs z$, so truncate near
$n\approx\abs z$. The least term is of order
$\abs z^{-1/2}\e^{-\abs z}$ up to a constant, matching the homogeneous
exponential scale.

\paragraph{(g) ODE verification.}
For a nonsingular ray,
\[
  E_\theta(z)=\int_0^{\e^{i\theta}\infty}
  \frac{\e^{-z\xi}}{1-\xi}\dd\xi.
\]
Then
\begin{align*}
  E_\theta'(z)+E_\theta(z)
  &=\int
    \frac{(1-\xi)\e^{-z\xi}}{1-\xi}\dd\xi\\
  &=\int_0^{\e^{i\theta}\infty}\e^{-z\xi}\dd\xi
   =\frac1z.
\end{align*}
Differentiation under the integral and convergence are analytic inputs. The full
solution adds $C\e^{-z}$, which solves the homogeneous equation exactly.

This one example contains the formal sector construction, factorial divergence,
Borel continuation, Stokes jump, parameter change, and numerical least-term
estimate in their simplest possible forms.

\chapter{Formula sheet and problem-solving checklist}
\label{ch:formula-sheet}

\section{Dominance rules worth memorizing}

For every fixed positive $a,b,c$,
\begin{align*}
  \e^{x^a}&\succdom x^b\succdom(\log x)^c\succdom1,\\
  x^{-b}&\succdom\e^{-x^a},\\
  \e^{x^a}&\succdom\e^{x^b}\quad(a>b),\\
  \e^{-x^a}&\precdom\e^{-x^b}\quad(a>b).
\end{align*}
More generally,
\[
  \e^A\succdom\e^B
  \quad\Longleftrightarrow\quad A-B\to+\infty
\]
for real asymptotic exponents. Take logarithms before trying to compare
complicated positive monomials.

\section{Normal forms}

For nonzero $T$, the additive leading form is
\[
  T=a\m+\text{smaller terms},
  \qquad a=\lc(T),\quad\m=\magop(T).
\]
If $T>0$, the multiplicative form is
\[
  T=a\m(1+U),
  \qquad a>0,
  \qquad U\precdom1.
\]
This one factorization unlocks all elementary functions:
\begin{align*}
  T^{-1}&=a^{-1}\m^{-1}\sum_{n\ge0}(-U)^n,\\
  T^r&=a^r\m^r\sum_{n\ge0}\binom rnU^n,\\
  \log T&=\log a+\log\m+
     \sum_{n\ge1}\frac{(-1)^{n+1}}nU^n.
\end{align*}
For $T=L+c+S$ with $L$ purely large, $c\in\R$, and $S\precdom1$,
\[
  \e^T=\e^L\e^c\sum_{n\ge0}\frac{S^n}{n!}.
\]

\section{Differential formulas}

\begin{align*}
  (TU)'&=T'U+TU',\\
  (T\compose U)'&=(T'\compose U)U',\\
  (\e^T)'&=T'\e^T,\\
  (\log T)'&=\frac{T'}T,\\
  (T^r)'&=rT^{r-1}T'.
\end{align*}
The logarithmic derivative
\[
  T^\dagger=\frac{T'}T
\]
is often easier to compare than $T'$ itself.

For asymptotic integration of $\e^{P(x)}$, set
\[
  I=\e^{P(x)}F(x),
  \qquad F'+P'F=Q
\]
when $I'=Q\e^P$. Begin with $F\sim Q/P'$ and iterate the derivative correction.

\section{Composition and inversion}

When $U\precdom x$ is small enough for Taylor expansion,
\[
  T(x+U)=\sum_{n\ge0}\frac{T^{(n)}(x)}{n!}U^n.
\]
To invert $T=x+R$ with $R\precdom x$, either solve
\[
  S=x-R\compose S
\]
by contraction or use Newton iteration
\[
  S_{n+1}=S_n-\frac{T\compose S_n-x}{T'\compose S_n}.
\]
Always verify by the residual $T\compose S-x$.

\section{Factorial divergence and optimal truncation}

If
\[
  a_n\sim K\frac{\Gamma(n+\beta)}{A^{n+\beta}},
\]
then
\[
  \frac{a_n}{a_{n-1}}\sim\frac{n+\beta-1}{A},
  \qquad N_*\approx\abs{Az},
\]
and the least term has an exponential factor $\e^{-\abs{Az}}$.

\section{Borel--Laplace formulas}

For
\[
  \phihat(z)=\sum_{n\ge0}a_nz^{-n-1},
\]
use
\begin{align*}
  \Borel\phihat(\xi)&=\sum_{n\ge0}\frac{a_n}{n!}\xi^n,\\
  \mathcal S_\theta\phihat(z)
  &=\int_0^{\e^{i\theta}\infty}
    \e^{-z\xi}\Borel\phihat(\xi)\dd\xi,\\
  \Borel(\phihat\widehat\psi)&=
  (\Borel\phihat)*(\Borel\widehat\psi).
\end{align*}
A singularity at $\xi=A$ predicts $\e^{-Az}$. A simple pole
$r/(\xi-A)$ predicts
\[
  a_n\sim-r\frac{n!}{A^{n+1}}.
\]

\section{A ten-question checklist}

When confronting a new problem, ask:
\begin{enumerate}
\item What is the limiting variable and in which complex sector does it move?
\item What are the candidate dominant balances?
\item Which monomials form the smallest closed asymptotic scale?
\item Is the proposed support well based?
\item Which coefficients or sector amplitudes are free?
\item Does substitution produce a recurrence with factorial late growth?
\item What actions are suggested by ratios or saddle differences?
\item Where are the Borel singularities and singular directions?
\item How must transseries parameters transform across those directions?
\item Have formal residuals, large-order predictions, and numerical values all
been checked independently?
\end{enumerate}


\chapter{Glossary}
\label{ch:glossary}

\begin{description}[style=nextline,leftmargin=2em]
\item[Action]
A quantity $A$ appearing in an exponential scale $\e^{-Az}$ or
$\e^{-A/\eps}$. In saddle-point problems it is often a difference of critical
values of an exponent.

\item[Alien derivative]
A derivation $\Delta_\omega$ that extracts resurgent singular information at the
Borel location $\omega$.

\item[Anti-Stokes line]
A convention-dependent term often used for a ray on which two exponentials have
equal magnitude. This book prefers the explicit phrase \emph{equal-magnitude
line}.

\item[Asymptotic couple]
An ordered value group together with a map recording the valuation of logarithmic
derivatives.

\item[Asymptotic expansion]
A formal series whose finite truncations approximate a function with successively
smaller remainders in a specified limit.

\item[Asymptotic scale]
An ordered sequence or family $\phi_0\succdom\phi_1\succdom\cdots$ used as the
coordinate axes of an expansion.

\item[Beyond all orders]
Smaller than every term in a specified algebraic scale, as $\e^{-x}$ is smaller
than every $x^{-N}$.

\item[Borel plane]
The $\xi$-plane in which a factorially divergent series becomes a locally
convergent Borel transform.

\item[Borel sum]
An analytic function obtained by Borel transformation, analytic continuation,
and a directional Laplace integral.

\item[Borel--Leroy transform]
A Borel transform dividing coefficients by $\Gamma(n+\beta)$ or a related gamma
factor adapted to an algebraic prefactor.

\item[Borel--Pad\'e method]
A numerical method that approximates a finite Borel series by a rational Pad\'e
approximant before Laplace integration.

\item[Bridge equation]
An identity equating alien differentiation of a transseries with a vector field
in its free parameters.

\item[Contractive map]
A formal map that makes differences asymptotically smaller. Iteration then fixes
successive transmonomials and produces a unique fixed point in a complete formal
setting.

\item[Dominance]
The comparison $f\precdom g$, meaning $f/g\to0$ analytically or that the leading
monomial of $f$ is smaller formally.

\item[Dulac series]
A generalized expansion involving powers and logarithms, common in local
analytic dynamics and return maps.

\item[Endless continuability]
The ability to analytically continue a Borel germ along arbitrary finite paths
after avoiding a locally finite singular set.

\item[Exponential height]
The number of recursive exponential levels needed to construct a transmonomial.

\item[Flat function]
A function smaller than every term of a chosen asymptotic scale. Flat functions
cause nonuniqueness of ordinary asymptotic expansions.

\item[Formal series]
A symbolic sum manipulated coefficientwise without an initial convergence claim.

\item[Froissart doublet]
A nearby pole-zero pair in a Pad\'e approximant, usually a numerical artifact.

\item[Gevrey order]
A measure of coefficient growth using bounds such as
$\abs{a_n}\le CA^n\Gamma(1+sn)$.

\item[Giant part]
The sum of terms in a transseries whose monomials dominate $1$.

\item[Grid]
A finitely generated set of monomials, typically
$\m_0\mu_1^{k_1}\cdots\mu_r^{k_r}$ with lower bounds on the integers $k_j$.

\item[Hahn field]
A field of generalized series with coefficients in a field and exponents or
monomials in an ordered abelian group, subject to a well-ordered support
condition.

\item[Hardy field]
A differential field of germs at infinity of real functions, ordered by eventual
sign.

\item[Height wins]
The rule that a genuinely higher exponential level dominates all lower-height
monomials, provided its top exponent is positive and purely large.

\item[Homogeneous sector]
A free solution of the linearized or homogeneous equation, often multiplied by
a transseries parameter.

\item[Hyperasymptotics]
Repeated optimal truncation and re-expansion of exponentially small remainders
using neighboring saddles or Borel singularities.

\item[Hyperseries]
Extensions of ordinary transseries allowing longer, often transfinite, logarithmic
or exponential hierarchies.

\item[$H$-field]
An ordered differential field axiomatizing basic eventual behavior of Hardy-field
germs and transseries.

\item[Instanton]
A nontrivial finite-action saddle in a semiclassical problem, contributing a
sector such as $\e^{-A/g}$.

\item[Laplace direction]
The ray along which the Borel transform is integrated against $\e^{-z\xi}$.

\item[Lateral sum]
A Borel--Laplace sum taken just above or below a singular direction.

\item[Leading coefficient]
The coefficient multiplying the largest monomial of a nonzero transseries.

\item[Leading monomial]
The largest monomial in the support of a nonzero transseries.

\item[Least term]
The smallest term in magnitude of a divergent asymptotic series at a fixed large
argument.

\item[Logarithmic depth]
The number of iterated logarithmic substitutions needed to express an element in
a chosen transseries hierarchy.

\item[Logarithmic derivative]
$T^\dagger=T'/T$, the relative rate of change of $T$.

\item[Median summation]
A symmetric resummation applying half a Stokes transformation so that conjugate
lateral ambiguities cancel under suitable reality conditions.

\item[Monomial group]
An ordered multiplicative abelian group whose elements label asymptotic sizes.

\item[Multisummability]
Compatible summation of several components of a formal solution having distinct
Gevrey levels.

\item[Non-Archimedean]
A setting with infinitesimals and infinitely large elements, where the largest
term usually controls a sum.

\item[Optimal truncation]
Stopping a divergent asymptotic series near its least term instead of summing
indefinitely.

\item[Purely large]
A transseries all of whose monomials dominate $1$ and with no constant or small
part.

\item[Ratio set]
A finite collection of small monomials used as generators for a grid.

\item[Resonance]
A coincidence among action combinations or characteristic exponents that forces
sector mixing and may require logarithmic terms.

\item[Resurgence]
The property that Borel singularities of one formal sector contain controlled
information about other sectors, yielding large-order relations and Stokes data.

\item[Sector]
A component of a transseries with a fixed exponential action and algebraic
prefactor, usually carrying its own fluctuation series.

\item[Singulant]
A possibly variable action function governing a late-term ansatz and an
exponential remainder, especially in exponential asymptotics.

\item[Small part]
The sum of terms whose monomials are dominated by $1$.

\item[Stokes automorphism]
The formal automorphism relating lower and upper lateral resummations across a
singular direction.

\item[Stokes constant]
A normalization-dependent coefficient measuring the coupling between resurgent
sectors across a Borel singularity.

\item[Stokes phenomenon]
The change of an asymptotic representation, often the coefficient of an
exponentially small sector, as a complex direction is crossed.

\item[Support]
The set of monomials having nonzero coefficients in a generalized series.

\item[Summable family]
A family of transseries for which the union of supports is well based and each
monomial receives contributions from only finitely many family members.

\item[Transasymptotic rearrangement]
A reorganization that treats a combined variable such as
$\tau=\sigma\e^{-Az}z^\beta$ as order one and sums infinitely many sectors into
coefficient functions of $\tau$.

\item[Transmonomial]
A formal symbol representing one asymptotic magnitude, built recursively from
powers, exponentials, and logarithms.

\item[Transseries]
A well-based generalized series of transmonomials; in resurgent analysis, also a
sectorial expansion combining divergent power series with exponential scales and
free parameters.

\item[Transexponential]
A scale beyond all finite iterates of the ordinary exponential, introduced to
solve certain iteration equations.

\item[Valuation]
An algebraic measure of smallness. In a transseries field it is determined by the
leading monomial, with larger valuation corresponding to smaller magnitude.

\item[Well-based support]
A support ordered so that every nonempty subset has a largest monomial, or
equivalently has no infinite chain increasing in dominance under the conventions
of this book.

\item[Witness]
In grid-based transseries calculations, finite ratio-set data certifying a
particular dominance or support relation.
\end{description}


\backmatter
\phantomsection
\addcontentsline{toc}{chapter}{Bibliography}

{\raggedright
\begin{thebibliography}{99}

\bibitem{ADHBook}
Matthias Aschenbrenner, Lou van den Dries, and Joris van der Hoeven,
\emph{Asymptotic Differential Algebra and Model Theory of Transseries},
Annals of Mathematics Studies, vol.~195, Princeton University Press, 2017.

\bibitem{AnicetoBasarSchiappa}
In\^es Aniceto, G\"ok\c{c}e Ba\c{s}ar, and Ricardo Schiappa,
``A Primer on Resurgent Transseries and Their Asymptotics,''
\emph{Physics Reports} \textbf{809} (2019), 1--135;
\url{https://arxiv.org/abs/1802.10441}.

\bibitem{BerarducciMantovaGerms}
Alessandro Berarducci and Vincenzo Mantova,
``Transseries as Germs of Surreal Functions,''
\emph{Transactions of the American Mathematical Society}
\textbf{371} (2019), 3549--3592;
\url{https://arxiv.org/abs/1703.01995}.

\bibitem{BerarducciMantovaSurreal}
Alessandro Berarducci and Vincenzo Mantova,
``Surreal Numbers, Derivations and Transseries,''
\emph{Journal of the European Mathematical Society}
\textbf{20} (2018), 339--390;
\url{https://arxiv.org/abs/1503.00315}.

\bibitem{ChenFangHoeven2026}
Shaoshi Chen, Hanqian Fang, and Joris van der Hoeven,
``A Zero-Test for $D$-Algebraic Transseries,'' preprint, 2026;
\url{https://arxiv.org/abs/2602.01188}.

\bibitem{CostinBook}
Ovidiu Costin,
\emph{Asymptotics and Borel Summability},
Chapman \& Hall/CRC Monographs and Surveys in Pure and Applied Mathematics,
vol.~141, CRC Press, 2009.

\bibitem{CostinRankOne}
Ovidiu Costin,
``On Borel Summation and Stokes Phenomena for Rank-One Nonlinear Systems of
Ordinary Differential Equations,''
\emph{International Mathematics Research Notices} 1995, no.~8, 377--417.

\bibitem{CostinTopological}
Ovidiu Costin,
``Topological Construction of Transseries and Introduction to Generalized
Borel Summability,'' in Ovidiu Costin, Martin D. Kruskal, and Angus Macintyre
(eds.), \emph{Analyzable Functions and Applications}, Contemporary
Mathematics, vol.~373, American Mathematical Society, 2005, pp.~137--175;
\url{https://arxiv.org/abs/math/0608309}.

\bibitem{DahnGoring}
Bernd I. Dahn and Peter G\"oring,
``Notes on Exponential-Logarithmic Terms,''
\emph{Fundamenta Mathematicae} \textbf{127} (1987), no.~1, 45--50.

\bibitem{DorigoniIntro}
Daniele Dorigoni,
``An Introduction to Resurgence, Trans-Series and Alien Calculus,''
\emph{Annals of Physics} \textbf{409} (2019), article 167914;
\url{https://arxiv.org/abs/1411.3585}.

\bibitem{EcalleAnalyzable}
Jean \'Ecalle,
\emph{Introduction aux fonctions analysables et preuve constructive de la
conjecture de Dulac}, Actualit\'es Math\'ematiques, Hermann, Paris, 1992.

\bibitem{EcalleResurgent}
Jean \'Ecalle,
\emph{Les Fonctions R\'esurgentes}, vols.~I--III,
Publications Math\'ematiques d'Orsay, 1981 and 1985.

\bibitem{EdgarBeginners}
Gerald A. Edgar,
``Transseries for Beginners,''
\emph{Real Analysis Exchange} \textbf{35} (2010), 253--310;
\url{https://arxiv.org/abs/0801.4877}.

\bibitem{EdgarComposition}
Gerald A. Edgar,
``Transseries: Composition, Recursion, and Convergence,'' preprint, 2009;
\url{https://arxiv.org/abs/0909.1259}.

\bibitem{EdgarFractional}
Gerald A. Edgar,
``Fractional Iteration of Series and Transseries,''
\emph{Transactions of the American Mathematical Society}
\textbf{365} (2013), no.~11, 5805--5832;
\url{https://arxiv.org/abs/1002.2378}.

\bibitem{Hahn1907}
Hans Hahn,
``\"Uber die nichtarchimedischen Gr\"o\ss ensysteme,''
\emph{Sitzungsberichte der Kaiserlichen Akademie der Wissenschaften in Wien,
Mathematisch-Naturwissenschaftliche Klasse} \textbf{116} (1907), 601--655.

\bibitem{HardyOrders}
G. H. Hardy,
\emph{Orders of Infinity: The ``Infinit\"arcalc\"ul'' of Paul du Bois-Reymond},
Cambridge Tracts in Mathematics and Mathematical Physics, no.~12,
Cambridge University Press, 1910.

\bibitem{KuhlmannTressl}
Salma Kuhlmann and Marcus Tressl,
``Comparison of Exponential-Logarithmic and Logarithmic-Exponential Series,''
\emph{Mathematical Logic Quarterly} \textbf{58} (2012), no.~6, 434--448.

\bibitem{Neumann1949}
B. H. Neumann,
``On Ordered Division Rings,''
\emph{Transactions of the American Mathematical Society}
\textbf{66} (1949), 202--252.

\bibitem{NumbersGermsTransseries}
Matthias Aschenbrenner, Lou van den Dries, and Joris van der Hoeven,
``On Numbers, Germs, and Transseries,'' expository preprint, 2017;
\url{https://arxiv.org/abs/1711.06936}.

\bibitem{PadgettTransexponential}
Adele Padgett,
``Sublogarithmic-Transexponential Series,'' preprint, 2022;
\url{https://arxiv.org/abs/2211.06736}.

\bibitem{SauzinIntro}
David Sauzin,
``Introduction to 1-Summability and Resurgence,'' preprint, 2014;
\url{https://arxiv.org/abs/1405.0356}.

\bibitem{SchmelingThesis}
Michael C. Schmeling,
\emph{Corps de transs\'eries}, Ph.D. thesis, Universit\'e Paris~VII, 2001.

\bibitem{vanDenDriesHoevenKaplan}
Lou van den Dries, Joris van der Hoeven, and Elliot Kaplan,
``Logarithmic Hyperseries,''
\emph{Transactions of the American Mathematical Society}
\textbf{372} (2019), 5199--5241;
\url{https://arxiv.org/abs/1810.01810}.

\bibitem{vanDenDriesMacintyreMarker1997}
Lou van den Dries, Angus Macintyre, and David Marker,
``Logarithmic-Exponential Power Series,''
\emph{Journal of the London Mathematical Society}
\textbf{56} (1997), 417--434.

\bibitem{vanDenDriesMacintyreMarker2001}
Lou van den Dries, Angus Macintyre, and David Marker,
``Logarithmic-Exponential Series,''
\emph{Annals of Pure and Applied Logic}
\textbf{111} (2001), 61--113.

\bibitem{vanDerHoevenBook}
Joris van der Hoeven,
\emph{Transseries and Real Differential Algebra},
Lecture Notes in Mathematics, vol.~1888, Springer, 2006.

\bibitem{vanDerHoevenOperators}
Joris van der Hoeven,
``Operators on Generalized Power Series,''
\emph{Illinois Journal of Mathematics} \textbf{45} (2001), no.~4,
1161--1190.

\end{thebibliography}
}

\end{document}
